\documentclass[11pt]{article}

\usepackage[margin=1.1in]{geometry}

\usepackage{amssymb,mathtools,natbib}

% Anchored formal environments for causalsmith papers.
% First mandatory argument = obj_id (machine anchor joining the paper to the
% Lean crosswalk; invisible in the rendered PDF). Optional argument = name.
% Each environment also defines \label{obj:<obj_id>} so prose can use
% \cref{obj:T-1}; cleveref derives the printed kind and number from the target.
\usepackage{amsmath}
\usepackage{mathrsfs}
\usepackage{amsthm}
\usepackage{booktabs}
% Long displays may break across pages; let TeX stretch a little harder before an
% overfull \hbox so wide formulas/tables are less likely to run off the page.
\allowdisplaybreaks
\setlength{\emergencystretch}{3em}
\newtheorem{theorem}{Theorem}
\newtheorem{lemma}{Lemma}
\newtheorem{assumption}{Assumption}
\newtheorem{definition}{Definition}
\newtheorem{proposition}{Proposition}
\newtheorem{remark}{Remark}
\newtheorem{citedresult}{Cited result}
% The amsthm counter is `algorithmbox` (not `algorithm`) so a later `\usepackage{algorithm}` cannot clash.
\newcounter{algorithmbox}
\renewcommand{\thealgorithmbox}{\arabic{algorithmbox}}
\NewDocumentEnvironment{theoremv}{m o s}{\begin{theorem}\label{obj:#1}{\normalfont[\texttt{\detokenize{#1}}]\IfValueT{#2}{\ (#2)}\IfBooleanT{#3}{\verificationfootnotemark}\ }}{\end{theorem}}
\NewDocumentEnvironment{lemmav}{m o s}{\begin{lemma}\label{obj:#1}{\normalfont[\texttt{\detokenize{#1}}]\IfValueT{#2}{\ (#2)}\IfBooleanT{#3}{\verificationfootnotemark}\ }}{\end{lemma}}
\NewDocumentEnvironment{assumptionv}{m o s}{\begin{assumption}\label{obj:#1}{\normalfont[\texttt{\detokenize{#1}}]\IfValueT{#2}{\ (#2)}\IfBooleanT{#3}{\verificationfootnotemark}\ }}{\end{assumption}}
\NewDocumentEnvironment{definitionv}{m o s}{\begin{definition}\label{obj:#1}{\normalfont[\texttt{\detokenize{#1}}]\IfValueT{#2}{\ (#2)}\IfBooleanT{#3}{\verificationfootnotemark}\ }}{\end{definition}}
% A `propositionv` is a proved result tiered as a Proposition (theorem-grade, machine-verified).
\NewDocumentEnvironment{propositionv}{m o s}{\begin{proposition}\label{obj:#1}{\normalfont[\texttt{\detokenize{#1}}]\IfValueT{#2}{\ (#2)}\IfBooleanT{#3}{\verificationfootnotemark}\ }}{\end{proposition}}
% A `remarkv` holds NON-load-bearing interpretive / open-question content (e.g. a causalsmith `oeq:`
% open question). It is neither proved nor assumed; no theorem may depend on it.
\NewDocumentEnvironment{remarkv}{m o s}{\begin{remark}\label{obj:#1}{\normalfont[\texttt{\detokenize{#1}}]\IfValueT{#2}{\ (#2)}\IfBooleanT{#3}{\verificationfootnotemark}\ }}{\end{remark}}
% An `algorithmv` is a DEFINITION-kind object presented as a boxed, numbered-step procedure (an
% estimator / selection rule built in stages). Same anchor contract as `definitionv`; its body is
% ordinary LaTeX whose steps are `\begin{enumerate}` items, so tex2html needs no extra machinery.
% Presented in the CONVENTIONAL algorithm style readers expect from `algorithm`+`algorithmic`:
% a rule above, an "Algorithm N (Title)" caption line, a rule under the caption, the numbered
% steps, and a closing rule. It is NOT a float — the anchor contract requires the object to stay
% where the outline places it, and floats drift. The obj-id tag is suppressed here (it is
% bookkeeping, and a caption line carrying it reads as debug output in a finished paper).
\NewDocumentEnvironment{algorithmv}{m o s}%
  {\par\addvspace{\medskipamount}\noindent\rule{\linewidth}{0.8pt}\par\nobreak\vspace{2pt}%
   \refstepcounter{algorithmbox}\label{obj:#1}%
   \noindent\textbf{Algorithm \thealgorithmbox}\IfValueT{#2}{ \textnormal{(#2)}}%
   \IfBooleanT{#3}{\verificationfootnotemark}\par\nobreak\vspace{2pt}%
   \noindent\rule{\linewidth}{0.4pt}\par\nobreak\vspace{2pt}\normalfont}%
  {\par\nobreak\vspace{2pt}\noindent\rule{\linewidth}{0.8pt}\par\addvspace{\medskipamount}}
% A `citedv` is an IMPORTED external result the paper relies on but does NOT prove
% (a causalsmith `gate_class:"cited"` node). It is machine-checked against the cited
% source at F2.5, never proved here; the fixed provenance note makes that explicit
% so the environment can never be read as a contribution of this paper. The \citep
% attribution is supplied by the surrounding prose (P2), keyed to references.bib.
\NewDocumentEnvironment{citedv}{m o s}{\begin{citedresult}\label{obj:#1}{\normalfont[\texttt{\detokenize{#1}}]\IfValueT{#2}{\ (#2)}\IfBooleanT{#3}{\verificationfootnotemark}\ \textit{(imported result; machine-checked against the cited source, not proved here.)}\ }}{\end{citedresult}}
% \leanref{obj_id}{display text}: an inline first-mention link to a from-note object's
% verified Lean code. In the PDF it renders the display text plainly (the Lean link is a
% web-only affordance); tex2html turns it into a drawer-opening span.
\newcommand{\leanref}[2]{#2}
% For a theorem/lemma whose Lean declaration accepts a source-matched published proposition as an
% input, the starred anchored environment puts the mark in its heading; the generated text remains
% outside the frozen mathematical environment. Keep the combined macro for older paper bundles.
\newcommand{\verificationfootnotemark}{\footnotemark}
\newcommand{\verificationfootnotetext}[1]{\footnotetext{#1}}
\newcommand{\verificationfootnote}[1]{\verificationfootnotemark\verificationfootnotetext{#1}}
% xcolor before hyperref (hyperref would otherwise pull in plain `color` for colorlinks, and a
% later xcolor load clashes). The page-1 identity stamp at the end of this file needs a grey.
\usepackage{xcolor}
\usepackage{graphicx}
\usepackage{eso-pic}
% hyperref follows natbib and the environment macros; cleveref must follow hyperref.
\usepackage[colorlinks=true,linkcolor=blue,citecolor=blue,urlcolor=blue]{hyperref}
\usepackage[capitalise,nameinlink,noabbrev]{cleveref}
% Pin the names explicitly: labels are placed inside the underlying amsthm environments, and these
% declarations make the PDF contract independent of cleveref's theorem-name inference. House style:
% reference names always print with a capital first letter ("Theorem 1", "Definition 2"), in any
% sentence position — hence `capitalise` above and capitalized \crefname forms below.
\crefname{theorem}{Theorem}{Theorems}
\Crefname{theorem}{Theorem}{Theorems}
\crefname{lemma}{Lemma}{Lemmas}
\Crefname{lemma}{Lemma}{Lemmas}
\crefname{assumption}{Assumption}{Assumptions}
\Crefname{assumption}{Assumption}{Assumptions}
\crefname{definition}{Definition}{Definitions}
\Crefname{definition}{Definition}{Definitions}
\crefname{proposition}{Proposition}{Propositions}
\Crefname{proposition}{Proposition}{Propositions}
\crefname{remark}{Remark}{Remarks}
\Crefname{remark}{Remark}{Remarks}
\crefname{citedresult}{Cited result}{Cited results}
\Crefname{citedresult}{Cited result}{Cited results}
\crefname{algorithmbox}{Algorithm}{Algorithms}
\Crefname{algorithmbox}{Algorithm}{Algorithms}
% ---------------------------------------------------------------------------
% Page-1 identity stamp (arXiv style) and the pinned title-block date.
%
% Split of responsibility: THIS file owns the FORMAT (placement, font, colour, punctuation),
% the generated `paper_stamp.tex` owns only the VALUES (number+version, area, revised date,
% DOI, link target). P4 refreshes this template on every emit, so a layout fix reaches every
% bundle from a recompile; and because `paper_stamp.tex` is not one of the P2 assembly
% digest's authored inputs, a DOI that arrives after publication needs only that one small
% file rewritten plus a recompile — never a redraft, never a reassembly.
%
% Defaults first, so a bundle WITHOUT a stamp file compiles exactly as it always did:
% `\cspaperdate` falls back to `\today` and no stamp is drawn.
\providecommand*{\cspaperdate}{\today}  % the \date{} target
\providecommand*{\csstampid}{}          % e.g. CSWP-2026-008v3 (empty ⇒ no stamp at all)
\providecommand*{\csstamparea}{}        % e.g. Stat
\providecommand*{\csstampdate}{}        % e.g. 27 Aug 2026
\providecommand*{\csstampdoi}{}         % e.g. 10.5281/zenodo.123 (empty ⇒ DOI part omitted)
\providecommand*{\csstamplink}{}        % href target (DOI url, else the paper's page; may be empty)
\IfFileExists{paper_stamp.tex}{% GENERATED by CausalSmith P4 — paper identity stamp values. Do not hand-edit.
%
% Field order on the stamp (owned by paper_macros.tex):
%   <number>v<version> · doi:<doi> · [<Area>] <date>   — the DOI is never the last token, so a
%   text extractor cannot run it into whatever follows. Without a DOI that part is omitted.
% Format lives in paper_macros.tex; only the values are here, and this file is NOT one of
% the P2 assembly digest's authored inputs. A DOI arriving after publication therefore needs
% this file rewritten plus a `latexmk` recompile — never a redraft or a reassembly.
\renewcommand*{\csstampid}{CSWP-2026-029v5}
\renewcommand*{\csstamparea}{Stat}
\renewcommand*{\csstampdate}{1 Oct 2026}
\renewcommand*{\csstampdoi}{}
\renewcommand*{\csstamplink}{https://causalsmith.org/papers/stat_rd_trueside_noise_frontier_v1}
\renewcommand*{\cspaperdate}{1 October 2026}
}{}
\makeatletter
% Field order is deliberate: the DOI is NEVER the last token on the line. A trailing identifier
% runs into whatever a text extractor emits next, which makes it unsafe to match mechanically
% (the Zenodo stamp matcher reads exactly this line); the date is a safe terminator.
%   <number>v<version> · doi:<doi> · [<Area>] <date>      — with a DOI
%   <number>v<version> · [<Area>] <date>                  — without
\newcommand{\cs@stampsep}{\,\textperiodcentered\,}
\newcommand{\cs@stamptext}{%
  \csstampid
  \ifx\csstampdoi\@empty\else\cs@stampsep doi:\csstampdoi\fi
  \ifx\csstamparea\@empty\else\cs@stampsep[\csstamparea]\fi
  \ifx\csstampdate\@empty\else\quad\csstampdate\fi
}
\ifx\csstampid\@empty\else
  \definecolor{csstampgrey}{gray}{0.45}
  % Background shipout material: it occupies no space, so the text block and the \thanks
  % footnote are byte-identical to an unstamped compile. The starred form applies to the
  % NEXT page shipped only — i.e. page 1. Placed in the left margin (the text block starts
  % at 1.1in), reading bottom-to-top, in a Times-like face at 8pt.
  \AddToShipoutPictureBG*{%
    \AtPageLowerLeft{%
      \put(\LenToUnit{0.42in},\LenToUnit{1.1in}){%
        \rotatebox{90}{%
          \fontfamily{ptm}\fontsize{8}{10}\selectfont
          \ifx\csstamplink\@empty
            \textcolor{csstampgrey}{\cs@stamptext}%
          \else
            \expandafter\href\expandafter{\csstamplink}{\textcolor{csstampgrey}{\cs@stamptext}}%
          \fi
        }%
      }%
    }%
  }
\fi
\makeatother


\title{Uniform estimation and honest inference in regression discontinuity with a mismeasured score}

\author{CausalSmith\thanks{Machine-generated by the CausalSmith research pipeline.}}

\date{\cspaperdate}

\begin{document}

\maketitle

\begin{abstract}
We characterize minimax estimation and honest inference for a latent-cutoff treatment effect when treatment records the true cutoff side and the running variable is measured with independent Gaussian error. The benchmark has known latent support \([-1,1]\), a continuous score density bounded between \(1/4\) and \(3/4\), and binary potential outcomes whose continuous conditional means lie between the same bounds and have full-interval Hölder seminorm at most two. For each fixed Hölder exponent \(\beta\in(0,1]\), minimax expected absolute error and minimax expected length of connected intervals with uniform \(90\%\) coverage have a common rate, uniformly over sample sizes \(n\ge2\) and known noise standard deviations \(\sigma\in[0,1]\). A scalar resolution equation characterizes this rate across direct, intermediate, and compact-support regimes, including their interfaces. The direct rate is \(n^{-\beta/(2\beta+1)}\). For each fixed positive noise level, the asymptotic rate as \(n\to\infty\) is \(\{\log\log(en)/\log(en)\}^{2\beta}\), with comparison constants permitted to depend on that fixed noise level. Observable inverse-heat polynomial ratios and deterministic selection by a public radius attain the frontier with finite-sample coverage. Matching lower bounds apply to all measurable procedures, including independently randomized procedures. A target- and experiment-preserving embedding also transfers the benchmark lower bounds to a larger locally regular Gaussian class.
\end{abstract}

\section{Introduction}

This paper establishes a common minimax precision frontier for estimation and honest inference at a latent regression discontinuity cutoff. Treatment assignment records which side of the latent cutoff an individual occupies, while the researcher observes a score contaminated by Gaussian measurement error. On a specified bounded-support benchmark, we characterize both worst-case expected absolute estimation error and worst-case expected length of connected confidence intervals with uniform \(90\%\) coverage. The comparison holds simultaneously for every sample size \(n\ge2\), where \(n\) denotes the number of independent observations, and every known noise standard deviation \(\sigma\in[0,1]\), with constants independent of sample size and noise scale.

The experiment observes treatment, a binary outcome, and the contaminated score. Assignment follows the latent score, and the target is the difference between the two conditional potential-outcome means at the latent cutoff. Recorded assignment therefore separates the sample into two true-side arms, each of which presents an endpoint estimation problem. This structure connects continuity-based regression discontinuity \citep{Hahn2001,CattaneoTitiunik2022Survey} to statistical recovery under measurement error. In the closely related observed-treatment experiment, \citet[Assumptions 1Y, 2C, and 7; Proposition 4(b)]{PeiShen2017DevilTails} establish identification under Gaussian measurement error and local density positivity. Our question concerns the precision available from a finite sample under quantitative restrictions on the latent law.

Those restrictions define the domain of the matched frontier. The latent score has known support \([-1,1]\) and a density continuous on that interval, bounded everywhere between \(1/4\) and \(3/4\). Both potential outcomes are binary. Their specified continuous conditional means also lie between \(1/4\) and \(3/4\), and each has full-interval Hölder seminorm at most two for a known exponent \(\beta\in(0,1]\), the smoothness parameter. The standardized Gaussian error is independent of the latent score and both potential outcomes. The error distribution, its scale, the support, and the smoothness exponent are public inputs. These conditions specify the benchmark in \cref{obj:def:model}; the uniform estimation and coverage criteria range over every law in that class.

Our first result constructs an observable estimator and a finite-sample honest interval on this benchmark. Positive polynomial weights concentrate near the cutoff within each reflected treatment arm. A finite inverse-heat transform corrects these weights for the known Gaussian contamination. Outcome-marked and unmarked sample averages then estimate the corresponding latent weighted averages, and stabilized, projected ratios recover the arm means. A public radius combines a Hölder approximation bound with a Gaussian second-moment bound. Deterministic minimization of that radius over a finite candidate set selects the estimator and interval. \Cref{obj:thm:finite-certificate} establishes uniform coverage of at least \(0.9\), bounds worst expected absolute error by the selected radius, and bounds worst expected interval length by twice that radius.

Our second result shows that this construction attains the optimal order for both decision criteria. \Cref{obj:thm:uniform-frontier} characterizes the common rate through a unique scalar resolution equation, with comparison constants depending only on the fixed Hölder exponent. Matching lower bounds concern the complete observed experiment and apply to every measurable estimator and every admissible honest connected interval, including procedures with independent randomization. The witnesses combine full-interval Hölder regularity, separation of the cutoff contrast, and polynomial cancellation under Gaussian convolution.

The resolution equation distinguishes three regimes. When the noise scale is at most \(n^{-1/(2\beta+1)}\), both criteria have the direct-experiment order
\[
n^{-\beta/(2\beta+1)}.
\]
On the intermediate branch specified in \cref{obj:thm:uniform-frontier}, their common order is
\[
\frac{\sigma^\beta}
{[1+\log(n\sigma^{2\beta+1})]^{3\beta/2}}.
\]
For each fixed positive noise scale, the compact-support branch eventually gives
\[
\left\{\frac{\log\log(en)}{\log(en)}\right\}^{2\beta},
\]
with comparison constants in this fixed-noise specialization permitted to depend on that noise level. The finite-sample root characterization covers both interfaces, zero noise, and every shrinking-noise sequence. For example, with Lipschitz conditional means and \(\sigma=n^{-1/6}\), the intermediate branch applies for every \(n\ge2\).

The construction draws on compact-support deconvolution and Gaussian moment inversion. \citet[Theorem~1(a), fixed-support case]{Meister2007CompactSupport} analyzes density recovery under known compact support, while \citet[Theorem~1; Section~4.1]{WuYang2020DenoisedMoments} uses Gaussian polynomial inversion to estimate finite mixing distributions. Here polynomial inversion targets positive endpoint averages that are converted into conditional means through arm ratios. The noise-scale comparison is also related to shrinking-noise density estimation \citep[Theorems~2.1 and 2.2]{Delaigle2008ShrinkingNoise} and distribution estimation under Wasserstein loss \citep[Section~2, Theorems~2.4--2.6]{DurotMukherjee2026VanishingNoise}. Our target and losses give these comparisons their specific decision-theoretic meaning. The honest interval construction follows the smoothness-based principle of combining worst-case approximation bias with stochastic uncertainty \citep{ArmstrongKolesar2018OptimalInference,ArmstrongKolesar2020SimpleHonest}.

The benchmark has two further interpretations. First, \cref{obj:prop:published-class-converse-transfer} embeds it into a larger Gaussian class with local density positivity and continuous conditional means near the cutoff, preserving both the target and the randomized observed experiment. This embedding transfers the benchmark lower bounds to the larger class. Second, a prospective eligibility-effect design links administrative true-side assignment, a contaminated eligibility score, and a prespecified binary outcome. Under externally justified Gaussian calibration, a known bounded latent range, independent sampling, and the benchmark's global density and conditional-mean restrictions, the results supply a precision benchmark for that design.

Supporting details for the main results appear in the appendix containing \cref{obj:thm:direct-reduction,obj:thm:uniform-frontier-resolution}.

The paper proceeds through related work; the experiment, assumptions, and decision criteria; observable estimation and finite-sample inference; uniform minimax rates across noise scales; and interpretation and extensions. The appendices develop the polynomial construction, observed-experiment lower bounds, rate comparisons and verification scope, and proofs of the main results.


\section{Related work}\label{sec:related-work}

The relationship between regression discontinuity with measurement error and the deconvolution literature depends on the assignment mechanism, the target, and the criterion used to evaluate statistical procedures. Our \leanref{S-1}{observed-treatment experiment} records the side of the latent cutoff and a Gaussian-contaminated score. The benchmark in \cref{obj:def:model} combines known compact support, quantitative density bounds, binary outcomes, and a fixed Hölder bound on the potential-outcome means. Its \leanref{S-2}{decision problem} evaluates estimators by worst-case expected absolute error and connected closed confidence intervals by worst-case expected length, subject to coverage of at least \(0.9\) for every benchmark law at each sample size. Three comparisons clarify how these criteria relate to existing results.

\subsection{Identification and cutoff decision criteria}

The closest identification comparison is \citet[Assumptions 1Y, 2C, and 7; Proposition 4(b)]{PeiShen2017DevilTails}. Their result identifies the sharp regression discontinuity effect from the observed outcome, treatment status, and mismeasured assignment variable under nondifferential measurement error, positive latent-score density near the cutoff, and Gaussian error. Observed treatment supplies information about the latent cutoff side, while Gaussian convolution permits recovery of the relevant latent distributions. Identification concerns recovery from the population observable law. The present benchmark evaluates the precision available from a finite sample under explicit quantitative restrictions. In particular, the known noise scale and the density and smoothness conditions in \cref{obj:def:model} specify the class over which estimation risk and interval length are compared.
% [Pei and Shen, Proposition 4(b)](https://arxiv.org/pdf/1609.01396)

A complementary inference comparison comes from \citet[Chapter~2, Assumptions~8, 9, 11, and 13; Theorem~6; printed pp.~49, 59, and 61]{Dong2018Thesis}. For observed treatment and a known error density, Theorem~6 establishes asymptotic normality of the deconvolution local-linear cutoff estimator after centering by its bias and standardizing by its variance. For Gaussian error, the prescribed bandwidth has order \((\log n)^{-1/2}\). The theorem also imposes kernel and transform-integrability conditions and a Lyapunov condition; Assumption~11 includes twice differentiable latent density and one-sided regression smoothness near the cutoff. These conditions support a pointwise distributional approximation for local-linear inference. Our criteria instead place the supremum over the specified benchmark class inside the risk and coverage requirements. The finite-sample guarantee in \cref{obj:thm:finite-certificate} and the rate comparison in \cref{obj:thm:uniform-frontier} address these uniform criteria under the benchmark's density bounds and Hölder restrictions. This theorem-level comparison relies on the cited-source attestation; its provenance limitation is recorded in the verification appendix, and it has no proof-use in our results.
% [Hao Dong, Essays in microeconometrics, Chapter 2](https://etheses.lse.ac.uk/3756/1/Dong__essays-in-microeconometrics.pdf)

Assignment and target distinctions also matter within the regression discontinuity literature. \citet{Eckles2025} exploit known exogenous noise in the score used to assign treatment. Their approach uses the resulting assignment randomness to estimate weighted treatment effects associated with changes in the assignment rule. In our experiment, assignment follows the latent score, and measurement error enters the score available to the researcher. The target is the potential-outcome contrast at the latent cutoff. Auxiliary information provides another route: \citet{DaveziesLeBarbanchon2017ContinuousError} study a two-sided fuzzy design with nondifferential measurement error and an auxiliary sample of treated individuals containing both true and noisy scores. Their target is a local average treatment effect, and their auxiliary observations identify features of the measurement process.
% [Eckles et al., Noise-induced randomization](https://arxiv.org/html/2004.09458)
% [Davezies and Le Barbanchon, Continuous measurement error](https://docs.iza.org/dp10801.pdf)

Observed-score interpretations supply a further comparison. Under correct classification of assignment and continuity of the potential-outcome means conditional on the observed score, \citet[Section~2.1, Lemma~2.1]{DongKolesar2023IgnoreMeasurementError} identify the treatment effect for units at the observed-score cutoff. Their additional condition on treatment-effect dependence on measurement error gives a weighted-average interpretation over latent scores. These conditions align conventional regression discontinuity analysis with an observed-score target. The latent-cutoff benchmark specifies a different conditioning event and uses the recorded treatment side together with the proxy distribution to recover its endpoint contrast.
% [Dong and Kolesár, Section 2.1](https://www.princeton.edu/~mkolesar/papers/rd_rounded.pdf)

\subsection{Compact support, boundary density estimation, and moments}

Compact-support deconvolution provides the second decision-problem comparison. \citet[Theorem~1(a), fixed-support case]{Meister2007CompactSupport} obtain uniform mean integrated squared error of order \((\log n/\log\log n)^{-2s}\) for density estimation over a compactly supported Fourier--Sobolev class of smoothness \(s\). Their procedure uses a known bound on the support and knowledge of the error characteristic function on a neighborhood of zero, where it is bounded away from zero. Thus, compact support supplies information that can improve density recovery even under partial knowledge of the error distribution. Our benchmark uses known support and a fully specified Gaussian measurement process, with Hölder smoothness imposed on conditional outcome means. The loss concerns an endpoint causal contrast, while the score density enters as a nuisance subject to continuity and quantitative bounds. These rates characterize different decision problems: integrated density recovery under Fourier--Sobolev restrictions and endpoint contrast recovery under conditional-mean Hölder restrictions.
% [Meister, Theorem 1(a)](https://www.f08.uni-stuttgart.de/.content/media/downloads/Mathematik/mathematische_berichte/2005/2005-007.pdf)

Boundary deconvolution and Gaussian moment methods clarify the constructional relationship. \citet{ZhangKarunamuni2009Boundary} study a density supported on a half-line under known independent additive error. They construct boundary deconvolution kernels and analyze pointwise mean squared error for density and density-derivative estimation under ordinary-smooth and supersmooth errors; their construction can also be adapted to compact support. This is a direct endpoint antecedent, but its target is the latent density rather than a marked-to-unmarked regression ratio, and its criterion does not include honest confidence-interval length. \citet[Theorem~1; Section~4.1]{WuYang2020DenoisedMoments} estimate finite Gaussian mixing distributions under Wasserstein loss, using bounded component locations and a specified number of components; their theorem distinguishes known and unknown common variance. For known variance, Hermite polynomials supply unbiased latent-moment estimates, which their procedure projects onto a feasible moment space before reconstructing the mixing distribution. The finite inverse-heat transform in \cref{obj:def:inverse-heat} uses the related polynomial inversion principle to estimate endpoint-weighted averages. The arm ratios in \cref{obj:def:ratio} then convert marked and unmarked averages into conditional-mean estimates. This connects Gaussian moment inversion to the cutoff target through an explicitly defined observable procedure.
% [Zhang and Karunamuni, Boundary deconvolution](https://doi.org/10.1016/j.jspi.2008.10.021)
% [Wu and Yang, Theorem 1 and Section 4.1](https://arxiv.org/html/1807.07237)

Deconvolution regression gives a closer antecedent for the ratio structure. \citet[Section~2, Theorems~2 and 5--6]{Fan1993} observe a response together with an error-contaminated covariate and estimate the latent regression function by normalizing a deconvolved response numerator by a deconvolved density denominator. They assume a known error law with nonvanishing characteristic function, a latent density bounded away from zero on a compact interval, and smoothness of both the density and regression function. For a supersmooth error of order \(\beta_e\), their uniform squared-risk upper rate is \(O((\log n)^{-2k/\beta_e})\); Gaussian error has \(\beta_e=2\), yielding the displayed convergence rate \((\log n)^{-k/2}\) in root mean square, with matching minimax results under their stated classes. That result concerns regression at points in a fixed interval with a fixed error law. The present target combines two recorded treatment arms at the latent support endpoint and evaluates expected absolute error and honest interval length uniformly over the full range of \(\sigma\).
% [Fan and Truong, Sections 2--3](https://archive.ymsc.tsinghua.edu.cn/pacm_download/468/10828-Jianqing_Fan_37.pdf)

Convolution-model linear functionals provide a complementary decision-theoretic comparison. \citet[Theorem~3.1 and Section~4.1]{Butucea2009} estimate \(\langle\psi,g\rangle\) from noisy observations under a known convolution density whose Fourier transform is nowhere zero. Their adaptive oracle bound balances the squared Fourier-tail bias against a cutoff-dependent penalty, and their pointwise-density specialization recovers the usual pointwise deconvolution regimes. \citet[Theorems~1--2 and Section~2.5]{Pensky2017} give upper and lower quadratic-risk bounds for general linear functionals over Sobolev signal classes with ordinary- or supersmooth Fourier behavior; the bounds match up to constants in the cases identified there and otherwise may differ by a logarithmic factor. Pointwise density is one of their applications. These results locate the inverse-problem difficulty for a single linear functional. Our endpoint estimand is instead assembled from four such marked and unmarked quantities through stabilized ratios, and the benchmark asks simultaneously for all-noise-scale absolute-risk and honest-length guarantees.
% [Butucea and Comte, Theorem 3.1 and Section 4.1](https://helios2.mi.parisdescartes.fr/~comte/ButuceaComte.pdf)
% [Pensky, Theorems 1--2 and Section 2.5](https://arxiv.org/pdf/1411.1660)

\subsection{Vanishing noise and the estimation target}

Shrinking-noise density estimation supplies the third comparison. \citet[Theorems~2.1 and 2.2]{Delaigle2008ShrinkingNoise} analyze mean integrated squared error for deconvolution density estimators as sample size increases and measurement-error scale decreases. Under their density-smoothness, kernel, and error assumptions, sufficiently small noise permits the corresponding error-free density rate; larger noise introduces a scale-dependent cost. This makes the comparison between noise scale and statistical resolution central to density estimation. For the present endpoint criterion, \cref{obj:def:uniform-rate,obj:thm:uniform-frontier} give the corresponding rate characterization across the benchmark noise scales, with constants uniform in sample size and noise scale at fixed benchmark smoothness.
% [Delaigle, Theorems 2.1 and 2.2](https://www3.stat.sinica.edu.tw/statistica/oldpdf/A18n311.pdf)

A recent minimax comparison is \citet[Section~2, Theorems~2.4--2.6]{DurotMukherjee2026VanishingNoise}. Their deconvolution analysis estimates a signal distribution under Wasserstein-one loss over a moment-bounded class, with known error distribution and known noise scale. Under Gaussian noise with standard deviation \(\sigma_n\), Theorem~2.4 gives the lower bound order \(\sigma_n(\log n)^{-1/2}+n^{-1/2}\). Their upper bounds match when \(\sigma_n\ge n^{-1/2+\eta}\), for fixed \(\eta>0\), and when \(\sigma_n\le n^{-1/2}\); the intervening regime carries separately stated logarithmic bounds under additional reciprocal-characteristic-function regularity. Our comparison concerns absolute error in a latent-cutoff contrast and expected length under uniform finite-sample coverage, so these distribution-recovery rates provide a regime comparison rather than a numerical benchmark for our loss.
% [Durot and Mukherjee, Section 2](https://arxiv.org/pdf/2404.09306)

\subsection{Honest inference under smoothness restrictions}

The interval criterion belongs to the optimal-recovery and smoothness-based regression literature. \citet{Donoho1994} connect statistical estimation over restricted function classes with optimal recovery. \citet{ArmstrongKolesar2018OptimalInference} derive finite-sample optimal confidence procedures for linear functionals in regression models with normal errors and known variance, together with uniform asymptotic extensions. \citet{ArmstrongKolesar2020SimpleHonest} construct honest intervals by incorporating worst-case smoothing bias into the critical value and bandwidth choice. Our finite-sample certificate similarly combines a public bias bound with a stochastic-error bound, here for inverse-heat averages and stabilized arm ratios. Coverage is imposed through \cref{obj:def:honest-class}, and precision is evaluated by worst expected length over that same class. These choices give the endpoint estimation and inference results a common decision-theoretic criterion.
% [Donoho, Statistical estimation and optimal recovery](https://digicoll.lib.berkeley.edu/record/85850/files/214.pdf)
% [Armstrong and Kolesár, Optimal inference](https://arxiv.org/abs/1511.06028)
% [Armstrong and Kolesár, Simple and honest confidence intervals](https://doi.org/10.3982/QE1199)


\section{Experiment, assumptions, and decision criteria}

Throughout, \(\lvert\cdot\rvert\) denotes absolute value and \(\|\cdot\|\) denotes the Euclidean norm. Generic positive constants, denoted by \leanref{sym:C}{\ensuremath{C}}, may depend on the fixed smoothness parameter unless another dependence is stated; their values may change between displays. We write \(a\lesssim b\) when \(a\le Cb\), and \(a\asymp b\) when both inequalities hold. The public smoothness parameter \leanref{sym:beta}{\ensuremath{\beta}} lies in \((0,1]\), the known noise scale \leanref{sym:sigma}{\ensuremath{\sigma}} lies in \([0,1]\), and the sample size \leanref{sym:n}{\ensuremath{n}} is an integer at least two. Treatment indices \leanref{sym:d}{\ensuremath{d}} range over \(\{0,1\}\), and sample indices \leanref{sym:i}{\ensuremath{i}} range from \(1\) to \(n\).

\subsection{Latent and observed experiments}

The experiment combines a latent running variable, binary potential outcomes, and an additive measurement error. Treatment follows the latent cutoff side, and the realized assignment is recorded alongside the outcome and the contaminated score. We first specify the latent law and the versions of its density and conditional means.

\begin{assumptionv}{ass:gaussian}[Gaussian measurement error]
Under \(P\), the ambient latent law, the standardized measurement error \(\varepsilon\) satisfies
\[
\varepsilon\sim N(0,1).
\]
\end{assumptionv}

Under the ambient latent law \leanref{sym:P}{\ensuremath{P}}, the running variable \leanref{sym:X}{\ensuremath{X}} has known support \([-1,1]\), the pair \leanref{sym:Ypot}{\ensuremath{(Y^0,Y^1)}} contains the untreated and treated potential outcomes, and \leanref{sym:epsilon}{\ensuremath{\varepsilon}} is the standardized measurement error. The density \leanref{sym:f}{\ensuremath{f(P,\cdot)}} and conditional-mean versions \leanref{sym:mu}{\ensuremath{\mu_d(P,\cdot)}} are part of this specification. Their continuous versions give a precise meaning to evaluation at the cutoff and at the support endpoints.

\begin{assumptionv}{ass:error-independence}[Independent measurement error]
Under \(P\), the measurement error \(\varepsilon\) is independent of the latent running variable \(X\) and the binary potential outcomes \(Y^0,Y^1\):
\[
\varepsilon\perp(X,Y^0,Y^1).
\]
\end{assumptionv}

The observation map retains the true-side assignment \leanref{sym:D}{\ensuremath{D}}, selects the observed outcome \leanref{sym:Y}{\ensuremath{Y}} from the potential outcomes, and reports the proxy score \leanref{sym:W}{\ensuremath{W}}:
\[
D=\mathbf1\{X\ge0\},\qquad
Y=(1-D)Y^0+DY^1,\qquad
W=X+\sigma\varepsilon.
\]
We consistently order the observed triple as
\[
O=(W,D,Y),\qquad
P_{\mathrm{obs}}=P\circ O^{-1}.
\]
The recorded assignment identifies the treatment arm of each observation, while the proxy supplies information about its latent distance from the cutoff. The causal target \leanref{sym:theta}{\ensuremath{\theta(P)}} is
\[
\theta(P)=\mu_1(P,0)-\mu_0(P,0).
\]
Continuity connects this contrast to the arm-specific limits at zero. Statistical procedures use an independent observed sample \leanref{sym:Sample}{\ensuremath{S_n}} and may randomize through an independent uniform seed \leanref{sym:U}{\ensuremath{U}}. Their joint law is
\[
\mathbb P_{P,n}=P_{\mathrm{obs}}^{\otimes n}\otimes\operatorname{Unif}[0,1],
\qquad
S_n=(O_1,\ldots,O_n),\quad U\sim\operatorname{Unif}[0,1].
\]
The seed permits randomized estimators and interval procedures within a common decision class. Deterministic procedures are included by taking their outputs to be constant in the seed.

\subsection{Law restrictions and benchmark}

The benchmark imposes five restrictions on the latent experiment: two determine the measurement-error structure, two give uniform density and mean bounds, and one controls the smoothness of the potential-outcome means. The Gaussian distribution and joint-independence conditions were stated with the observation map above. We now fix the ambient latent law and its continuous conditional-mean versions before imposing the quantitative density, mean, and Hölder restrictions.

\begin{definitionv}{synth_2}[Continuous conditional potential-outcome means]
For a latent law \(P\) with score \(X\in[-1,1]\) and binary potential outcomes \(Y^0,Y^1\), let \(\mu_d(P,\cdot)\) be a specified continuous version on \([-1,1]\) of the conditional mean of \(Y^d\) given \(X\), for \(d\in\{0,1\}\). Thus
\[
\mu_d(P,x)=\mathbb E_P[Y^d\mid X=x],
\qquad x\in[-1,1],
\]
with the conditional expectation interpreted through this specified continuous version, including at the endpoints and at the cutoff \(x=0\).
\end{definitionv}

The Gaussian classical measurement-error condition supplies a calibrated error distribution, with the public scale determining the magnitude of contamination \citep{PeiShen2017DevilTails}. Independence specifies how this error is coupled to the latent score and potential outcomes.

\begin{definitionv}{synth_1}[Latent law and continuous versions]
An ambient latent law \(P\) is a probability law on \(\mathbb R\times\{0,1\}\times\{0,1\}\times\mathbb R\), with coordinates \((X,Y^0,Y^1,\varepsilon)\), together with a density \(f(P,\cdot):\mathbb R\to\mathbb R\) and conditional potential-outcome mean versions \(\mu_d(P,\cdot):\mathbb R\to\mathbb R\), \(d\in\{0,1\}\), satisfying:
\begin{itemize}
\item \textbf{(Latent interval.)} \(P(X\in[-1,1])=1\).
\item \textbf{(Score density.)} The function \(f(P,\cdot)\) is continuous on \([-1,1]\), nonnegative on \(\mathbb R\), and zero outside \([-1,1]\). Moreover,
\[
\int_{-1}^{1} f(P,x)\,dx=1,
\qquad
P(X\in B)=\int_B f(P,x)\,dx
\]
for every measurable set \(B\subseteq\mathbb R\).
\item \textbf{(Continuous mean versions.)} For each \(d\in\{0,1\}\), the function \(\mu_d(P,\cdot)\) is continuous on \([-1,1]\), and for every measurable set \(B\subseteq\mathbb R\),
\[
\int_{\{X\in B\}} Y^d\,dP
=
\int_{B\cap[-1,1]} \mu_d(P,x)f(P,x)\,dx.
\]
\end{itemize}
\end{definitionv}

The joint-independence condition in \cref{obj:ass:error-independence} is specific to this analysis. It makes the same calibrated error distribution apply across latent scores and potential-outcome realizations, so the observation map has a common additive contamination mechanism. We next fix uniform bounds on the density throughout the known support.

\begin{assumptionv}{ass:density-bounds}[Latent density bounds]
The latent-score density \(f(P,x)\) satisfies
\[
\frac14\le f(P,x)\le\frac34
\qquad\text{for every }x\in[-1,1].
\]
\end{assumptionv}

The density bounds are specific to this analysis: they maintain positive latent-score density across the entire support and control its maximum. Continuity is imposed on the latent interval through \cref{obj:synth_1}; together, these requirements specify both the domain and the amount of score mass available near either side of the cutoff. A separate restriction keeps the binary conditional means uniformly interior.

\begin{assumptionv}{ass:mean-bounds}[Potential-outcome mean bounds]
The conditional potential-outcome means \(\mu_d(P,x)\) satisfy
\[
\frac14\le\mu_d(P,x)\le\frac34
\qquad\text{for every }d\in\{0,1\}\text{ and }x\in[-1,1].
\]
\end{assumptionv}

These mean bounds are also specific to this analysis. They keep each conditional success probability away from zero and one throughout the latent interval and imply \(\theta(P)\in[-1/2,1/2]\). Smoothness is specified separately through the full-interval Hölder seminorm \leanref{sym:Hseminorm}{\ensuremath{[\mu_d(P,\cdot)]_\beta}}, which measures variation between every pair of latent-score values.

\begin{definitionv}{synth_3}[Full-interval Hölder seminorm]
For \(0<\beta\le1\) and a continuous function \(v:[-1,1]\to\mathbb R\), define its full-interval Hölder seminorm by
\[
[v]_\beta
=
\sup_{\substack{x,t\in[-1,1]\\x\ne t}}
\frac{|v(x)-v(t)|}{|x-t|^\beta}.
\]
In particular, for \(d\in\{0,1\}\),
\[
[\mu_d(P,\cdot)]_\beta
=
\sup_{\substack{x,t\in[-1,1]\\x\ne t}}
\frac{|\mu_d(P,x)-\mu_d(P,t)|}{|x-t|^\beta}.
\]
The supremum is taken over the entire latent-score interval and may equal \(+\infty\).
\end{definitionv}

Using this seminorm, the smoothness restriction fixes a common radius for both potential-outcome mean functions.

\begin{assumptionv}{ass:mean-holder}[Potential-outcome Hölder bounds]
The full-interval Hölder seminorms \([\mu_d(P,\cdot)]_\beta\) satisfy
\[
[\mu_d(P,\cdot)]_\beta\le2
\qquad\text{for every }d\in\{0,1\}.
\]
\end{assumptionv}

The standard Hölder-ball restriction gives a fractional-order extension of the Lipschitz smoothness condition used in honest inference \citep{ArmstrongKolesar2020SimpleHonest}. Here the maintained radius is two, and the restriction applies on the entire latent interval, including pairs of points on opposite sides of the cutoff. Thus the density is continuous and globally bounded on its support, while the potential-outcome means satisfy an explicit full-domain Hölder bound. These conditions define the exact benchmark class.

\begin{definitionv}{def:model}[Latent-law class \(\mathcal M_\beta(\sigma)\)]
\(\mathcal M_\beta(\sigma)\) is the class of individual latent laws \(P\) satisfying \cref{obj:ass:gaussian,obj:ass:error-independence,obj:ass:density-bounds,obj:ass:mean-bounds,obj:ass:mean-holder} and the following conditions:
\begin{itemize}
\item \textbf{(Latent support.)} \(P\) is a probability law on
\[
[-1,1]\times\{0,1\}^{2}\times\mathbb R
\]
with coordinates \((X,Y^0,Y^1,\varepsilon)\).
\item \textbf{(Density continuity.)} The latent-score density \(f(P,\cdot)\) is continuous on \([-1,1]\).
\item \textbf{(Density normalization.)}
\[
\int_{-1}^{1}f(P,x)\,dx=1.
\]
\end{itemize}
The Gaussian-contaminated observed score \(W\) is
\[
W=X+\sigma\varepsilon.
\]
For every integer \(n\ge2\), the independent observed sample \(S_n\) has distribution
\[
S_n\sim P_{\mathrm{obs}}^{\otimes n},
\]
where \(P_{\mathrm{obs}}\) denotes the induced observed law.
\end{definitionv}

The benchmark class \leanref{sym:Model}{\ensuremath{\mathcal M_\beta(\sigma)}} treats the Gaussian calibration, noise scale, latent support, and smoothness exponent as known inputs. Its density and conditional means vary across admissible laws. Fixing these public inputs makes the statistical comparison uniform over the latent laws satisfying the displayed restrictions, with the smoothness radius and numerical bounds held fixed.

\subsection{Estimation and honest inference}

For estimation, the decision class \leanref{sym:Estimators}{\ensuremath{\mathcal T_n}} consists of all measurable real-valued procedures on the randomized observed experiment. Its minimax absolute risk is
\[
R_\beta(n,\sigma)
=
\inf_{\widehat\theta\in\mathcal T_n}
\sup_{P\in\mathcal M_\beta(\sigma)}
\mathbb E_{P,n}\!\left[|\widehat\theta(S_n,U)-\theta(P)|\right].
\]
The infimum permits every measurable estimator, including procedures chosen using the public inputs. The supremum requires a single chosen procedure to control risk across the entire benchmark class. For inference, the corresponding requirement is uniform coverage. The honest connected-interval class \leanref{sym:Intervals}{\ensuremath{\mathcal I_{\beta,n,\sigma}}} imposes this requirement at each sample size.

\begin{definitionv}{def:honest-class}[Honest interval class $\mathcal I_{\beta,n,\sigma}$]
The class of honest connected intervals is
\[
\mathcal I_{\beta,n,\sigma}
=
\left\{
I:(S_n,U)\mapsto[\ell,\upsilon]\subset[-1,1]:
\ell\le\upsilon,\quad
I\text{ is measurable},\quad
\inf_{P\in\mathcal M_\beta(\sigma)}
\mathbb P_{P,n}\{\theta(P)\in I(S_n,U)\}\ge0.9
\right\}.
\]
Here \(S_n\) is the observed sample, \(U\) is the independent procedure seed, \(\mathbb P_{P,n}\) is their joint law under the latent law \(P\), \(\mathcal M_\beta(\sigma)\) is the benchmark law class, and \(\theta(P)\) is the cutoff treatment contrast. The endpoints \(\ell\) and \(\upsilon\) are the outputs of the interval procedure.
\end{definitionv}

Honesty requires coverage of at least \(0.9\) for every law in the benchmark, with probability taken over both sampling and procedure randomization. Precision is measured by
\[
\mathcal L_\beta(n,\sigma)
=
\inf_{I\in\mathcal I_{\beta,n,\sigma}}
\sup_{P\in\mathcal M_\beta(\sigma)}
\mathbb E_{P,n}\!\left[\operatorname{length}(I(S_n,U))\right].
\]
The interval \([-1,1]\) belongs to the honest class because it contains every admissible cutoff contrast. The length criterion therefore compares the precision of uniformly valid connected intervals over a nonempty decision class. Together, the minimax absolute risk \(R_\beta(n,\sigma)\) and minimax worst expected length \(\mathcal L_\beta(n,\sigma)\) place estimation and honest inference on the same latent-law benchmark, with common public calibration, support, and smoothness inputs.


\section{Observable estimation and finite-sample inference}

For the benchmark class in \cref{obj:def:model}, the \leanref{S-3}{public polynomial procedure} combines positive averaging near the latent cutoff with an observable correction for Gaussian measurement error. The recorded assignment identifies the latent cutoff side, allowing the two arms to be treated separately. A stabilized ratio then estimates each arm's endpoint mean, and deterministic comparisons of public radii select the effect estimate and confidence interval.

The construction begins with a positive integer \leanref{sym:L}{\ensuremath{L}}, the polynomial index. We use the Chebyshev polynomial \leanref{sym:Chebyshev}{\ensuremath{T_L}} with the following normalization.

\begin{definitionv}{synth_4}[Chebyshev polynomials]
For each nonnegative integer \(L\), let \(T_L\) be the Chebyshev polynomial of the first kind, characterized by
\[
T_L(\cos t)=\cos(Lt),
\qquad t\in\mathbb R.
\]
This normalization gives \(T_0(x)=1\), \(T_1(x)=x\), and \(T_L(1)=1\). Each \(T_L\) is evaluated as a polynomial on all of \(\mathbb R\); the endpoint-kernel construction uses integers \(L\ge1\).
\end{definitionv}

From this family, we construct the normalized endpoint weight \leanref{sym:Kernel}{\ensuremath{K_L}}. Its nonnegative values on the unit interval support latent averaging, while its polynomial extension permits evaluation at every observed proxy value.

\begin{definitionv}{def:endpoint-kernel}[Endpoint polynomial weight $K_L$]
The endpoint polynomial weight is
\[
K_L(x)
=
\frac{
\left(\dfrac{1-T_L(1-2x)}{2x}\right)^3
}{
\displaystyle\int_0^1
\left(\dfrac{1-T_L(1-2t)}{2t}\right)^3\,dt
}.
\]
Here \(T_L\) is the Chebyshev polynomial. Removable values are filled by continuity, and \(K_L\) is extended polynomially to every real argument.
\end{definitionv}

The normalization makes \(K_L\) integrate to one on \([0,1]\). The polynomial inequalities governing its concentration near zero are developed in \cref{obj:lem:positive-endpoint}. To specify the observable correction, first record the polynomial degree \leanref{sym:Degree}{\ensuremath{m_L}} and derivative convention.



Gaussian averaging acts on a polynomial through finitely many even derivatives. The inverse-heat polynomial \leanref{sym:Inverse}{\ensuremath{Q_{L,\sigma}}} reverses that averaging using the known measurement-error scale.

\begin{definitionv}{def:inverse-heat}[Inverse-heat polynomial $Q_{L,\sigma}$]
The inverse-heat polynomial is
\[
Q_{L,\sigma}(w)
=
\sum_{k=0}^{\lfloor m_L/2\rfloor}
\frac{(-\sigma^2/2)^k}{k!}\,
K_L^{(2k)}(w).
\]
Here \(m_L\) is the degree of the endpoint polynomial \(K_L\), \(\sigma\) is the measurement-error scale, and \(K_L^{(2k)}\) denotes its derivative of order \(2k\).
\end{definitionv}

Under \cref{obj:ass:gaussian,obj:ass:error-independence}, averaging this correction over the measurement error recovers the corresponding latent polynomial weight. At zero noise, the correction equals \(K_L\). For positive noise, its observable values may have either sign even though the recovered latent weight is nonnegative. The ratio construction therefore stabilizes the empirical denominator explicitly.

Both arms use the same endpoint weight after reflection onto the positive latent interval. The arm sign \leanref{sym:Sign}{\ensuremath{s_d}} fixes this convention.

\begin{definitionv}{synth_5}[Arm-reflection convention]
For the treatment-arm index \(d\in\{0,1\}\), define
\[
s_d=2d-1,
\qquad s_1=1,\qquad s_0=-1.
\]
With \(D=\mathbf1\{X\ge0\}\), the reflected latent score \(s_dX\) lies in \([0,1]\) on the event \(\{D=d\}\). The observable polynomial weight for arm \(d\) is evaluated at the reflected proxy \(s_dW\).
\end{definitionv}

The population marked average \leanref{sym:Amean}{\ensuremath{A_{d,L}(P)}} and unmarked average \leanref{sym:Bmean}{\ensuremath{B_{d,L}(P)}} describe the endpoint averaging that the observable correction targets. We define them before introducing their sample counterparts.

\[
A_{d,L}(P)
=\int_0^1 K_L(x)\,\mu_d(P,s_dx)\,f(P,s_dx)\,dx,
\qquad
B_{d,L}(P)
=\int_0^1 K_L(x)\,f(P,s_dx)\,dx.
\]

The ratio \(A_{d,L}(P)/B_{d,L}(P)\) averages the arm's conditional mean with a positive, density-adjusted endpoint weight. Gaussian correction makes these population quantities accessible through averages of the observed triple. The outcome-marked sample average \leanref{sym:Ahat}{\ensuremath{\widehat A_{d,L}}} estimates the numerator, and the unmarked sample average \leanref{sym:Bhat}{\ensuremath{\widehat B_{d,L}}} estimates its normalizing mass. Their stabilized ratio gives the projected arm estimate \leanref{sym:muhat}{\ensuremath{\widehat\mu_{d,L}}}. The algorithm lists the coordinates in the computational order \((D_i,Y_i,W_i)\), but their law is the consistently ordered observed triple \(O_i=(W_i,D_i,Y_i)\) defined above.

\begin{algorithmv}{def:ratio}[Projected arm estimate $\widehat\mu_{d,L}$]
Inputs are the observed sample \((D_i,Y_i,W_i)_{i=1}^n\), the arm \(d\), the index \(L\), the measurement-error scale \(\sigma\), the inverse-heat polynomial \(Q_{L,\sigma}\), and the arm-reflection sign \(s_d\).
\begin{enumerate}
\item Compute the outcome-marked average:
\[
\widehat A_{d,L}
=
\frac1n\sum_{i=1}^n
\mathbf1\{D_i=d\}Y_iQ_{L,\sigma}(s_dW_i).
\]
\item Compute the unmarked average:
\[
\widehat B_{d,L}
=
\frac1n\sum_{i=1}^n
\mathbf1\{D_i=d\}Q_{L,\sigma}(s_dW_i).
\]
\item Output the projected arm estimate:
\[
\widehat\mu_{d,L}
=
\Pi_{[1/4,3/4]}
\left(
\frac{\widehat A_{d,L}}
{\max\{\widehat B_{d,L},1/8\}}
\right).
\]
Here \(\Pi_{[1/4,3/4]}\) denotes ordinary projection onto \([1/4,3/4]\).
\end{enumerate}
\end{algorithmv}

The two averages retain normalization by the full sample size, so their ratio adjusts for the arm's weighted population mass. Their expectations are \(A_{d,L}(P)\) and \(B_{d,L}(P)\), respectively. The denominator floor makes the estimate well defined for every observed sample, including samples with small or negative corrected mass. Projection enforces the conditional-mean range in \cref{obj:ass:mean-bounds}.

To assemble the effect estimate and interval, we use two public quantities: the endpoint bias moment \leanref{sym:Moment}{\ensuremath{M_{L,\beta}}} and the derivative variance cost \leanref{sym:Variance}{\ensuremath{V_L(\sigma)}}. The former measures the displacement of the latent weight from the cutoff under the maintained Hölder exponent; the latter governs the sampling term in the certificate. The candidate effect estimate \leanref{sym:thetahat}{\ensuremath{\widehat\theta_L}}, public radius \leanref{sym:radius}{\ensuremath{\rho_L}}, and connected interval \leanref{sym:Ipoly}{\ensuremath{I_L}} are defined together, including their constant fallback values.

\begin{definitionv}{synth_6}[Candidate effect estimates, public radii, and intervals]
Fix public parameters \(0<\beta\le1\), an integer \(n\ge2\), and \(0\le\sigma\le1\). For each integer \(L\ge1\), let \(K_L\) be the endpoint kernel, put \(m_L=3(L-1)\), and define the public moment and covariance quantities
\[
M_{L,\beta}=\int_0^1x^\beta K_L(x)\,dx,
\qquad
V_L(\sigma)=\frac34\sum_{j=0}^{m_L}
\frac{\sigma^{2j}}{j!}
\int_0^1[K_L^{(j)}(x)]^2\,dx.
\]
Here \(K_L^{(j)}\) denotes the ordinary \(j\)th derivative, and the \(j=0\) coefficient is one, including at \(\sigma=0\). Using the projected arm estimates, define
\[
\widehat\theta_L=\widehat\mu_{1,L}-\widehat\mu_{0,L},
\qquad
\rho_L=12M_{L,\beta}
+28\sqrt{\frac{40V_L(\sigma)}{n}},
\]
and the connected closed interval
\[
I_L=
[\widehat\theta_L-\rho_L,\widehat\theta_L+\rho_L]
\cap[-1,1].
\]
The index-zero candidate is
\[
\widehat\theta_0=0,\qquad
\rho_0=\frac12,\qquad
I_0=\left[-\frac12,\frac12\right].
\]
Thus each radius is a deterministic function of \((\beta,n,\sigma)\), while each positive-index interval is centered at its observable candidate estimate before intersection with \([-1,1]\).
\end{definitionv}

The fallback interval contains the full target range implied by \cref{obj:ass:mean-bounds}. Positive-index candidates combine the two projected arm estimates and add a radius accounting for endpoint approximation and sampling variability. Intersecting with \([-1,1]\) preserves coverage of the admissible target and gives a connected closed interval. Polynomial derivative bounds and the Gaussian covariance calculation are presented in \cref{obj:lem:factorial-derivatives,obj:lem:exact-covariance}.

Selection takes place over a finite candidate set. Its public upper limit \leanref{sym:Lmax}{\ensuremath{L_{\max}}} is specified before observing the sample.

\[
L_{\max}=\left\lceil n^{1/(4\beta+2)}\right\rceil.
\]

The selected index \leanref{sym:Lstar}{\ensuremath{L_\star}} minimizes the public radius over this set, with the smallest index resolving ties. The selected certificate \leanref{sym:Certificate}{\ensuremath{E_\beta(n,\sigma)}} is that minimum radius.

\begin{algorithmv}{def:public-selection}[Public selection $L_\star$]
Inputs are \((\beta,n,\sigma)\), the finite index cap \(L_{\max}\), and the candidate radii \(\rho_L\), effect estimates \(\widehat\theta_L\), and intervals \(I_L\); index zero uses the specified constant estimate and the full target-range interval.
\begin{enumerate}
\item Select the smallest index minimizing the radius:
\[
L_\star
=
\min\operatorname*{arg\,min}_{L\in\{0,\ldots,L_{\max}\}}
\rho_L.
\]
All comparisons are deterministic functions of \((\beta,n,\sigma)\).
\item Define the selected radius:
\[
E_\beta(n,\sigma)=\rho_{L_\star}.
\]
\item Output the selected estimate and interval:
\[
\bigl(\widehat\theta_{L_\star},I_{L_\star}\bigr).
\]
\end{enumerate}
\end{algorithmv}

The displayed public cap and selection rule define the procedures used in \cref{obj:thm:uniform-frontier}. For fixed public inputs, the selected index is fixed across samples. Consequently, its coverage guarantee follows from the corresponding fixed-index guarantee, or from the full target-range fallback when index zero is selected. Including that fallback also ensures \(E_\beta(n,\sigma)\le1/2\). The complete procedure has the following finite-sample certificate.

\begin{theoremv}{thm:finite-certificate}[Finite sample certificate]
Suppose that:
\begin{itemize}
\item \textbf{(Smoothness.)} The public smoothness parameter satisfies \(0<\beta\le1\).
\item \textbf{(Noise.)} The noise scale satisfies \(0\le\sigma\le1\).
\item \textbf{(Sample and degree.)} The sample size \(n\) and polynomial index \(L\) are integers satisfying \(n\ge2\) and \(L\ge1\).
\item \textbf{(Latent law.)} \(P\in\mathcal M_\beta(\sigma)\).
\end{itemize}
Write
\[
\theta(P)=\mu_1(P,0)-\mu_0(P,0),
\qquad
B_{d,L}(P)=\mathbb E_{P,n}\widehat B_{d,L},
\]
where \(\mathbb P_{P,n}\) is the joint law of the independent observed sample and independent uniform procedure seed, and \(\mathbb E_{P,n}\) denotes expectation under this law. Define the endpoint bias moment and derivative variance cost by
\[
M_{L,\beta}
=\int_0^1 x^\beta K_L(x)\,dx,
\qquad
V_L(\sigma)
=\frac34\sum_{j=0}^{m_L}
\frac{\sigma^{2j}}{j!}
\int_0^1\bigl(K_L^{(j)}(x)\bigr)^2\,dx,
\]
where \(m_L\) is the degree of the endpoint polynomial \(K_L\).
For the effect estimate
\[
\widehat\theta_L=\widehat\mu_{1,L}-\widehat\mu_{0,L}
\]
and its connected closed interval \(I_L\), the following guarantees hold:
\[
B_{d,L}(P)\ge\frac14
\quad\text{for every }d\in\{0,1\},
\]
\[
\mathbb E_{P,n}\bigl|\widehat\theta_L-\theta(P)\bigr|
\le
12M_{L,\beta}
+28\sqrt{\frac{V_L(\sigma)}{n}},
\qquad
\mathbb P_{P,n}\{\theta(P)\in I_L\}\ge\frac9{10}.
\]
Moreover, the publicly selected interval \(I_{L_\star}\) belongs to the uniformly honest class \(\mathcal I_{\beta,n,\sigma}\) of \cref{obj:def:honest-class}, and the selected estimator and interval satisfy
\[
\sup_{P\in\mathcal M_\beta(\sigma)}
\mathbb E_{P,n}
\bigl|\widehat\theta_{L_\star}-\theta(P)\bigr|
\le E_\beta(n,\sigma),
\qquad
\sup_{P\in\mathcal M_\beta(\sigma)}
\mathbb E_{P,n}\operatorname{length}(I_{L_\star})
\le2E_\beta(n,\sigma).
\]
\end{theoremv}

The certificate separates endpoint approximation from sampling variability. Positive latent averaging and the density lower bound give each population denominator mass of at least \(1/4\), placing the empirical floor below that mass. Smoothness controls approximation through the endpoint moment, while the derivative variance cost accounts for the observable Gaussian correction. Together, these terms yield both an absolute-error bound and coverage of at least \(0.9\) at each positive index.

Public selection converts these candidate guarantees into a single uniformly honest procedure over the benchmark class. Its deterministic radius bounds worst expected absolute error, and twice that radius bounds worst expected interval length. Under the stated calibration, support, density, and conditional-mean restrictions, \cref{obj:thm:finite-certificate} therefore supplies a finite-sample comparison of candidate precision together with a coverage guarantee for the selected interval.


\section{Uniform minimax rates across noise scales}

The public certificate in \cref{obj:thm:finite-certificate} bounds both estimation error and honest interval length. We now characterize the optimal scale of these two criteria across sample sizes and calibrated noise levels. The characterization uses a scalar resolution equation that balances endpoint approximation against the cost of Gaussian inversion. Its definition also identifies the observable procedures that attain the resulting rate.

\begin{definitionv}{def:uniform-rate}[Uniform rate \(r_\beta(n,\sigma)\)]
The uniform resolution rate and selected observable procedures are
\[
r_\beta(n,\sigma)=h_\star(n,\sigma)^\beta,
\qquad
\widetilde\theta_{\beta,n,\sigma}=\widehat\theta_{L_\star},
\qquad
\widetilde I_{\beta,n,\sigma}=I_{L_\star},
\]
where \(L_\star\) is selected by the finite public minimization. With the local abbreviations
\[
\nu=2\beta+1,\qquad h_0=n^{-1/\nu},
\]
define the noise cost, for \(0<h\le1\), by
\[
\Psi(h,\sigma)=
\begin{cases}
0,&\sigma\le h,\\
(\sigma/h)^{2/3}-1,&\sigma^4\le h<\sigma,\\
h^{-1/2}\{1+\log(\sigma^2/\sqrt h)\}-1,&h<\sigma^4.
\end{cases}
\]
At \(\sigma=0\), the first branch applies throughout \(0<h\le1\). The resolution \(h_\star(n,\sigma)\) is specified as the root in \([h_0,1]\) of
\[
\log n+\nu\log h_\star(n,\sigma)
=\Psi(h_\star(n,\sigma),\sigma).
\]
The left side is continuous and increasing in the resolution, and the right side is continuous and nonincreasing. Root existence and uniqueness, risk bounds, and interval honesty are subsequent assertions. All tuning parameters belong to the procedures; the law class retains its defining restrictions.
\end{definitionv}

The noise cost \leanref{sym:Noisecost}{\ensuremath{\Psi(h,\sigma)}} summarizes Gaussian inversion cost for the rate comparison. The selected resolution \leanref{sym:Rateresolution}{\ensuremath{h_\star(n,\sigma)}} translates this cost into an endpoint approximation scale, and the rate \leanref{sym:Frontier}{\ensuremath{r_\beta(n,\sigma)}} is its Hölder power. The selected estimator \leanref{sym:Attainer}{\ensuremath{\widetilde\theta_{\beta,n,\sigma}}} and selected interval \leanref{sym:Honestattainer}{\ensuremath{\widetilde I_{\beta,n,\sigma}}} are the procedures chosen by \cref{obj:def:public-selection}. Thus the scalar equation describes the precision achieved by the finite public selection rule.

The two noisy branches come from distinct parts of the same endpoint mechanism. In the intermediate range, the noise cost grows as \((\sigma/h)^{2/3}-1\), linking finer endpoint resolution to a larger Gaussian inversion cost. Below the compact-support interface \(h=\sigma^4\), the cost becomes \(h^{-1/2}\{1+\log(\sigma^2/\sqrt h)\}-1\). These expressions agree at the interface. The balance equation therefore connects endpoint approximation to inversion cost continuously across the two noisy regimes, using the same resolution to compare both decision criteria.

On the benchmark class of \cref{obj:def:model}, the following result characterizes minimax absolute risk and minimax honest expected length simultaneously. Its comparison constants are uniform in both sample size and noise scale.

\begin{theoremv}{thm:uniform-frontier}[Uniform estimation and interval frontier]
Fix \(\beta\in(0,1]\). There exist constants \(c,C>0\), depending only on \(\beta\), such that the following conclusions hold simultaneously for every \(n\ge2\) and every \(\sigma\in[0,1]\).

Write \(\nu=2\beta+1\) and \(h_0=n^{-1/\nu}\). The equation defining the resolution \(h_\star(n,\sigma)\) has a unique root. With \(r_\beta(n,\sigma)=h_\star(n,\sigma)^\beta\), the selected observable estimator \(\widetilde\theta_{\beta,n,\sigma}\), the selected observable interval \(\widetilde I_{\beta,n,\sigma}\), and their public radius \(E_\beta(n,\sigma)\) satisfy
\[
c r_\beta(n,\sigma)
\le R_\beta(n,\sigma)
\le \sup_{P\in\mathcal M_\beta(\sigma)}
\mathbb E_{P,n}
\left|\widetilde\theta_{\beta,n,\sigma}-\theta(P)\right|
\le E_\beta(n,\sigma)
\le C r_\beta(n,\sigma)
\]
and
\[
c r_\beta(n,\sigma)
\le \mathcal L_\beta(n,\sigma)
\le \sup_{P\in\mathcal M_\beta(\sigma)}
\mathbb E_{P,n}
\operatorname{length}(\widetilde I_{\beta,n,\sigma})
\le 2E_\beta(n,\sigma)
\le C r_\beta(n,\sigma),
\qquad
\widetilde I_{\beta,n,\sigma}\in\mathcal I_{\beta,n,\sigma}.
\]
The minimax criteria range over all procedures in their stated decision classes, including procedures with arbitrary independent randomization and arbitrary degree.

The following branches exhaust every sample size. Throughout, \(\asymp\) denotes comparison by positive constants depending only on \(\beta\), except where dependence on an additional fixed exponent is stated.
\begin{itemize}
\item \textbf{(Direct branch.)}
If \(\sigma\le h_0\), then \(h_\star=h_0\), including at \(\sigma=0\).

\item \textbf{(Intermediate branch.)}
If \(\sigma>h_0\), then \(h_0<h_\star<\sigma\). Within this noisy region, the intermediate branch occurs exactly when
\[
\log(n\sigma^{4\nu})\le \sigma^{-2}-1.
\]
There is a unique \(z\in(1,\sigma^{-2}]\) satisfying
\[
z+\frac{3\nu}{2}\log z=1+\log(n\sigma^\nu),
\]
and
\[
h_\star=\sigma z^{-3/2}
\asymp
\frac{\sigma}{[1+\log(n\sigma^\nu)]^{3/2}}.
\]

\item \textbf{(Compact branch.)}
If \(\sigma>h_0\) and
\[
\log(n\sigma^{4\nu})>\sigma^{-2}-1,
\]
there is a unique \(\tau>\sigma^{-2}\) satisfying
\[
2\nu\log\tau+
\tau\{1+\log(\sigma^2\tau)\}
=\log n+1.
\]
On this branch,
\[
h_\star=\tau^{-2},
\qquad
\tau\asymp
\frac{\log(en)}{\log(e+\sigma^2\log(en))},
\qquad
r_\beta(n,\sigma)\asymp
\left\{
\frac{\log(e+\sigma^2\log(en))}{\log(en)}
\right\}^{2\beta},
\]
uniformly over the branch.
\end{itemize}
At the intermediate--compact equality, \(h_\star=\sigma^4\), and both scalar formulas give this value. At the direct--noisy equality, \(h_\star=\sigma=h_0\), and the adjacent formulas agree.

For every fixed \(a\) with \(0<a<1/(2\beta+1)\),
\[
r_\beta(n,n^{-a})
\asymp n^{-a\beta}(\log n)^{-3\beta/2}
\]
as \(n\) tends to infinity, with comparison constants that may additionally depend on \(a\). For every fixed \(a\ge1/(2\beta+1)\), the sequence \(\sigma=n^{-a}\) lies in the exact direct branch for every \(n\ge2\). For every fixed \(\sigma\in(0,1]\),
\[
r_\beta(n,\sigma)
\asymp
\left\{\frac{\log\log(en)}{\log(en)}\right\}^{2\beta}
\]
as \(n\) tends to infinity, with comparison constants allowed to depend on the fixed noise level. The root and finite-sample branch formulas characterize every shrinking-noise sequence.

In particular, for \(\beta=1\) and \(\sigma=n^{-1/6}\), the intermediate branch holds for every \(n\ge2\), and
\[
r_1(n,n^{-1/6})
=h_\star(n,n^{-1/6})
\asymp
\frac{n^{-1/6}}{[1+\tfrac12\log n]^{3/2}}.
\]
\end{theoremv}

\Cref{obj:thm:uniform-frontier} contains the exact branch conditions, both interfaces, and the relevant shrinking- and fixed-noise specializations. Its scalar balance connects endpoint approximation with Gaussian inversion cost, while the matched upper and lower bounds show that this one resolution governs both absolute risk and honest expected length. In the direct regime, the inversion cost vanishes at the direct resolution. In the intermediate regime, resolving the endpoint more finely than the noise scale incurs the fractional-power cost in \cref{obj:def:uniform-rate}. The compact-support regime uses the finer-resolution cost below \(\sigma^4\), yielding the logarithmic comparison in \cref{obj:thm:uniform-frontier}.

For reference, \cref{tab:noise-regimes} collects the exact branch conditions and resolutions. The auxiliary scalars \(z\) and \(\tau\) are the unique solutions specified in \cref{obj:thm:uniform-frontier}.

\begin{table}[htbp]
\centering
\caption{Exact resolution regimes for \(n\ge2\) and \(\sigma\in[0,1]\), with \(\nu=2\beta+1\) and \(h_0=n^{-1/\nu}\).}
\label{tab:noise-regimes}
\begin{tabular}{@{}p{0.17\linewidth}p{0.49\linewidth}p{0.25\linewidth}@{}}
\hline
Regime & Exact condition & Exact resolution \\
\hline
Direct &
\(\sigma\le h_0\) &
\(h_\star=h_0\) \\[0.6em]
Intermediate &
\(\sigma>h_0\) and
\(\log(n\sigma^{4\nu})\le\sigma^{-2}-1\) &
\(h_\star=\sigma z^{-3/2}\) \\[0.6em]
Compact support &
\(\sigma>h_0\) and
\(\log(n\sigma^{4\nu})>\sigma^{-2}-1\) &
\(h_\star=\tau^{-2}\) \\
\hline
\end{tabular}

\smallskip
\parbox{0.95\linewidth}{\small
At the direct--noisy interface, \(\sigma=h_0\) and \(h_\star=\sigma=h_0\).
At the intermediate--compact interface,
\(\log(n\sigma^{4\nu})=\sigma^{-2}-1\) within the noisy region, and
\(h_\star=\sigma^4\); both noisy scalar formulas agree there.
The common rate is \(h_\star^\beta\), and its comparison with both minimax criteria uses constants depending only on \(\beta\), uniformly in \(n\) and \(\sigma\).}
\end{table}

The uniform characterization also distinguishes asymptotic paths. For polynomially shrinking noise \(\sigma=n^{-a}\), a fixed exponent \(0<a<1/(2\beta+1)\) gives the order \(n^{-a\beta}(\log n)^{-3\beta/2}\), whereas \(a\ge1/(2\beta+1)\) places every sample size in the exact direct branch. With fixed positive noise, the rate has order \(\{\log\log(en)/\log(en)\}^{2\beta}\). The Lipschitz example \(\beta=1\), \(\sigma=n^{-1/6}\) sharpens this comparison: the intermediate condition holds for every \(n\ge2\), so its displayed formula describes the entire sequence. These are comparisons of rates; the exact resolution is determined by the scalar equation, and the minimax inequalities retain their comparison constants. Constants for the uniform branch comparisons depend only on \(\beta\), while those for a specified asymptotic path may also depend on its fixed exponent or fixed noise level as stated in \cref{obj:thm:uniform-frontier}.

At zero noise, endpoint approximation and sampling variability balance at \(h_0=n^{-1/(2\beta+1)}\), giving the Hölder power \(n^{-\beta/(2\beta+1)}\). The direct-experiment specialization in \cref{obj:thm:direct-reduction}, presented in the verification appendix, establishes this order for both criteria and its attainment by the selected estimator and honest interval. Together with the exact direct branch in \cref{obj:thm:uniform-frontier}, it shows how sufficiently precise score measurement recovers the direct-experiment rate throughout the region \(\sigma\le h_0\).


\section{Interpretation and extensions}

The rate characterization in \cref{obj:thm:uniform-frontier} describes the information available when assignment records the latent cutoff side, the latent score has known bounded support, and the Gaussian measurement-error distribution is calibrated. Recorded assignment separates the observations into the two latent-score arms. Within each arm, the cutoff is an endpoint, and the contaminated score supplies information about distance from that endpoint. These features give the endpoint weights in \cref{obj:def:endpoint-kernel} their statistical interpretation.

Known Gaussian calibration supplies the inverse-heat transform in \cref{obj:def:inverse-heat}. The finite-sample guarantees in \cref{obj:thm:finite-certificate} combine this correction with the global density and conditional-mean restrictions defining \cref{obj:def:model}. The density bounds control the mass entering the arm ratios, while the full-interval Hölder restrictions control their approximation error. Together, these conditions support observable estimation and honest inference uniformly over the benchmark, with matching minimax orders across the stated noise scales.

\subsection{Converse transfer to a larger Gaussian class}

The benchmark also supplies lower bounds for a larger class specified through local density and conditional-mean conditions. To express this implication precisely, the following definition fixes the larger class, its target, and its decision criteria immediately before the embedding result.



\begin{propositionv}{prop:published-class-converse-transfer}[Gaussian class converse transfer]
For \(\sigma\in[0,1]\), let \(\mathcal G(\sigma)\) consist of probability laws \(Q\) on \(\mathbb R\times\{0,1\}^2\times\mathbb R\), with coordinates \((Z,V^0,V^1,E)\), a score density \(a\), and conditional-mean versions
\[
v_d(z)=\mathbb E_Q[V^d\mid Z=z],\qquad d\in\{0,1\}.
\]
Define the assignment, observed outcome, and observed score by
\[
A=\mathbf 1\{Z\ge0\},\qquad
B=AV^1+(1-A)V^0,\qquad
R=Z+\sigma E.
\]
Membership requires the following conditions:
\begin{itemize}
\item \textbf{(Local density.)} There exists \(r>0\) such that \(a\) is continuous on \((-r,r)\) and \(a(z)>0\) for every \(z\in(-r,r)\).
\item \textbf{(Conditional means.)} Each \(v_d\) is a conditional-mean version continuous on an open neighborhood of zero, with \(0\le v_d(0)\le1\).
\item \textbf{(Gaussian error.)} \(E\sim N(0,1)\).
\item \textbf{(Joint independence.)} \(E\perp(Z,B)\).
\end{itemize}
The target and randomized experiment are
\[
\vartheta(Q)=v_1(0)-v_0(0),\qquad
\bigl(Q\circ(R,A,B)^{-1}\bigr)^{\otimes n}
 \otimes\operatorname{Unif}[0,1].
\]
Let \(\mathcal R_{\mathcal G}(n,\sigma)\) denote the infimum, over measurable estimates based on this experiment, of
\[
\sup_{Q\in\mathcal G(\sigma)}
\mathbb E_Q\!\left[\,|T-\vartheta(Q)|\,\right].
\]
Let \(\mathcal L_{\mathcal G}(n,\sigma)\) denote the infimum of
\[
\sup_{Q\in\mathcal G(\sigma)}
\mathbb E_Q[\upsilon-\ell]
\]
over measurable connected closed intervals \([\ell,\upsilon]\subseteq[-1,1]\) based on the same experiment and satisfying
\[
\inf_{Q\in\mathcal G(\sigma)}
Q\{\vartheta(Q)\in[\ell,\upsilon]\}\ge0.9.
\]
These expectations and coverage probabilities include the independent uniform procedure seed.

Fix \(\beta\in(0,1]\). For every \(\sigma\in[0,1]\) and every \(P\in\mathcal M_\beta(\sigma)\), the law \(P\), equipped with the density \(a=f(P,\cdot)\) and versions \(v_d=\mu_d(P,\cdot)\), belongs to \(\mathcal G(\sigma)\). This embedding preserves the target,
\[
\vartheta(Q)=\theta(P)
=\mu_1(P,0)-\mu_0(P,0),
\]
the observed law,
\[
Q\circ(R,A,B)^{-1}=P_{\mathrm{obs}},
\]
and, for every integer \(n\ge0\), the entire randomized experiment,
\[
\bigl(Q\circ(R,A,B)^{-1}\bigr)^{\otimes n}
 \otimes\operatorname{Unif}[0,1]
=\mathbb P_{P,n}.
\]

There exists \(c>0\), depending on \(\beta\), such that for every integer \(n\ge2\) and every \(\sigma\in[0,1]\),
\[
c\,r_\beta(n,\sigma)
\le R_\beta(n,\sigma)
\le \mathcal R_{\mathcal G}(n,\sigma),
\qquad
c\,r_\beta(n,\sigma)
\le \mathcal L_\beta(n,\sigma)
\le \mathcal L_{\mathcal G}(n,\sigma),
\]
where \(r_\beta(n,\sigma)\) is the benchmark frontier rate, and \(R_\beta(n,\sigma)\) and \(\mathcal L_\beta(n,\sigma)\) are the corresponding minimax absolute risk and minimax expected length under uniform \(0.9\) coverage on \(\mathcal M_\beta(\sigma)\).

For every \(\sigma\in(0,1]\) and every \(Q\in\mathcal G(\sigma)\), the locally defined binary-response Pei--Shen conditions hold: the standardized error is independent jointly of the latent score and observed outcome (Assumption 1Y), the score density is strictly positive on an open neighborhood of zero (Assumption 2C), and the scaled error satisfies
\[
\sigma E\sim N(0,\sigma^2)
\]
(Assumption 7).
\end{propositionv}

The intuition behind \cref{obj:prop:published-class-converse-transfer} is inclusion of statistical decision problems with exactly the same target and data distribution on the embedded benchmark. Any estimator evaluated over the larger class faces every benchmark law. Likewise, an interval honest over the larger class satisfies the benchmark coverage requirement. Preserving the independent procedure seed extends both comparisons to randomized procedures. The benchmark lower bounds therefore remain necessary precision bounds for the larger class.

The local positivity and continuity requirements in the larger Gaussian class of \cref{obj:prop:published-class-converse-transfer} give the cutoff contrast a well-defined meaning, while joint independence specifies a common Gaussian contamination mechanism for the latent score and observed outcome. For positive noise scales, that proposition connects these restrictions to Assumptions~1Y, 2C, and 7 of \citet{PeiShen2017DevilTails}. The resulting interpretation is an embedding-based converse for the larger Gaussian class. Matching attainment and finite-sample honesty have the exact benchmark in \cref{obj:def:model} as their domain, as established by \cref{obj:thm:finite-certificate,obj:thm:uniform-frontier}.

\subsection{A prospective eligibility-effect benchmark}

Following the Medicaid/CHIP setting in \citet[Chapter~2, Section~2.6]{Dong2018Thesis}, the proposed eligibility application would be organized around an eligibility cutoff.
% [Dong dissertation, Chapter 2, Section 2.6](https://etheses.lse.ac.uk/3756/1/Dong__essays-in-microeconometrics.pdf)
The proposed design would link administrative eligibility assignment to a contaminated score and a prespecified binary outcome. The score would be oriented so that eligibility corresponds to its nonnegative latent side. Potential outcomes would be indexed by eligibility assignment, giving the cutoff contrast the interpretation of an eligibility effect on the chosen outcome.

Applying the benchmark would require an administrative record that identifies the latent cutoff side, together with externally justified Gaussian calibration of the score error and its known scale. The proposed error mechanism would satisfy joint independence from the latent score and both potential outcomes. A prespecified population would have a known bounded latent range, normalized to \([-1,1]\) with the cutoff at zero, and observations would follow the independent sampling experiment in \cref{obj:def:model}.

The remaining proposed design conditions are substantive global restrictions: a continuous latent density between \(1/4\) and \(3/4\) throughout the normalized interval, continuous binary potential-outcome means between the same bounds, and full-interval Hölder seminorms at most two for a specified exponent in \((0,1]\). These are the restrictions in \cref{obj:ass:density-bounds,obj:ass:mean-bounds,obj:ass:mean-holder}. Under these maintained conditions, \cref{obj:thm:finite-certificate} supplies an observable interval with uniform finite-sample coverage, and \cref{obj:thm:uniform-frontier} characterizes its attainable precision. The application thus provides a prospective design benchmark conditional on administrative linkage, calibration, support, and global regularity.

\subsection{Limitations and future work}

The analysis treats the Gaussian error distribution and its scale as known. Estimated calibration would introduce uncertainty into both the inverse-heat weights and the public interval radius. Extending the coverage guarantee to that setting requires an experiment that includes the calibration data and conditions controlling its estimation error.

Smoothness is also a public input. Adaptation to an unknown exponent would require procedures and coverage arguments that account for selection across smoothness classes. Empirical assessment of benchmark membership presents a separate challenge: the support, density bounds, and global potential-outcome mean restrictions concern latent quantities. Administrative and validation data could inform their plausibility, but the prospective eligibility design remains conditional on these maintained restrictions.

Additional outcome models would broaden the application beyond binary responses. Such extensions require suitable outcome-tail or moment conditions and corresponding changes to the finite-sample second-moment bounds. The present results establish estimation and honest inference for the binary-outcome benchmark.

Finally, the focused comparison in \cref{sec:related-work} separates three antecedents: boundary kernels estimate a truncated latent density, deconvolution regression uses a numerator--denominator ratio at fixed error law, and convolution linear-functional theory gives quadratic-risk bounds for individual functionals. None of these cited results supplies the joint criterion used here: a recorded-arm endpoint contrast with absolute-risk and honest-length guarantees uniform over the benchmark noise scales. The comparison with \citet{ZhangKarunamuni2009Boundary} is qualitative, covering the boundary-density target, error classes, and pointwise mean squared error criterion described in the available primary-source material. The operational size of the explicit finite-sample certificate at applied sample sizes is likewise left for a separate calibrated numerical study; the present paper establishes its validity and optimal order without reporting a design-specific grid.


\appendix

\section*{Appendices}

\section{Polynomial construction and the finite-sample certificate}

The achievability argument for \cref{obj:thm:finite-certificate} combines endpoint concentration, control of polynomial derivatives, and exact Gaussian second moments. These ingredients connect the observable correction in \cref{obj:def:inverse-heat} to the stabilized ratios and public intervals in \cref{obj:def:ratio,obj:synth_6}. We present the auxiliary results in that mathematical order.

\subsection{Hermite polynomials and Gaussian averaging}

The Gaussian identities use the probabilists' Hermite polynomials. We fix their normalization before stating the orthogonality result.



This normalization determines the factorial weights in Gaussian polynomial second moments. The relevant classical identity is the following.

\begin{theoremv}{lem:classical-hermite-facts}[Gaussian Hermite orthogonality]
Let \(H_j\) be the probabilists' Hermite polynomials, defined by
\[
\exp(zt-t^2/2)
=\sum_{j=0}^{\infty}H_j(z)\frac{t^j}{j!}.
\]
For a standard Gaussian variable \(Z\) and nonnegative integers \(j,k\),
\[
\mathbb E[H_j(Z)H_k(Z)]
=j!\mathbf1\{j=k\},
\qquad H_0=1.
\]
\end{theoremv}

Orthogonality separates the contributions of distinct polynomial orders under Gaussian averaging. Consequently, a finite Hermite expansion has a second moment expressed as a sum of squared coefficients with factorial weights. This structure supplies the exact second-moment identity in \cref{obj:lem:exact-covariance}.

\subsection{Positive endpoint concentration}

The latent averaging step uses the endpoint polynomial weight in \cref{obj:def:endpoint-kernel}. Its positivity and concentration determine how accurately a weighted arm mean approximates the corresponding cutoff mean. The following bounds quantify that concentration and its squared-integral cost.

\begin{lemmav}{lem:positive-endpoint}[Positive endpoint kernel bounds]
Let \(0<\beta\le1\), and let \(K_L\) be the endpoint kernel defined in \cref{obj:def:endpoint-kernel}. Define its bias moment by
\[
M_{L,\beta}=\int_0^1 x^\beta K_L(x)\,dx.
\]
There exists a constant \(C>0\), depending on \(\beta\), such that for every integer \(L\ge1\),
\[
K_L(x)\ge0\quad\text{for all }x\in[0,1],
\qquad
\int_0^1 K_L(x)\,dx=1,
\]
and
\[
M_{L,\beta}\le C L^{-2\beta},
\qquad
\int_0^1 K_L(x)^2\,dx\le C L^2.
\]
\end{lemmav}

\begin{proof}[Proof of \cref{obj:lem:positive-endpoint}]
Fix \(0<\beta\le1\), and put
\[
c_{\mathrm{env}}:=\frac{\pi^6}{4}.
\]
We first establish positivity, unit mass, and the square-integral estimate with constant \(c_{\mathrm{env}}\), and then prove the moment estimate with a positive constant \(C_{\mathrm{bias}}\). The common constant will be
\[
C:=\max\{C_{\mathrm{bias}},c_{\mathrm{env}}\}.
\]
Since \(L^{-2\beta}\) and \(L^2\) are nonnegative, enlarging either constant to \(C\) preserves its respective bound.
% lean: CausalSmith.Stat.RdTruesideNoiseFrontier.positive_endpoint

\begin{enumerate}
\item Fix an integer \(L\ge1\). Define the polynomial quotient and its normalizer by
\[
F_L(x):=
\begin{cases}
\dfrac{1-T_L(1-2x)}{2x},&x\ne0,\\[6pt]
L^2,&x=0,
\end{cases}
\qquad x\in\mathbb R,
\qquad
c_{\mathrm{norm},L}:=\int_0^1 F_L(t)^3\,dt.
\]
Indeed, \(T_L(1)=1\), so the numerator is divisible by \(x\). The standard Chebyshev derivative identity gives
\[
T_L'(1)=L^2,
\qquad
2xF_L(x)=1-T_L(1-2x),
\qquad
F_L(0)=T_L'(1)=L^2,
\]
where the last equality also follows by differentiating the preceding polynomial identity at zero.

The normalization in \cref{obj:synth_4} implies, for \(0<x\le1\),
\[
T_L(1-2x)
=
\cos\!\bigl(L\arccos(1-2x)\bigr)
\le1.
\]
Thus \(F_L(x)\ge0\) on \((0,1]\), and continuity extends this inequality to zero. Consequently,
\[
F_L(x)\ge0\quad(0\le x\le1),
\qquad
c_{\mathrm{norm},L}>0.
\]
Strict positivity of the integral follows from continuity, nonnegativity, and \(F_L(0)^3=L^6>0\). By \cref{obj:def:endpoint-kernel},
\[
K_L(x)=c_{\mathrm{norm},L}^{-1}F_L(x)^3,
\]
and hence
\[
K_L(x)\ge0\quad(0\le x\le1),
\qquad
\int_0^1K_L(x)\,dx
=
c_{\mathrm{norm},L}^{-1}c_{\mathrm{norm},L}
=1.
\]
% lean: CausalSmith.Stat.RdTruesideNoiseFrontier.endpointKernel_integral_one

\item We next obtain a quantitative lower bound for the normalizer. For \(0<x\le1\), define
\[
\varphi_x:=\arcsin\sqrt{x}.
\]
Then
\[
0\le\varphi_x\le\frac{\pi}{2},
\qquad
\sin\varphi_x=\sqrt{x}.
\]
The polynomial identity above, the Chebyshev identity in \cref{obj:synth_4}, and the double-angle formula give
\[
2\sin^2\varphi_x\,F_L(\sin^2\varphi_x)
=
1-\cos(2L\varphi_x)
=
2\sin^2(L\varphi_x).
\]
In particular,
\[
xF_L(x)=\sin^2(L\varphi_x).
\]

Define the endpoint-window width by
\[
\delta_{\mathrm{win},L}:=\frac{1}{16L^2}.
\]
For \(0<x\le\delta_{\mathrm{win},L}\), we have
\[
L\sqrt{x}\le\frac14.
\]
Using \(\sin v\ge(2/\pi)v\) for \(0\le v\le\pi/2\) at \(v=\varphi_x\), we obtain
\[
\frac{2}{\pi}L\varphi_x
\le L\sin\varphi_x
=L\sqrt{x}
\le\frac14.
\]
Thus \(0\le L\varphi_x\le\pi/2\), so the same sine inequality applies at \(L\varphi_x\). Together with \(\sin\varphi_x\le\varphi_x\), it yields
\[
\frac{2}{\pi}L\sin\varphi_x
\le\frac{2}{\pi}L\varphi_x
\le\sin(L\varphi_x).
\]
Squaring and using \(xF_L(x)=\sin^2(L\varphi_x)\), then dividing by \(x>0\), gives
\[
F_L(x)\ge\left(\frac{2}{\pi}\right)^2L^2.
\]
At \(x=0\), this inequality follows from \(F_L(0)=L^2\) and \(2/\pi\le1\). Therefore it holds throughout \([0,\delta_{\mathrm{win},L}]\). Since \(F_L^3\) is nonnegative on \([0,1]\),
\[
\begin{aligned}
c_{\mathrm{norm},L}
&\ge
\int_0^{\delta_{\mathrm{win},L}}F_L(t)^3\,dt\\
&\ge
\int_0^{\delta_{\mathrm{win},L}}
\left(\left(\frac{2}{\pi}\right)^2L^2\right)^3\,dt\\
&=
\frac{1}{16L^2}\frac{64L^6}{\pi^6}
=
\frac{4}{\pi^6}L^4.
\end{aligned}
\]
% lean: CausalSmith.Stat.RdTruesideNoiseFrontier.endpointNormalizer_lower_bound

\item To obtain the uniform kernel bound, we use
\[
|\sin(L\varphi_x)|\le L|\sin\varphi_x|.
\]
This follows by induction on the nonnegative integer \(L\). The zero case is immediate, and the sine addition formula gives the induction step:
\[
\begin{aligned}
|\sin((L+1)\varphi_x)|
&\le
|\sin(L\varphi_x)|\,|\cos\varphi_x|
+
|\cos(L\varphi_x)|\,|\sin\varphi_x|\\
&\le
|\sin(L\varphi_x)|+|\sin\varphi_x|\\
&\le
(L+1)|\sin\varphi_x|.
\end{aligned}
\]
For \(0<x\le1\), squaring this bound gives
\[
xF_L(x)=\sin^2(L\varphi_x)
\le L^2\sin^2\varphi_x
=L^2x,
\]
so \(F_L(x)\le L^2\). At zero, equality holds. Hence
\[
F_L(x)^3\le L^6
\le c_{\mathrm{norm},L}\,c_{\mathrm{env}}L^2
\qquad(0\le x\le1),
\]
where the second inequality is the normalizer lower bound multiplied by \(c_{\mathrm{env}}L^2\). Dividing by the positive normalizer proves
\[
K_L(x)\le c_{\mathrm{env}}L^2
\qquad(0\le x\le1).
\]
% lean: CausalSmith.Stat.RdTruesideNoiseFrontier.endpointKernel_le_index_sq

\item Multiplying the uniform bound by the nonnegative kernel and integrating gives
\[
\begin{aligned}
\int_0^1K_L(x)^2\,dx
&\le
\int_0^1 c_{\mathrm{env}}L^2K_L(x)\,dx\\
&=
c_{\mathrm{env}}L^2\int_0^1K_L(x)\,dx
=
c_{\mathrm{env}}L^2.
\end{aligned}
\]
% lean: CausalSmith.Stat.RdTruesideNoiseFrontier.endpointKernel_square_integral_bound

\item For the moment estimate, we also need a tail bound. For \(0<x\le1\), the Chebyshev identity gives
\[
T_L(1-2x)
=
\cos\!\bigl(L\arccos(1-2x)\bigr)
\ge-1.
\]
Consequently,
\[
2xF_L(x)=1-T_L(1-2x)\le2,
\qquad
F_L(x)\le x^{-1}.
\]
Cubing this inequality and using the normalizer lower bound yields
\[
F_L(x)^3
\le x^{-3}
\le c_{\mathrm{norm},L}\,c_{\mathrm{env}}L^{-4}x^{-3}.
\]
Division by \(c_{\mathrm{norm},L}>0\) proves
\[
K_L(x)\le c_{\mathrm{env}}L^{-4}x^{-3}
\qquad(0<x\le1).
\]
% lean: CausalSmith.Stat.RdTruesideNoiseFrontier.endpointKernel_tail_bound

\item Define the splitting point by
\[
\delta_L:=L^{-2},
\qquad
0<\delta_L\le1.
\]
On \([0,\delta_L]\), monotonicity of \(x^\beta\), nonnegativity of \(K_L\), and the uniform kernel bound give
\[
x^\beta K_L(x)
\le
\delta_L^\beta c_{\mathrm{env}}L^2.
\]
Therefore,
\[
\begin{aligned}
\int_0^{\delta_L}x^\beta K_L(x)\,dx
&\le
\int_0^{\delta_L}
\delta_L^\beta c_{\mathrm{env}}L^2\,dx\\
&=
c_{\mathrm{env}}\delta_L^\beta,
\end{aligned}
\]
because \(\delta_LL^2=1\).
% lean: CausalSmith.Stat.RdTruesideNoiseFrontier.endpointKernel_moment_bound

\item On \([\delta_L,1]\), the tail bound gives
\[
x^\beta K_L(x)
\le
c_{\mathrm{env}}\delta_L^2x^{\beta-3}.
\]
Since \(\delta_L>0\) and \(0<\beta\le1\), integration of this power yields
\[
\begin{aligned}
\int_{\delta_L}^1x^\beta K_L(x)\,dx
&\le
c_{\mathrm{env}}\delta_L^2
\int_{\delta_L}^1x^{\beta-3}\,dx\\
&=
c_{\mathrm{env}}\delta_L^2
\frac{\delta_L^{\beta-2}-1}{2-\beta}.
\end{aligned}
\]
Here \(2-\beta\ge1\) and \(\delta_L^{\beta-2}\ge0\), so
\[
\delta_L^{\beta-2}-1
\le
(2-\beta)\delta_L^{\beta-2}.
\]
It follows that
\[
\int_{\delta_L}^1x^\beta K_L(x)\,dx
\le
c_{\mathrm{env}}\delta_L^2\delta_L^{\beta-2}
=
c_{\mathrm{env}}\delta_L^\beta.
\]
These identities also cover \(L=1\): then \(\delta_L=1\), and the integral over \([\delta_L,1]\) is zero.
% lean: CausalSmith.Stat.RdTruesideNoiseFrontier.endpointKernel_moment_bound

\item The function \(x^\beta K_L(x)\) is continuous on \([0,1]\), so its integral splits at \(\delta_L\). Combining the two bounds gives
\[
\begin{aligned}
M_{L,\beta}
&=
\int_0^{\delta_L}x^\beta K_L(x)\,dx
+
\int_{\delta_L}^1x^\beta K_L(x)\,dx\\
&\le
2c_{\mathrm{env}}\delta_L^\beta
=
2c_{\mathrm{env}}L^{-2\beta}.
\end{aligned}
\]
Thus we may take
\[
C_{\mathrm{bias}}:=2c_{\mathrm{env}}>0,
\qquad
C:=\max\{C_{\mathrm{bias}},c_{\mathrm{env}}\}>0.
\]
The moment and square-integral estimates then become
\[
M_{L,\beta}\le CL^{-2\beta},
\qquad
\int_0^1K_L(x)^2\,dx\le CL^2.
\]
Together with positivity and unit mass, these establish all the assertions for every integer \(L\ge1\).
% lean: CausalSmith.Stat.RdTruesideNoiseFrontier.positive_endpoint
\end{enumerate}
\end{proof}

Positivity and unit integral make the latent weight a normalized averaging measure on the reflected arm interval. Its Hölder moment decreases at order \(L^{-2\beta}\), giving the endpoint construction spatial resolution of order \(L^{-2}\). The squared-integral bound records the corresponding sampling cost before measurement-error correction.

The density and mean restrictions enter distinct parts of this averaging argument. The density lower bound in \cref{obj:ass:density-bounds} supplies the population denominator bound \(B_{d,L}(P)\ge1/4\) in \cref{obj:thm:finite-certificate}. The upper density bound controls the amount of mass multiplying the polynomial weight. Variation of the potential-outcome mean around its cutoff value is controlled by \cref{obj:ass:mean-holder}, with \(M_{L,\beta}\) measuring the resulting endpoint approximation error. Thus the approximation term is governed by potential-mean smoothness, while the density bounds control normalization and weighted mass.

\subsection{Derivative growth and inverse-heat second moments}

Gaussian correction introduces derivatives of the endpoint polynomial. Their growth determines how measurement error changes the second-moment cost. The required inequality applies to a real polynomial \leanref{sym:Polynomial}{\ensuremath{p}} of bounded degree.

\begin{lemmav}{lem:factorial-derivatives}[Factorial polynomial derivative bound]
Let \(m,j\) be nonnegative integers with \(1\le j\le m\), and let \(p\) be a real polynomial of degree at most \(m\). Writing \(p^{(j)}\) for its \(j\)-th ordinary derivative, we have
\[
\left(\int_0^1 \bigl(p^{(j)}(x)\bigr)^2\,dx\right)^{1/2}
\le
\frac{4^j(m+1)^{2j}}{j!}
\left(\int_0^1 p(x)^2\,dx\right)^{1/2}.
\]
\end{lemmav}

\begin{proof}[Proof of \cref{obj:lem:factorial-derivatives}]
\begin{enumerate}
\item We first treat \(j=1\), recording identities that will also be used for higher derivatives. Define, for every \(k\ge0\),
\[
R_k(x)=P_k^{\mathrm{Leg}}(2x-1),\qquad
\omega_k=2k+1,\qquad
c_k=\omega_k\int_0^1 p(x)R_k(x)\,dx,
\qquad
\mathcal E_p=\int_0^1 p(x)^2\,dx.
\]
Changing variables in \cref{obj:lem:classical-legendre-facts} gives
\[
\int_0^1 R_k(x)R_l(x)\,dx
=
\begin{cases}
1/\omega_k,&k=l,\\
0,&k\ne l.
\end{cases}
\]
The polynomials \((-1)^kR_k\), \(0\le k\le m\), have respective degrees \(0,\ldots,m\), so they span the polynomials of degree at most \(m\). The same is therefore true of the \(R_k\). Integrating a basis expansion against \(R_k\) identifies its coefficient as \(c_k\). Consequently,
\[
p(x)=\sum_{k=0}^m c_kR_k(x),
\qquad
\mathcal E_p=\sum_{k=0}^m\frac{c_k^2}{\omega_k}\ge0.
\]
The second identity follows by expanding the square and applying orthogonality.
% lean: CausalSmith.Stat.RdTruesideNoiseFrontier.polynomial_shiftedLegendre_parseval

\item Differentiating the expansion and applying weighted Cauchy--Schwarz yields, for every real \(x\),
\[
p'(x)=\sum_{k=0}^m c_kR_k'(x),
\qquad
p'(x)^2
\le
\left(\sum_{k=0}^m\frac{c_k^2}{\omega_k}\right)
\left(\sum_{k=0}^m\omega_kR_k'(x)^2\right)
=
\mathcal E_p\sum_{k=0}^m\omega_kR_k'(x)^2.
\]
To integrate this bound, first observe that
\[
\sum_{i=0}^{k-1}\omega_i=k^2\qquad(k\ge0);
\]
this follows by induction, since adding \(\omega_k=2k+1\) changes \(k^2\) to \((k+1)^2\). Define the parity-restricted sums
\[
\pi_k=\sum_{\substack{0\le i<k\\k-i\ \mathrm{odd}}}\omega_i
\qquad(k\ge0).
\]
The parity classes alternate when \(k\) increases, and the new index \(i=k\) belongs to the sum for \(k+1\). Thus
\[
\pi_0=0,\qquad
\pi_{k+1}+\pi_k=(k+1)^2,
\qquad
2\pi_k=k(k+1),
\]
where the last identity follows by induction from the preceding recurrence.

For \(k\ge1\), \cref{obj:lem:shifted-legendre-derivative} supplies the derivative expansion. Orthogonality then gives
\[
\int_0^1 R_k'(x)^2\,dx
=
4\sum_{\substack{0\le i<k\\k-i\ \mathrm{odd}}}\omega_i
=
2k(k+1).
\]
For \(k=0\), the same identity holds because \(R_0(x)=1\) and \(R_0'(x)=0\). Integrating the pointwise bound therefore gives
\[
\int_0^1 p'(x)^2\,dx
\le
\mathcal E_p\sum_{k=0}^m\omega_k\,2k(k+1).
\]
Since \(0\le k\le m\) implies \(k(k+1)\le(m+1)^2\),
\[
\sum_{k=0}^m\omega_k\,2k(k+1)
\le
2(m+1)^2\sum_{k=0}^m\omega_k
=
2(m+1)^4
\le
4(m+1)^4.
\]
It follows that
\[
\int_0^1 p'(x)^2\,dx
\le
4(m+1)^4\mathcal E_p
\le
16(m+1)^4\mathcal E_p
=
\left(4(m+1)^2\sqrt{\mathcal E_p}\right)^2.
\]
Taking nonnegative square roots proves the required inequality when \(j=1\).
% lean: CausalSmith.Stat.RdTruesideNoiseFrontier.polynomial_first_derivative_norm_bound

\item Now suppose \(j\ne1\). Since \(j\ge1\), this means \(j\ge2\). Define the nonnegative constant
\[
\Lambda_{m,j}=\frac{4^j(m+1)^{2j}}{j!}.
\]
On the index set \(0,\ldots,m\), define the differentiation matrix
\[
(\mathsf J_m)_{kl}
=
\begin{cases}
2\omega_k,&k<l\ \text{and }l-k\text{ is odd},\\
0,&\text{otherwise}.
\end{cases}
\]
Its powers are defined by
\[
(\mathsf J_m^0)_{kl}
=
\begin{cases}
1,&k=l,\\
0,&k\ne l,
\end{cases}
\qquad
(\mathsf J_m^{r+1})_{kl}
=
\sum_{i=0}^m(\mathsf J_m^r)_{ki}(\mathsf J_m)_{il}
\quad(r\ge0).
\]
Also define the derivative coefficients, for \(r\ge0\) and \(0\le k\le m\), by
\[
c_k^{[r]}=\sum_{l=0}^m(\mathsf J_m^r)_{kl}c_l.
\]
The derivative expansion in \cref{obj:lem:shifted-legendre-derivative}, together with \(R_0'=0\), gives
\[
R_l'(x)=\sum_{k=0}^m(\mathsf J_m)_{kl}R_k(x)
\qquad(0\le l\le m).
\]
Differentiating finite sums applies this coefficient matrix once; induction on \(r\) therefore gives
\[
p^{(r)}(x)=\sum_{k=0}^m c_k^{[r]}R_k(x).
\]
Orthogonality now identifies both square integrals:
\[
\int_0^1 p(x)^2\,dx
=
\sum_{k=0}^m\frac{c_k^2}{\omega_k},
\qquad
\int_0^1 p^{(j)}(x)^2\,dx
=
\sum_{k=0}^m\frac{(c_k^{[j]})^2}{\omega_k}.
\]
% lean: CausalSmith.Stat.RdTruesideNoiseFrontier.polynomial_iterated_derivative_parseval

\item We next bound the entries of the matrix powers. The prefix estimate needed below is
\[
(r+1)\sum_{i=0}^{k-1}\omega_i(i^2)^r
\le
(k^2)^{r+1}
\qquad(k,r\ge0).
\]
To establish it, use the elementary inequality
\[
\xi^{r+1}+(r+1)\zeta\xi^r
\le
(\xi+\zeta)^{r+1}
\qquad(\xi,\zeta\ge0,\ r\ge0).
\]
For \(r=0\) this is equality. Its inductive step follows from
\[
\xi^{r+2}+(r+2)\zeta\xi^{r+1}
\le
\bigl(\xi^{r+1}+(r+1)\zeta\xi^r\bigr)(\xi+\zeta)
\le
(\xi+\zeta)^{r+2},
\]
because the difference in the first inequality is
\((r+1)\zeta^2\xi^r\ge0\).
For fixed \(r\), the prefix estimate follows by induction on \(k\). The empty sum at \(k=0\) is zero, and the inductive step is
\[
\begin{aligned}
(r+1)\sum_{i=0}^{k}\omega_i(i^2)^r
&\le
(k^2)^{r+1}+(r+1)(2k+1)(k^2)^r\\
&\le
\bigl(k^2+2k+1\bigr)^{r+1}
=
\bigl((k+1)^2\bigr)^{r+1}.
\end{aligned}
\]
Here, as usual, the zeroth power is \(1\), including at zero.
% lean: CausalSmith.Stat.RdTruesideNoiseFrontier.shiftedLegendre_prefix_power_bound

\item Every entry of every matrix power is nonnegative, by its defining recurrence. We claim that
\[
0\le(\mathsf J_m^r)_{kl}
\le
\frac{2^r\omega_k(l^2)^{r-1}}{(r-1)!}
\qquad(r\ge1,\ 0\le k,l\le m).
\]
For \(r=1\), this follows directly from
\(0\le(\mathsf J_m)_{kl}\le2\omega_k\).
Suppose the bound holds for \(r\). Entries \((\mathsf J_m)_{il}\) vanish unless \(i<l\), and are bounded by \(2\omega_i\). The recurrence and the prefix estimate therefore give
\[
\begin{aligned}
(\mathsf J_m^{r+1})_{kl}
&\le
\sum_{i=0}^{l-1}
\frac{2^r\omega_k(i^2)^{r-1}}{(r-1)!}\,2\omega_i\\
&=
\frac{2^{r+1}\omega_k}{(r-1)!}
\sum_{i=0}^{l-1}\omega_i(i^2)^{r-1}\\
&\le
\frac{2^{r+1}\omega_k(l^2)^r}{r(r-1)!}
=
\frac{2^{r+1}\omega_k(l^2)^r}{r!}.
\end{aligned}
\]
This also covers \(l=0\), for which the sum is empty. Define
\[
\Gamma_{m,j}
=
\frac{2^j(m+1)^{2(j-1)}}{(j-1)!}\ge0.
\]
Since \(0\le l\le m\), the entry estimate implies
\[
0\le(\mathsf J_m^j)_{kl}\le\Gamma_{m,j}\omega_k,
\qquad
\bigl((\mathsf J_m^j)_{kl}\bigr)^2
\le
\Gamma_{m,j}^2\omega_k^2.
\]
% lean: CausalSmith.Stat.RdTruesideNoiseFrontier.legendreDerivativePower_entry_bound

\item Weighted Cauchy--Schwarz transfers this entry bound to the coefficients. For every \(0\le k\le m\),
\[
\begin{aligned}
(c_k^{[j]})^2
&\le
\left(\sum_{l=0}^m
\bigl((\mathsf J_m^j)_{kl}\bigr)^2\omega_l\right)
\left(\sum_{l=0}^m\frac{c_l^2}{\omega_l}\right)\\
&\le
\Gamma_{m,j}^2\omega_k^2
\left(\sum_{l=0}^m\omega_l\right)
\left(\sum_{l=0}^m\frac{c_l^2}{\omega_l}\right).
\end{aligned}
\]
Dividing by the positive weight \(\omega_k\), summing over \(k\), and using the total-weight identity gives
\[
\sum_{k=0}^m\frac{(c_k^{[j]})^2}{\omega_k}
\le
\Gamma_{m,j}^2
\left(\sum_{k=0}^m\omega_k\right)^2
\left(\sum_{l=0}^m\frac{c_l^2}{\omega_l}\right)
=
\Gamma_{m,j}^2(m+1)^4
\sum_{l=0}^m\frac{c_l^2}{\omega_l}.
\]
% lean: CausalSmith.Stat.RdTruesideNoiseFrontier.weighted_derivative_matrix_action_bound

\item For every integer \(r\ge0\), induction gives \(r\le2^r\): the base case is immediate, and
\[
r+1\le2^r+1\le2^{r+1}.
\]
Using \(j!=j(j-1)!\) and \(4^j=2^j2^j\), we obtain
\[
\frac{2^j}{(j-1)!}
=
\frac{j\,2^j}{j!}
\le
\frac{4^j}{j!}.
\]
Consequently,
\[
\Gamma_{m,j}(m+1)^2
=
\frac{2^j}{(j-1)!}(m+1)^{2j}
\le
\Lambda_{m,j}.
\]
Both sides are nonnegative, so squaring and substituting into the coefficient estimate yields
\[
\sum_{k=0}^m\frac{(c_k^{[j]})^2}{\omega_k}
\le
\Lambda_{m,j}^2
\sum_{k=0}^m\frac{c_k^2}{\omega_k}.
\]
% lean: CausalSmith.Stat.RdTruesideNoiseFrontier.legendreDerivativeAction_factorial_bound

\item Finally, the two square-integral identities give
\[
\int_0^1 p^{(j)}(x)^2\,dx
\le
\Lambda_{m,j}^2\mathcal E_p
=
\left(\Lambda_{m,j}\sqrt{\mathcal E_p}\right)^2.
\]
Since \(\mathcal E_p\ge0\) and \(\Lambda_{m,j}\ge0\), taking nonnegative square roots proves
\[
\left(\int_0^1 p^{(j)}(x)^2\,dx\right)^{1/2}
\le
\frac{4^j(m+1)^{2j}}{j!}
\left(\int_0^1 p(x)^2\,dx\right)^{1/2}.
\]
The argument includes \(\mathcal E_p=0\) and the endpoint order \(j=m\).
% lean: CausalSmith.Stat.RdTruesideNoiseFrontier.factorial_derivatives
\end{enumerate}
\end{proof}

The factorial dependence on \leanref{sym:j}{\ensuremath{j}}, the derivative order, controls the accumulation of high-order terms in the inverse-heat second moment. Together with the squared-integral bound in \cref{obj:lem:positive-endpoint}, it gives a degree-dependent bound for every derivative contributing to the correction.

The inverse-heat polynomial has both an exact conditional expectation and an exact conditional second moment. These identities apply at every real latent-score argument and throughout the declared noise range.

\begin{lemmav}{lem:exact-covariance}[Exact Gaussian moment identities]
There exists an absolute constant \(C>0\) such that, for every integer \(L\ge1\) and every \(\sigma\in[0,1]\), the following holds. Let \(K_L\) and \(Q_{L,\sigma}\) be defined by \cref{obj:def:endpoint-kernel,obj:def:inverse-heat}, and set
\[
m_L=3(L-1),\qquad
V_L(\sigma)=\frac34\sum_{j=0}^{m_L}\frac{\sigma^{2j}}{j!}
\int_0^1 [K_L^{(j)}(x)]^2\,dx.
\]
For a standard Gaussian random variable \(\varepsilon\), at every \(x\in\mathbb R\),
\[
\mathbb E\!\left[Q_{L,\sigma}(x+\sigma\varepsilon)\right]=K_L(x),
\qquad
\mathbb E\!\left[Q_{L,\sigma}(x+\sigma\varepsilon)^2\right]
=\sum_{j=0}^{m_L}\frac{\sigma^{2j}[K_L^{(j)}(x)]^2}{j!}.
\]
Moreover,
\[
V_L(\sigma)\le C L^2
\sum_{j=0}^{m_L}\frac{(C\sigma^2L^4)^j}{(j!)^3}.
\]
\end{lemmav}

\begin{proof}[Proof of \cref{obj:lem:exact-covariance}]
\begin{enumerate}
\item We first establish the variance bound and choose its constant. Applying \cref{obj:lem:positive-endpoint} with \(\beta=1\), choose an absolute constant \(C_{\mathrm{end}}>0\) such that
\[
\int_0^1 K_L(x)^2\,dx\le C_{\mathrm{end}}L^2
\qquad(L\ge1).
\]
Define
\[
C=\max\{C_{\mathrm{end}},1296\}.
\]
Then \(C>0\). Fix \(L\ge1\) and \(\sigma\in[0,1]\). We have
\[
m_L+1=3L-2\le3L,
\qquad
\int_0^1 K_L(x)^2\,dx\le CL^2.
\]
The numerator \(1-T_L(1-2x)\) has degree at most \(L\); its polynomial quotient by \(2x\) has degree at most \(L-1\). Cubing and multiplying by the normalization constant in \cref{obj:def:endpoint-kernel} therefore gives
\[
\deg K_L\le3(L-1)=m_L.
\]

For \(1\le j\le m_L\), squaring the nonnegative sides of \cref{obj:lem:factorial-derivatives}, applied to \(K_L\), gives
\[
\int_0^1 [K_L^{(j)}(x)]^2\,dx
\le
\frac{16^j(m_L+1)^{4j}}{(j!)^2}
\int_0^1 K_L(x)^2\,dx.
\]
For \(j=0\), this inequality is an equality. Thus it holds throughout \(0\le j\le m_L\), including when \(L=1\). Moreover,
\[
16^j(m_L+1)^{4j}
=\bigl(16(m_L+1)^4\bigr)^j
\le(1296L^4)^j
\le(CL^4)^j.
\]
Consequently, for each such \(j\),
\[
\begin{aligned}
\frac{\sigma^{2j}}{j!}
\int_0^1 [K_L^{(j)}(x)]^2\,dx
&\le
\frac{\sigma^{2j}}{j!}
\frac{16^j(m_L+1)^{4j}CL^2}{(j!)^2}\\
&\le
CL^2\frac{(C\sigma^2L^4)^j}{(j!)^3}.
\end{aligned}
\]
Every summand on the left is nonnegative. Summing and using \(3/4\le1\) proves
\[
V_L(\sigma)
\le
\sum_{j=0}^{m_L}\frac{\sigma^{2j}}{j!}
\int_0^1 [K_L^{(j)}(x)]^2\,dx
\le
CL^2\sum_{j=0}^{m_L}
\frac{(C\sigma^2L^4)^j}{(j!)^3}.
\]
% lean: CausalSmith.Stat.RdTruesideNoiseFrontier.kernelVariance_factorial_envelope

\item Suppose first that \(\sigma=0\). In \cref{obj:def:inverse-heat}, the term \(k=0\) equals \(K_L\), and every term with \(k>0\) vanishes. Hence
\[
Q_{L,0}(w)=K_L(w).
\]
Since \(m_L\ge0\), the derivative sum contains its zeroth term. With the usual convention \(0^0=1\), its remaining terms vanish:
\[
\sum_{j=0}^{m_L}\frac{0^{2j}[K_L^{(j)}(x)]^2}{j!}
=K_L(x)^2.
\]
The two Gaussian expectations are therefore
\[
\mathbb E[Q_{L,0}(x+0\varepsilon)]=K_L(x),
\qquad
\mathbb E[Q_{L,0}(x+0\varepsilon)^2]=K_L(x)^2,
\]
as required.
% lean: CausalSmith.Stat.RdTruesideNoiseFrontier.inverseHeat_zero_covariance

\item Assume now that \(\sigma\ne0\), and fix \(x\in\mathbb R\). For a nonnegative integer \(m_{\mathrm{cut}}\), a real number \(\sigma_{\mathrm{heat}}\), and a real polynomial \(p\), define the polynomial operator
\[
(\mathcal T_{m_{\mathrm{cut}},\sigma_{\mathrm{heat}}}p)(w)
=
\sum_{k=0}^{\lfloor m_{\mathrm{cut}}/2\rfloor}
\frac{(-\sigma_{\mathrm{heat}}^2/2)^k}{k!}\,
p^{(2k)}(w).
\]
Thus \cref{obj:def:inverse-heat} gives
\[
Q_{L,\sigma}=\mathcal T_{m_L,\sigma}K_L.
\]
Define the affine polynomial and its monomial coefficients by
\[
p_{\mathrm{aff}}(z)=K_L(x+\sigma z)
=\sum_{j=0}^{m_L}c_{\mathrm{aff},j}z^j
\qquad(z\in\mathbb R).
\]
Its degree is at most \(m_L\), by the degree bound established above.

Repeated differentiation of the affine substitution gives, for every nonnegative integer \(j\),
\[
p_{\mathrm{aff}}^{(j)}(z)
=\sigma^jK_L^{(j)}(x+\sigma z).
\]
Substituting this identity into the definition of the operator yields
\[
\begin{aligned}
(\mathcal T_{m_L,1}p_{\mathrm{aff}})(z)
&=
\sum_{k=0}^{\lfloor m_L/2\rfloor}
\frac{(-1/2)^k}{k!}\,
\sigma^{2k}K_L^{(2k)}(x+\sigma z)\\
&=
(\mathcal T_{m_L,\sigma}K_L)(x+\sigma z).
\end{aligned}
\]
Differentiation distributes over finite sums and commutes with scalar multiplication. Applying this to the monomial decomposition of \(p_{\mathrm{aff}}\), we obtain
\[
Q_{L,\sigma}(x+\sigma z)
=
\sum_{j=0}^{m_L}
c_{\mathrm{aff},j}\,
(\mathcal T_{m_L,1}(z^j))(z).
\]
Here the operator acts on the polynomial \(z\mapsto z^j\) before evaluation.
% lean: CausalSmith.Stat.RdTruesideNoiseFrontier.finiteInverseHeat_affine_eval

\item We identify the transformed monomials with the probabilists' Hermite polynomials \(H_j\) of \cref{obj:lem:classical-hermite-facts}. For \(0\le j\le m_L\), all derivatives of \(z^j\) of order greater than \(j\) vanish. Enlarging the cutoff from \(j\) to \(m_L\) therefore adds zero terms, and
\[
\mathcal T_{m_L,1}(z^j)
=
\mathcal T_{j,1}(z^j)
=
\sum_{k=0}^{\lfloor j/2\rfloor}
\frac{(-1/2)^k}{k!}\,
\frac{j!}{(j-2k)!}\,z^{j-2k}.
\]

To compare coefficients, define the odd double factorial by
\[
(2k-1)!!=\prod_{k'=1}^{k}(2k'-1)
\qquad(k\ge0),
\]
where the empty product equals \(1\), so that \((-1)!!=1\).
Separating the even and odd factors of \((2k)!\) gives
\[
(2k)!=2^k k!(2k-1)!!.
\]
Expanding the generating function in \cref{obj:lem:classical-hermite-facts} coefficientwise gives, for nonnegative integers \(i,j\),
\[
\text{coefficient of \(z^i\) in \(H_j(z)\)}
=
\begin{cases}
(-1)^k\binom ji(2k-1)!!,
& j=i+2k,\quad k\ge0,\\
0,
& i>j\ \text{or }j-i\text{ is odd}.
\end{cases}
\]
Indeed, the coefficient of \(z^it^{i+2k}\) in
\(\exp(zt)\exp(-t^2/2)\) is
\((-1/2)^k/(i!k!)\); multiplying by \((i+2k)!\) and using the factorial identity above gives the displayed coefficient.

When \(j=i+2k\), this coefficient equals
\[
(-1)^k\binom ji(2k-1)!!
=
\frac{(-1/2)^k}{k!}\frac{j!}{i!}.
\]
This is exactly the coefficient of \(z^i\) in
\(\mathcal T_{j,1}(z^j)\). For \(i>j\), or for \(j-i\) odd, both coefficients are zero. Thus
\[
\mathcal T_{m_L,1}(z^j)=H_j(z)
\qquad(0\le j\le m_L).
\]
% lean: CausalSmith.Stat.RdTruesideNoiseFrontier.finiteInverseHeat_unit_monomial

Finally, finite Taylor expansion followed by scaling gives
\[
K_L(x+\sigma z)
=
\sum_{j=0}^{m_L}
\frac{\sigma^jK_L^{(j)}(x)}{j!}\,z^j,
\qquad
c_{\mathrm{aff},j}
=
\frac{\sigma^jK_L^{(j)}(x)}{j!}.
\]
Substituting these coefficients and the transformed monomials into the preceding decomposition proves the finite identity
\[
Q_{L,\sigma}(x+\sigma z)
=
\sum_{j=0}^{m_L}
\frac{\sigma^jK_L^{(j)}(x)}{j!}\,H_j(z)
\qquad(z\in\mathbb R).
\]
% lean: CausalSmith.Stat.RdTruesideNoiseFrontier.finiteInverseHeat_hermite_expansion

\item By \cref{obj:lem:classical-hermite-facts}, \(H_0=1\), and orthogonality with this constant polynomial gives
\[
\mathbb E[H_j(\varepsilon)]
=
\mathbb E[H_j(\varepsilon)H_0(\varepsilon)]
=
\begin{cases}
1,&j=0,\\
0,&j\ge1.
\end{cases}
\]
The Hermite polynomials are square integrable under the Gaussian law, so expectations pass through the finite sum. Its zeroth term is present because \(m_L+1>0\). Hence
\[
\mathbb E[Q_{L,\sigma}(x+\sigma\varepsilon)]
=
c_{\mathrm{aff},0}
=
K_L(x).
\]
% lean: CausalSmith.Stat.RdTruesideNoiseFrontier.hermite_finite_sum_mean

\item Expanding the square of the same finite sum and applying \cref{obj:lem:classical-hermite-facts} gives
\[
\begin{aligned}
\mathbb E[Q_{L,\sigma}(x+\sigma\varepsilon)^2]
&=
\sum_{j=0}^{m_L}\sum_{k=0}^{m_L}
c_{\mathrm{aff},j}c_{\mathrm{aff},k}
\mathbb E[H_j(\varepsilon)H_k(\varepsilon)]\\
&=
\sum_{j=0}^{m_L}c_{\mathrm{aff},j}^{\,2}j!\\
&=
\sum_{j=0}^{m_L}
\frac{\sigma^{2j}[K_L^{(j)}(x)]^2}{j!}.
\end{aligned}
\]
The products are integrable by square integrability and the Cauchy--Schwarz inequality, which justifies the finite expansion under expectation. The last equality uses \(j!>0\) for every \(j\ge0\). Together with the zero-noise case and the variance bound, this proves the result with the chosen absolute constant \(C\).
% lean: CausalSmith.Stat.RdTruesideNoiseFrontier.hermite_finite_sum_square
\end{enumerate}
\end{proof}

The first identity identifies the latent endpoint average targeted by the observable correction. The second accounts for the full Gaussian distribution of the proxy, including its tails at positive noise scales. The cubic factorial in the final bound combines the factorial derivative control with the Gaussian second-moment weights.

Under \cref{obj:ass:gaussian,obj:ass:error-independence}, these identities apply conditionally on the latent score and potential outcomes, using the reflection convention in \cref{obj:synth_5}. The density upper bound makes \(V_L(\sigma)\) a common second-moment bound for the arm summands. Binary outcomes also place the outcome-marked summands under that bound. Accordingly, the sampling term in \cref{obj:thm:finite-certificate} accounts for both empirical averages over the entire observed experiment.

\subsection{The finite-sample certificate}

The auxiliary results supply the endpoint approximation and sampling controls used by \cref{obj:thm:finite-certificate}. The construction in \cref{obj:def:ratio} addresses the remaining issue of turning empirical marked and unmarked averages into arm estimates. Its population denominator is bounded below by \(1/4\), while the empirical denominator is floored at \(1/8\). The floor therefore lies below the guaranteed population mass and keeps the ratio defined for every observed sample. Projection onto \([1/4,3/4]\) aligns each arm estimate with the maintained mean range in \cref{obj:ass:mean-bounds}.

For every fixed positive index, \cref{obj:thm:finite-certificate} bounds expected absolute effect-estimation error by
\[
12M_{L,\beta}+28\sqrt{\frac{V_L(\sigma)}{n}}.
\]
The public radius in \cref{obj:synth_6} uses the same approximation term and the larger sampling term \(28\sqrt{40V_L(\sigma)/n}\), delivering coverage of at least \(9/10\). The certificate thus translates positive latent averaging and exact observable second moments into finite-sample estimation and inference guarantees under the benchmark restrictions.

Interval geometry and measurability are also part of the declared decision problem. The projected arm estimates place the estimated contrast in \([-1/2,1/2]\). Hence the intersection defining \(I_L\) in \cref{obj:synth_6} is a nonempty connected closed interval. The target belongs to the same contrast range under \cref{obj:ass:mean-bounds}, so intersection with \([-1,1]\) preserves its coverage. Finite polynomial averages, the denominator floor, projection, and the interval endpoints are measurable functions of the observed sample. Their outputs are constant in the independent procedure seed, placing them within the randomized experiment of \cref{obj:def:honest-class}.

Finally, \cref{obj:def:public-selection} compares radii over a finite public candidate set and resolves ties by choosing the smallest index. For fixed public inputs, the selected index is fixed across samples; its coverage is therefore the coverage of its selected candidate. The index-zero values in \cref{obj:synth_6} provide an estimate of zero and the interval \([-1/2,1/2]\), which contains every admissible target. This fallback supplies coverage one, absolute error at most \(1/2\), and interval length one.

Including index zero ensures \(E_\beta(n,\sigma)\le1/2\). For a selected positive index, the candidate risk bound is bounded by its public radius, and interval length is bounded by twice that radius. These guarantees, together with the fallback, are the selected risk and expected-length bounds in \cref{obj:thm:finite-certificate}. The selected interval consequently belongs to the uniformly honest class in \cref{obj:def:honest-class}, with precision controlled by a deterministic certificate over the full benchmark.


\section{Lower bounds for the observed experiment}

The \leanref{S-4}{lower-bound witnesses} combine full-interval regularity with separation of the cutoff treatment contrast. Their statistical comparison concerns the complete observed triple: the noisy score, recorded assignment, and binary outcome. The construction proceeds from polynomial cancellation to observed-law proximity and then to lower bounds for estimation risk and honest expected interval length.

\subsection{Polynomial cancellation and legal alternatives}

We first fix the polynomial normalization used throughout the construction. The normalized Legendre polynomials \leanref{sym:Legendre}{\ensuremath{P_j^{\mathrm{Leg}}}} and their shifts to the unit interval are defined as follows.



The normalized Legendre polynomials \leanref{sym:Legendre}{\ensuremath{P_j^{\mathrm{Leg}}}} supply orthogonality, endpoint values, and uniform bounds on the reference interval. The following identities collect these properties.

\begin{theoremv}{lem:classical-legendre-facts}[Classical Legendre polynomial identities]
For the normalized Legendre polynomials \(P_j^{\mathrm{Leg}}\), every pair of nonnegative integers \(j,k\) satisfies
\[
\int_{-1}^{1} P_j^{\mathrm{Leg}}(t)P_k^{\mathrm{Leg}}(t)\,dt
=
\begin{cases}
\dfrac{2}{2j+1}, & j=k,\\
0, & j\ne k.
\end{cases}
\]
For every nonnegative integer \(j\), reflection symmetry holds for every \(t\in\mathbb R\):
\[
P_j^{\mathrm{Leg}}(-t)=(-1)^jP_j^{\mathrm{Leg}}(t).
\]
The same polynomial satisfies
\[
|P_j^{\mathrm{Leg}}(t)|\le 1
\quad\text{for every }t\in[-1,1],
\qquad
P_j^{\mathrm{Leg}}(-1)=(-1)^j.
\]
\end{theoremv}

Orthogonality controls the integrated squared size of a polynomial block, while the endpoint identities align its components at the cutoff. The uniform bound controls its amplitude throughout the interval. Together, these properties support a perturbation with a prescribed cutoff value and small integrated squared mass.

Full-interval smoothness also requires control of derivatives after shifting the polynomials to the unit interval. The relevant bound retains the quadratic dependence on polynomial degree.

\begin{lemmav}{lem:shifted-legendre-derivative}[Shifted Legendre derivative bounds]
Let \(P_k^{\mathrm{Leg}}\) denote the degree-\(k\) Legendre polynomial normalized by \(P_k^{\mathrm{Leg}}(1)=1\), and define
\[
R_k(t)=P_k^{\mathrm{Leg}}(2t-1),\qquad k\in\mathbb N,\quad t\in\mathbb R.
\]
For every integer \(j\ge1\),
\[
R_j'(t)=2\sum_{\substack{0\le k<j\\j-k\ \mathrm{odd}}}(2k+1)R_k(t)
\qquad\text{for all }t\in\mathbb R,
\]
and
\[
|R_j'(t)|\le j(j+1)
\qquad\text{for all }t\in[0,1].
\]
\end{lemmav}

\begin{proof}[Proof of \cref{obj:lem:shifted-legendre-derivative}]
Fix \(j\ge1\). We first establish the derivative expansion, then bound it on the unit interval.

\begin{enumerate}
\item By the change of variables from \([0,1]\) to \([-1,1]\), the orthogonality in \cref{obj:lem:classical-legendre-facts} gives, for all nonnegative integers \(k,\ell\),
\[
\int_0^1 R_k(t)R_\ell(t)\,dt
=\frac12\int_{-1}^1
P_k^{\mathrm{Leg}}(t)P_\ell^{\mathrm{Leg}}(t)\,dt
=
\begin{cases}
\dfrac{1}{2k+1},&k=\ell,\\
0,&k\ne\ell.
\end{cases}
\]
Each polynomial \((-1)^kR_k\) has degree \(k\) and nonzero leading coefficient. Consequently,
\[
\operatorname{span}_{\mathbb R}\{(-1)^kR_k:0\le k<j\}
=
\{p\in\mathbb R[t]:\deg p<j\},
\]
where the set on the right includes the zero polynomial. Orthogonality and linearity therefore imply
\[
\int_0^1 p(t)R_j(t)\,dt=0
\qquad
\text{for every real polynomial \(p\) of degree below \(j\)}.
\]
Indeed, this integral vanishes on each spanning polynomial \((-1)^kR_k\), and hence on their linear combinations.
% lean: CausalSmith.Stat.RdTruesideNoiseFrontier.polynomial_shiftedLegendre_orthogonal

\item The same spanning argument gives the expansion
\[
p(t)=
\sum_{k=0}^{j-1}(2k+1)
\left(\int_0^1p(x)R_k(x)\,dx\right)R_k(t),
\qquad t\in\mathbb R,
\]
for every real polynomial \(p\) of degree below \(j\), including zero. To verify it on a spanning polynomial, take \(p=(-1)^\ell R_\ell\), where \(0\le\ell<j\). The right-hand side becomes
\[
\sum_{k=0}^{j-1}(2k+1)(-1)^\ell
\left(\int_0^1R_\ell(x)R_k(x)\,dx\right)R_k(t)
=(-1)^\ell R_\ell(t)=p(t).
\]
Both sides are linear in \(p\), so the identity extends to the entire span.
% lean: CausalSmith.Stat.RdTruesideNoiseFrontier.polynomial_shiftedLegendre_expansion

\item Polynomial differentiation gives
\[
\deg R_j'<j\quad(j\ge1),
\qquad R_0'=0.
\]
Applying the preceding expansion to \(R_j'\) yields
\[
R_j'(t)=
\sum_{k=0}^{j-1}(2k+1)
\left(\int_0^1R_j'(x)R_k(x)\,dx\right)R_k(t).
\]
For \(0\le k<j\), the polynomial \(R_k'\) has degree below \(j\), or is zero when \(k=0\). Thus the orthogonality established above supplies
\[
\int_0^1R_k'(x)R_j(x)\,dx=0.
\]
Integration by parts now gives
\[
\int_0^1R_j'(x)R_k(x)\,dx
=
R_k(1)R_j(1)-R_k(0)R_j(0)
-\int_0^1R_k'(x)R_j(x)\,dx.
\]
The endpoint identity and reflection symmetry in
\cref{obj:lem:classical-legendre-facts} imply, for every \(k\ge0\),
\[
R_k(0)=P_k^{\mathrm{Leg}}(-1)=(-1)^k,
\qquad
R_k(1)=P_k^{\mathrm{Leg}}(1)
=(-1)^kP_k^{\mathrm{Leg}}(-1)=1.
\]
Since \(j=(j-k)+k\), the endpoint product satisfies
\[
(-1)^k(-1)^j=(-1)^{j-k}(-1)^{2k}=(-1)^{j-k}.
\]
Consequently,
\[
\int_0^1R_j'(x)R_k(x)\,dx
=1-(-1)^{j-k}
=
\begin{cases}
2,&j-k\text{ odd},\\
0,&j-k\text{ even}.
\end{cases}
\]
Substituting these coefficients into the expansion proves, for every real \(t\),
\[
R_j'(t)=
2\sum_{\substack{0\le k<j\\j-k\ \mathrm{odd}}}(2k+1)R_k(t).
\]
% lean: CausalSmith.Stat.RdTruesideNoiseFrontier.shiftedLegendre_derivative_expansion

\item Fix \(t\in[0,1]\). Then \(2t-1\in[-1,1]\), so the unit bound in \cref{obj:lem:classical-legendre-facts} gives
\[
|R_k(t)|=|P_k^{\mathrm{Leg}}(2t-1)|\le1
\qquad(k\ge0).
\]
Using the expansion, the triangle inequality, and \(2k+1>0\), we obtain
\[
\begin{aligned}
|R_j'(t)|
&\le
2\sum_{\substack{0\le k<j\\j-k\ \mathrm{odd}}}
|(2k+1)R_k(t)|\\
&=
2\sum_{\substack{0\le k<j\\j-k\ \mathrm{odd}}}
(2k+1)|R_k(t)|\\
&\le
2\sum_{\substack{0\le k<j\\j-k\ \mathrm{odd}}}(2k+1).
\end{aligned}
\]
% lean: CausalSmith.Stat.RdTruesideNoiseFrontier.shiftedLegendre_deriv_bound_of_expansion

\item To evaluate the last sum, define, for every integer \(\ell\ge0\),
\[
\omega_\ell:=
\sum_{\substack{0\le k<\ell\\\ell-k\ \mathrm{odd}}}(2k+1).
\]
First,
\[
\sum_{k=0}^{\ell-1}(2k+1)=\ell^2
\qquad(\ell\ge0).
\]
This follows by induction: the empty sum at \(\ell=0\) is zero, and the induction step is
\[
\sum_{k=0}^{\ell}(2k+1)
=\ell^2+(2\ell+1)=(\ell+1)^2.
\]
For each \(k<\ell\), the integers \(\ell-k\) and \(\ell+1-k\) have opposite parity. The additional term \(k=\ell\) belongs to \(\omega_{\ell+1}\). Hence
\[
\omega_{\ell+1}+\omega_\ell
=
\sum_{k=0}^{\ell-1}(2k+1)+(2\ell+1)
=(\ell+1)^2.
\]
We now prove \(2\omega_\ell=\ell(\ell+1)\) by induction. The base case follows from \(\omega_0=0\); if the identity holds at \(\ell\), then
\[
2\omega_{\ell+1}
=2(\ell+1)^2-2\omega_\ell
=2(\ell+1)^2-\ell(\ell+1)
=(\ell+1)(\ell+2).
\]
In particular, \(2\omega_j=j(j+1)\). Substitution into the preceding bound proves
\[
|R_j'(t)|\le j(j+1)
\qquad(t\in[0,1]).
\]
% lean: CausalSmith.Stat.RdTruesideNoiseFrontier.shiftedLegendre_parity_weight_sum
\end{enumerate}
\end{proof}

The derivative estimate in \cref{obj:lem:shifted-legendre-derivative} determines the spatial scale of the witness. We now define the polynomial block and its normalization together with its extension to the latent-score interval. Let \leanref{sym:m}{\ensuremath{m}} be the polynomial order, with \(m\ge2\), and let \leanref{sym:b}{\ensuremath{b}} be the support width, with \(0<b\le1\).

\begin{definitionv}{def:legal-extension}[Cancellation block and full-interval extension $v_{b,m}$]
For an integer \(m\ge2\), define the block normalizer, cancellation polynomial, and fixed amplitude by
\[
N_m=(2m+1)^2-m^2,
\qquad
g_m(t)=\frac{1-t}{N_m}
\sum_{j=m}^{2m}(2j+1)(-1)^jP_j^{\mathrm{Leg}}(2t-1),
\qquad
\kappa=\frac1{100}.
\]
For \(0<b\le1\), the full-interval extension is
\[
v_{b,m}(x)
=
\begin{cases}
1, & -1\le x<0,\\
g_m(x/b), & 0\le x\le b,\\
0, & b<x\le1.
\end{cases}
\]
\end{definitionv}

The block normalizer \leanref{sym:Blocknorm}{\ensuremath{N_m}}, cancellation polynomial \leanref{sym:Block}{\ensuremath{g_m}}, and fixed amplitude constant \leanref{sym:kappa}{\ensuremath{\kappa}} in \cref{obj:def:legal-extension} specify the perturbation on the unit interval. The extension \leanref{sym:Extension}{\ensuremath{v_{b,m}}} in the same definition specifies the treated potential-outcome mean on both sides of the cutoff. Its constant continuation to the left and zero continuation beyond the support allow full-interval regularity to be assessed under the benchmark conditions in \cref{obj:def:model}. The statistical witnesses embed this profile in uniform-score laws with Bernoulli potential outcomes.

\begin{definitionv}{def:alternatives}[Alternative laws \(P^\pm_{b,m}\)]
The latent laws \(P^\pm_{b,m}\) are specified by:
\begin{itemize}
\item \textbf{(Latent score.)}
\[
X\sim\operatorname{Uniform}[-1,1].
\]
\item \textbf{(Potential outcomes.)} Conditional on \(X=x\), the binary potential outcomes \(Y^0\) and \(Y^1\) are independent, with
\[
Y^0\mid X=x\sim\operatorname{Bernoulli}(1/2),
\qquad
Y^1\mid X=x\sim
\operatorname{Bernoulli}\!\left(
\frac12\pm\kappa(b/m^2)^\beta v_{b,m}(x)
\right).
\]
\item \textbf{(Measurement error.)} The standardized measurement error \(\varepsilon\) is independent of \((X,Y^0,Y^1)\), with
\[
\varepsilon\sim N(0,1).
\]
\end{itemize}
The induced observed law uses the benchmark observation map
\[
W=X+\sigma\varepsilon,\qquad
D=\mathbf1\{X\ge0\},\qquad
Y=(1-D)Y^0+DY^1.
\]
\end{definitionv}

The alternative pair \leanref{sym:Alt}{\ensuremath{P^\pm_{b,m}}} in \cref{obj:def:alternatives} varies the treated conditional mean while retaining the same latent-score distribution, untreated mean, and independent Gaussian error. The next result establishes both membership in the benchmark class and separation of the causal target.

\begin{lemmav}{lem:legal-cancellation}[Legal cancellation and separation]
Let \(\sigma\in\mathbb R\), and suppose that:
\begin{itemize}
\item \textbf{(Smoothness.)} \(0<\beta\le1\).
\item \textbf{(Support width.)} \(0<b\le1\).
\item \textbf{(Polynomial order.)} \(m\in\mathbb N\) and \(m\ge2\).
\end{itemize}
For the normalizer \(N_m\), cancellation polynomial \(g_m\), amplitude \(\kappa\), and extension \(v_{b,m}\) defined in \cref{obj:def:legal-extension},
\[
g_m(0)=1,\qquad g_m(1)=0,\qquad
\sup_{t\in[0,1]}|g_m(t)|\le1,\qquad
\int_0^1g_m(t)^2\,dt\le\frac1{N_m},
\]
and, for every integer \(k\) with \(0\le k\le m-2\),
\[
\int_0^1t^kg_m(t)\,dt=0.
\]
Moreover,
\[
\sup_{t\in[0,1]}|g_m^{(1)}(t)|\le7m^2.
\]
For the extension \(v_{b,m}\) in \cref{obj:def:legal-extension}, its full-interval Hölder seminorm satisfies
\[
[v_{b,m}]_\beta
=
\sup_{\substack{x,t\in[-1,1]\\x\ne t}}
\frac{|v_{b,m}(x)-v_{b,m}(t)|}{|x-t|^\beta}
\le7\left(\frac{b}{m^2}\right)^{-\beta}.
\]
The two latent alternatives \(P^\pm_{b,m}\) of \cref{obj:def:alternatives} both belong to \(\mathcal M_\beta(\sigma)\). Their cutoff causal contrasts, where \(\theta(P)\) is the treated conditional potential-outcome mean at zero minus the untreated conditional potential-outcome mean at zero, satisfy
\[
\left|\theta(P^+_{b,m})-\theta(P^-_{b,m})\right|
=
2\kappa\left(\frac{b}{m^2}\right)^\beta.
\]
\end{lemmav}

\begin{proof}[Proof of \cref{obj:lem:legal-cancellation}]
\begin{enumerate}
\item Define the shifted polynomials and the untapered normalized block by
\[
R_j(t)=P_j^{\mathrm{Leg}}(2t-1),
\qquad
q_m^{\mathrm{pol}}(t)
=\frac1{N_m}\sum_{j=m}^{2m}(2j+1)(-1)^jR_j(t),
\qquad j\in\mathbb N,\quad t\in\mathbb R.
\]
Thus \(g_m(t)=(1-t)q_m^{\mathrm{pol}}(t)\). Summing the odd weights gives
\[
\sum_{j=m}^{2m}(2j+1)
=\sum_{j=0}^{2m}(2j+1)-\sum_{j=0}^{m-1}(2j+1)
=(2m+1)^2-m^2
=N_m=3m^2+4m+1>0.
\]
The endpoint identity in \cref{obj:lem:classical-legendre-facts} therefore yields
\[
q_m^{\mathrm{pol}}(0)
=\frac1{N_m}\sum_{j=m}^{2m}(2j+1)(-1)^j(-1)^j
=1.
\]
Consequently,
\[
g_m(0)=1,\qquad g_m(1)=0.
\]
% lean: CausalSmith.Stat.RdTruesideNoiseFrontier.block_zero
% lean: CausalSmith.Stat.RdTruesideNoiseFrontier.block_one

\item For \(t\in[0,1]\), the argument \(2t-1\) lies in \([-1,1]\). The unit bound in \cref{obj:lem:classical-legendre-facts} gives
\[
\left|\sum_{j=m}^{2m}(2j+1)(-1)^jR_j(t)\right|
\le\sum_{j=m}^{2m}(2j+1)|R_j(t)|
\le\sum_{j=m}^{2m}(2j+1)=N_m.
\]
Since \(N_m>0\) and \(0\le1-t\le1\), it follows that
\[
|q_m^{\mathrm{pol}}(t)|\le1,
\qquad
|g_m(t)|=(1-t)|q_m^{\mathrm{pol}}(t)|\le1.
\]
% lean: CausalSmith.Stat.RdTruesideNoiseFrontier.block_abs_le_one

\item Applying the substitution \(u=2t-1\) to the orthogonality identity in \cref{obj:lem:classical-legendre-facts} gives, for \(j,k\in\mathbb N\),
\[
\int_0^1R_j(t)R_k(t)\,dt
=
\begin{cases}
(2j+1)^{-1},&j=k,\\
0,&j\ne k.
\end{cases}
\]
Expanding the square of the finite sum, the off-diagonal terms vanish and each diagonal term contributes \(2j+1\). Hence
\[
\begin{aligned}
\int_0^1[q_m^{\mathrm{pol}}(t)]^2\,dt
&=\frac1{N_m^2}
\sum_{j=m}^{2m}\sum_{k=m}^{2m}
(2j+1)(2k+1)(-1)^{j+k}
\int_0^1R_j(t)R_k(t)\,dt\\
&=\frac1{N_m^2}\sum_{j=m}^{2m}(2j+1)
=\frac1{N_m}.
\end{aligned}
\]
All these integrands are polynomials, so finite sums and integrals may be interchanged. Finally,
\[
g_m(t)^2=(1-t)^2[q_m^{\mathrm{pol}}(t)]^2
\le[q_m^{\mathrm{pol}}(t)]^2
\qquad(t\in[0,1]),
\]
and integration proves
\[
\int_0^1g_m(t)^2\,dt\le\frac1{N_m}.
\]
% lean: CausalSmith.Stat.RdTruesideNoiseFrontier.block_sq_integral_le

\item Fix an integer \(k\) with \(0\le k\le m-2\), and define
\[
\varphi_k(t)=t^k(1-t),\qquad t\in\mathbb R.
\]
Its degree is at most \(k+1\). For every \(j\in\{m,\ldots,2m\}\),
\[
\deg\varphi_k\le k+1<m\le j.
\]
The reversed shifted polynomials \((-1)^iR_i\), \(0\le i<j\), have respective degrees \(i\) and nonzero leading coefficients. Successively cancelling the leading term of a polynomial of degree less than \(j\) therefore expresses it as a linear combination of these polynomials; the zero polynomial is represented by zero coefficients. In particular, there exist real coefficients \(c_{k,j,i}\), \(0\le i<j\), such that
\[
\varphi_k(t)=\sum_{i=0}^{j-1}c_{k,j,i}(-1)^iR_i(t)
\qquad(t\in\mathbb R).
\]
The shifted orthogonality identity established above now gives
\[
\int_0^1\varphi_k(t)R_j(t)\,dt
=\sum_{i=0}^{j-1}c_{k,j,i}(-1)^i
\int_0^1R_i(t)R_j(t)\,dt
=0.
\]
Using the defining finite sum for \(g_m\), we conclude that
\[
\int_0^1t^kg_m(t)\,dt
=\frac1{N_m}\sum_{j=m}^{2m}(2j+1)(-1)^j
\int_0^1\varphi_k(t)R_j(t)\,dt
=0.
\]
This includes \(m=2\), for which the permitted moment is \(k=0\).
% lean: CausalSmith.Stat.RdTruesideNoiseFrontier.block_moment_zero

\item Differentiating the finite polynomial sum gives
\[
(q_m^{\mathrm{pol}})'(t)
=\frac1{N_m}\sum_{j=m}^{2m}(2j+1)(-1)^jR_j'(t).
\]
Every index in this sum satisfies \(j\ge m\ge2\), so \cref{obj:lem:shifted-legendre-derivative} supplies
\[
|R_j'(t)|\le j(j+1)\le(2m)(2m+1)
\qquad(t\in[0,1]).
\]
Thus
\[
\begin{aligned}
|(q_m^{\mathrm{pol}})'(t)|
&\le\frac1{N_m}\sum_{j=m}^{2m}(2j+1)|R_j'(t)|\\
&\le\frac1{N_m}\sum_{j=m}^{2m}(2j+1)(2m)(2m+1)
=(2m)(2m+1).
\end{aligned}
\]
The product rule yields
\[
g_m'(t)=-q_m^{\mathrm{pol}}(t)
+(1-t)(q_m^{\mathrm{pol}})'(t).
\]
Using the unit bound and \(0\le1-t\le1\), we obtain
\[
|g_m'(t)|
\le1+(2m)(2m+1)
=1+4m^2+2m
\le7m^2.
\]
The final calibration follows from
\[
7m^2-(1+4m^2+2m)=(3m+1)(m-1)\ge0
\qquad(m\ge2).
\]
% lean: CausalSmith.Stat.RdTruesideNoiseFrontier.block_deriv_bound

\item Define the clamped, rescaled argument by
\[
\mathfrak c_b(x)=\max\{0,\min\{1,x/b\}\},
\qquad x\in\mathbb R.
\]
It belongs to \([0,1]\). For \(x<0\) it equals zero, for \(0\le x\le b\) it equals \(x/b\), and for \(x>b\) it equals one. The endpoint identities and \cref{obj:def:legal-extension} therefore give
\[
v_{b,m}(x)=g_m(\mathfrak c_b(x)),
\qquad
|v_{b,m}(x)|\le1
\qquad(x\in[-1,1]).
\]
Both minimum and maximum with a fixed real number are contractions. More explicitly, for \(x,t\in\mathbb R\),
\[
\begin{aligned}
|\mathfrak c_b(x)-\mathfrak c_b(t)|
&\le|\min\{1,x/b\}-\min\{1,t/b\}|\\
&\le|x/b-t/b|
=\frac{|x-t|}{b}.
\end{aligned}
\]
The mean value bound for the polynomial \(g_m\) on the convex interval \([0,1]\) gives
\[
|g_m(t_1)-g_m(t_2)|
\le7m^2|t_1-t_2|
\qquad(t_1,t_2\in[0,1]).
\]
Combining these two inequalities proves the bound across all pieces of the extension:
\[
|v_{b,m}(x)-v_{b,m}(t)|
\le\frac{7m^2}{b}|x-t|
\qquad(x,t\in[-1,1]).
\]

For the Hölder estimate, define
\[
\Lambda_{b,m}=\frac{m^2}{b}>0,
\qquad
\Lambda_{b,m}^{\beta}=(b/m^2)^{-\beta}.
\]
Fix distinct \(x,t\in[-1,1]\). If \(\Lambda_{b,m}|x-t|\le1\), then \(0<\beta\le1\) implies
\[
\Lambda_{b,m}|x-t|
\le(\Lambda_{b,m}|x-t|)^\beta,
\]
so the Lipschitz estimate gives
\[
|v_{b,m}(x)-v_{b,m}(t)|
\le7\Lambda_{b,m}|x-t|
\le7(\Lambda_{b,m}|x-t|)^\beta.
\]
If \(\Lambda_{b,m}|x-t|>1\), then
\[
1\le(\Lambda_{b,m}|x-t|)^\beta,
\qquad
|v_{b,m}(x)-v_{b,m}(t)|
\le|v_{b,m}(x)|+|v_{b,m}(t)|\le2.
\]
Consequently, in this case as well,
\[
|v_{b,m}(x)-v_{b,m}(t)|
\le2\le7(\Lambda_{b,m}|x-t|)^\beta.
\]
Since the factors are positive,
\[
(\Lambda_{b,m}|x-t|)^\beta
=\Lambda_{b,m}^\beta|x-t|^\beta.
\]
Dividing by \(|x-t|^\beta\) and taking the supremum proves
\[
[v_{b,m}]_\beta\le7(b/m^2)^{-\beta}.
\]
The same argument applies at \(\beta=1\).
% lean: CausalSmith.Stat.RdTruesideNoiseFrontier.legalExtension_holder_bound

\item For the probability profiles in \cref{obj:def:alternatives}, write
\[
p_\pm(x)=\frac12\pm\kappa(b/m^2)^\beta v_{b,m}(x),
\qquad x\in[-1,1].
\]
The clamp representation makes \(v_{b,m}\), and hence \(p_\pm\), continuous and Borel measurable on this interval. Since \(m\ge2\) and \(0<b\le1\),
\[
0<\frac b{m^2}\le1,
\qquad
0<(b/m^2)^\beta\le1.
\]
Together with the unit bound, this implies
\[
|(b/m^2)^\beta v_{b,m}(x)|\le1,
\qquad
\frac12-\frac1{100}\le p_\pm(x)\le\frac12+\frac1{100}.
\]
In particular,
\[
\frac14\le p_\pm(x)\le\frac34,
\]
so both profiles define valid Bernoulli probabilities.

The construction in \cref{obj:def:alternatives} uses an independent Gaussian factor. Its Gaussian marginal and product structure give
\[
\varepsilon\sim N(0,1),
\qquad
\varepsilon\perp(X,Y^0,Y^1),
\]
as required by \cref{obj:ass:gaussian,obj:ass:error-independence}. The uniform score and binary Bernoulli outcomes give the latent support in \cref{obj:def:model}. Moreover,
\[
f(P^\pm_{b,m},x)=\frac12
\qquad(x\in[-1,1]),
\qquad
\int_{-1}^1f(P^\pm_{b,m},x)\,dx=1.
\]
This density is continuous on the latent interval and satisfies \cref{obj:ass:density-bounds}. The conditional Bernoulli means are
\[
\mu_0(P^\pm_{b,m},x)=\frac12,
\qquad
\mu_1(P^\pm_{b,m},x)=p_\pm(x),
\]
which are continuous and satisfy \cref{obj:ass:mean-bounds}.

For distinct \(x,t\in[-1,1]\), the preceding Hölder estimate gives
\[
\begin{aligned}
|p_\pm(x)-p_\pm(t)|
&=\kappa(b/m^2)^\beta
|v_{b,m}(x)-v_{b,m}(t)|\\
&\le7\kappa(b/m^2)^\beta(b/m^2)^{-\beta}|x-t|^\beta\\
&=\frac7{100}|x-t|^\beta
\le2|x-t|^\beta,
\end{aligned}
\]
where
\[
(b/m^2)^\beta(b/m^2)^{-\beta}=1
\]
because \(b/m^2>0\). When \(x=t\), the same inequality follows from the zero difference. The untreated mean is constant, so
\[
[\mu_0(P^\pm_{b,m},\cdot)]_\beta=0,
\qquad
[\mu_1(P^\pm_{b,m},\cdot)]_\beta\le\frac7{100}\le2.
\]
Thus \cref{obj:ass:mean-holder} holds, completing all conditions of \cref{obj:def:model} for each sign and every \(\sigma\in\mathbb R\):
\[
P^\pm_{b,m}\in\mathcal M_\beta(\sigma).
\]
% lean: CausalSmith.Stat.RdTruesideNoiseFrontier.altLaw_model

\item At the cutoff,
\[
v_{b,m}(0)=g_m(0)=1.
\]
Evaluating the conditional means therefore gives
\[
\theta(P^\pm_{b,m})
=\mu_1(P^\pm_{b,m},0)-\mu_0(P^\pm_{b,m},0)
=\pm\kappa(b/m^2)^\beta.
\]
Since \(\kappa(b/m^2)^\beta\ge0\), subtraction and the absolute value yield
\[
\left|\theta(P^+_{b,m})-\theta(P^-_{b,m})\right|
=\left|2\kappa(b/m^2)^\beta\right|
=2\kappa(b/m^2)^\beta.
\]
% lean: CausalSmith.Stat.RdTruesideNoiseFrontier.altLaw_separation
\end{enumerate}
\end{proof}

The endpoint value and moment cancellations in \cref{obj:lem:legal-cancellation} serve distinct purposes. The former preserves a cutoff contrast of the stated size; the latter supplies cancellation under Gaussian convolution. The derivative and extension bounds tie the perturbation amplitude to full-interval Hölder regularity. Thus the witnesses belong to the exact benchmark class in \cref{obj:def:model}, with effective endpoint scale \(b/m^2\).

\subsection{Full observed-law comparison and decision transfer}

The likelihood comparison uses the treated-score convolution and its signed perturbation. We define these quantities together with the probability-distance conventions before stating the observed-law bound.



The reference convolution \leanref{sym:q}{\ensuremath{q_\sigma}} and signed convolution \leanref{sym:u}{\ensuremath{u_{b,m,\sigma}}}, defined in \cref{obj:lem:full-marked-likelihood}, describe the treated outcome cells across all real proxy values. The support-to-noise scalar \leanref{sym:lambda}{\ensuremath{\lambda_{b,\sigma}}} indexes the factorial tail in the comparison. Probability distances are taken between complete laws under the following conventions.



The extended divergence \leanref{sym:Chisq}{\ensuremath{\chi^2}}, defined in \cref{obj:lem:full-marked-likelihood}, provides a likelihood-based measure of observed-law proximity. The next bound incorporates assignment and both outcome marks in that comparison.

\begin{lemmav}{lem:full-marked-likelihood}[Full marked likelihood bound]
Consider the alternative laws \(P^\pm_{b,m}\) of \cref{obj:def:alternatives}, with the cancellation polynomial \(g_m\) of \cref{obj:lem:legal-cancellation}, under the following parameter conditions:
\begin{itemize}
\item \textbf{(Smoothness.)} \(0<\beta\le1\).
\item \textbf{(Support.)} \(0<b\le1\).
\item \textbf{(Order.)} \(m\) is an integer with \(m\ge2\).
\item \textbf{(Noise.)} \(0<\sigma\le1\).
\end{itemize}
Set
\[
\kappa=\frac1{100},
\qquad
N_m=(2m+1)^2-m^2,
\qquad
\lambda_{b,\sigma}=\frac{b^2}{4\sigma^2},
\]
and define the likelihood-comparison quantities by
\[
q_\sigma(w)
=\int_0^1
\frac{\exp\!\bigl(-(w-x)^2/(2\sigma^2)\bigr)}
{\sqrt{2\pi}\,\sigma}\,dx,
\qquad
u_{b,m,\sigma}(w)
=\int_0^b g_m(x/b)
\frac{\exp\!\bigl(-(w-x)^2/(2\sigma^2)\bigr)}
{\sqrt{2\pi}\,\sigma}\,dx.
\]
Here \(\chi^2\) denotes the extended chi-squared divergence: for probability laws, it is the integral of the squared Radon--Nikodym derivative minus one when the first law is absolutely continuous with respect to the second, and \(+\infty\) otherwise.

The full observed laws satisfy
\[
\begin{aligned}
&\chi^2\!\left(
(P^+_{b,m})_{\mathrm{obs}},
(P^-_{b,m})_{\mathrm{obs}}
\right)\\
&\quad\le
16\kappa^2\left(\frac{b}{m^2}\right)^{2\beta}
\int_{\mathbb R}
\frac{u_{b,m,\sigma}(w)^2}{q_\sigma(w)}\,dw\\
&\quad\le
16\kappa^2\left(\frac{b}{m^2}\right)^{2\beta}
\frac{b}{N_m}
\exp\!\left(\frac{b^2}{24\sigma^2}\right)
\sum_{j=m-1}^{\infty}\frac{\lambda_{b,\sigma}^{\,j}}{j!}.
\end{aligned}
\]
Moreover,
\[
(P^+_{b,m})_{\mathrm{obs}}
\ll
(P^-_{b,m})_{\mathrm{obs}}.
\]
The squared likelihood-ratio deviation
\[
\left(
\frac{d(P^+_{b,m})_{\mathrm{obs}}}
{d(P^-_{b,m})_{\mathrm{obs}}}-1
\right)^2
\]
is integrable under \((P^-_{b,m})_{\mathrm{obs}}\), the function
\(w\mapsto u_{b,m,\sigma}(w)^2/q_\sigma(w)\) is Lebesgue integrable, and the displayed factorial series converges. The divergence compares the full observed triples, including both binary outcome cells and all real proxy values.
\end{lemmav}

\begin{proof}[Proof of \cref{obj:lem:full-marked-likelihood}]
\begin{enumerate}
\item Define the perturbation amplitude and Gaussian density by
\[
\eta_{\mathrm{pert}}
=\kappa\left(\frac{b}{m^2}\right)^\beta,
\qquad
\varphi_\sigma(t)
=\frac{\exp(-t^2/(2\sigma^2))}{\sqrt{2\pi}\,\sigma},
\quad t\in\mathbb R.
\]
Since \(m\ge2\), \(0<b\le1\), and \(\beta>0\),
\[
0<\frac{b}{m^2}\le1,
\qquad
0<\left(\frac{b}{m^2}\right)^\beta\le1,
\qquad
0<\eta_{\mathrm{pert}}\le\frac1{100}\le\frac14.
\]
The polynomial bound in \cref{obj:lem:legal-cancellation}, obtained from the identities of \cref{obj:lem:classical-legendre-facts}, and the extension in \cref{obj:def:legal-extension} give
\[
|v_{b,m}(x)|\le1,
\qquad
\frac14\le\frac12\pm\eta_{\mathrm{pert}}v_{b,m}(x)\le\frac34,
\quad x\in[-1,1].
\]
Consequently, the conditional mass of either value of \(Y^1\) under the positive alternative is at most three times its conditional mass under the negative alternative. The score, \(Y^0\), and error distributions in \cref{obj:def:alternatives} are common to the two laws. Integrating the conditional mass comparison and then applying the observation map therefore gives
\[
P^+_{b,m}\le3P^-_{b,m},
\qquad
(P^+_{b,m})_{\mathrm{obs}}
\le3(P^-_{b,m})_{\mathrm{obs}}.
\]
In particular,
\[
(P^+_{b,m})_{\mathrm{obs}}
\ll(P^-_{b,m})_{\mathrm{obs}}.
\]
Define the likelihood ratio, up to almost-everywhere equality, by
\[
\mathsf{LR}(w,d,y)
=
\frac{d(P^+_{b,m})_{\mathrm{obs}}}
     {d(P^-_{b,m})_{\mathrm{obs}}}(w,d,y),
\qquad
(w,d,y)\in\mathbb R\times\{0,1\}^2.
\]
Domination implies
\[
0\le\mathsf{LR}\le3,
\qquad
(\mathsf{LR}-1)^2\le4
\quad
(P^-_{b,m})_{\mathrm{obs}}\text{-almost everywhere}.
\]
Both observed laws are probability laws, so the squared likelihood-ratio deviation is integrable. The extended divergence consequently has the finite integral representation
\[
\chi^2\!\left((P^+_{b,m})_{\mathrm{obs}},
             (P^-_{b,m})_{\mathrm{obs}}\right)
=
\int(\mathsf{LR}-1)^2\,d(P^-_{b,m})_{\mathrm{obs}}.
\]
% lean: CausalSmith.Stat.RdTruesideNoiseFrontier.altLaw_observed_absolutelyContinuous
% lean: CausalSmith.Stat.RdTruesideNoiseFrontier.altLaw_observed_likelihood_square_integrable
% lean: CausalSmith.Stat.RdTruesideNoiseFrontier.chiSqExt_eq_ofReal_chiSqDiv

\item The positivity of \(\varphi_\sigma\), the bound \(b\le1\), and \cref{obj:lem:legal-cancellation} imply, for every \(w\in\mathbb R\),
\[
q_\sigma(w)>0,
\qquad
\begin{aligned}
|u_{b,m,\sigma}(w)|
&\le\int_0^b |g_m(x/b)|\varphi_\sigma(w-x)\,dx\\
&\le\int_0^b\varphi_\sigma(w-x)\,dx
\le q_\sigma(w).
\end{aligned}
\]
Gaussian normalization and integration over the unit latent interval give
\[
\int_{\mathbb R}q_\sigma(w)\,dw
=
\int_0^1\left(\int_{\mathbb R}\varphi_\sigma(w-x)\,dw\right)dx
=1.
\]
Thus
\[
0\le
\frac{u_{b,m,\sigma}(w)^2}{q_\sigma(w)}
\le q_\sigma(w),
\]
which proves the asserted Lebesgue integrability.

To compute the full observed likelihood, let
\[
\mathsf p_\pm(w,d,y)
=
\frac{d(P^\pm_{b,m})_{\mathrm{obs}}}
     {d(\text{Lebesgue measure}\otimes\text{counting measure on }\{0,1\}^2)}
     (w,d,y).
\]
Versions defined for every \(w\in\mathbb R\) and \(d,y\in\{0,1\}\) are
\[
\begin{aligned}
\mathsf p_\pm(w,1,y)
&=\frac12\int_0^1
 \left(\frac12\pm(2y-1)\eta_{\mathrm{pert}}v_{b,m}(x)\right)
 \varphi_\sigma(w-x)\,dx,\\
\mathsf p_\pm(w,0,y)
&=\frac14\int_{-1}^0\varphi_\sigma(w-x)\,dx.
\end{aligned}
\]
These formulas follow by conditioning on the uniform score and using the independent Gaussian error and the conditional Bernoulli probabilities in \cref{obj:def:alternatives}. Splitting the first integral at \(b\) and using \cref{obj:def:legal-extension} gives
\[
\begin{aligned}
\int_0^1v_{b,m}(x)\varphi_\sigma(w-x)\,dx
&=\int_0^bg_m(x/b)\varphi_\sigma(w-x)\,dx
  +\int_b^1 0\,dx\\
&=u_{b,m,\sigma}(w).
\end{aligned}
\]
This identity also covers \(b=1\), when the second integral is zero. Hence
\[
\begin{aligned}
\mathsf p_\pm(w,1,1)
&=\frac{q_\sigma(w)}4
  \pm\frac{\eta_{\mathrm{pert}}u_{b,m,\sigma}(w)}2,\\
\mathsf p_\pm(w,1,0)
&=\frac{q_\sigma(w)}4
  \mp\frac{\eta_{\mathrm{pert}}u_{b,m,\sigma}(w)}2.
\end{aligned}
\]
The untreated densities agree, and all four cell densities are positive. Therefore the likelihood ratio has the version
\[
\mathsf{LR}(w,d,y)
=
\begin{cases}
\mathsf p_+(w,1,y)/\mathsf p_-(w,1,y),&d=1,\\
1,&d=0.
\end{cases}
\]
Integrating with the negative-law density and summing over the binary cells yields
\[
\begin{aligned}
&\chi^2\!\left((P^+_{b,m})_{\mathrm{obs}},
              (P^-_{b,m})_{\mathrm{obs}}\right)\\
&\quad=
\int_{\mathbb R}\sum_{y\in\{0,1\}}
\frac{\bigl(\mathsf p_+(w,1,y)-\mathsf p_-(w,1,y)\bigr)^2}
     {\mathsf p_-(w,1,y)}\,dw.
\end{aligned}
\]
The untreated contributions vanish because their likelihood ratio is one.
% lean: CausalSmith.Stat.RdTruesideNoiseFrontier.markedConvolution_square_div_integrable
% lean: CausalSmith.Stat.RdTruesideNoiseFrontier.altLaw_observed_rnDeriv
% lean: CausalSmith.Stat.RdTruesideNoiseFrontier.altLikelihoodRatio_sq_integral_eq

\item For every \(w\in\mathbb R\),
\[
|\eta_{\mathrm{pert}}u_{b,m,\sigma}(w)|
\le\frac{q_\sigma(w)}4,
\qquad
\frac{q_\sigma(w)}4
\pm\frac{\eta_{\mathrm{pert}}u_{b,m,\sigma}(w)}2
\ge\frac{q_\sigma(w)}8.
\]
The two treated-cell differences are
\[
\begin{aligned}
\mathsf p_+(w,1,1)-\mathsf p_-(w,1,1)
&=\eta_{\mathrm{pert}}u_{b,m,\sigma}(w),\\
\mathsf p_+(w,1,0)-\mathsf p_-(w,1,0)
&=-\eta_{\mathrm{pert}}u_{b,m,\sigma}(w).
\end{aligned}
\]
Thus their squared likelihood contributions satisfy
\[
\begin{aligned}
&\sum_{y\in\{0,1\}}
\frac{\bigl(\mathsf p_+(w,1,y)-\mathsf p_-(w,1,y)\bigr)^2}
     {\mathsf p_-(w,1,y)}\\
&\quad=
\frac{(\eta_{\mathrm{pert}}u_{b,m,\sigma}(w))^2}
     {q_\sigma(w)/4-\eta_{\mathrm{pert}}u_{b,m,\sigma}(w)/2}
+
\frac{(\eta_{\mathrm{pert}}u_{b,m,\sigma}(w))^2}
     {q_\sigma(w)/4+\eta_{\mathrm{pert}}u_{b,m,\sigma}(w)/2}\\
&\quad\le
2\,\frac{(\eta_{\mathrm{pert}}u_{b,m,\sigma}(w))^2}{q_\sigma(w)/8}
=
16\eta_{\mathrm{pert}}^2
\frac{u_{b,m,\sigma}(w)^2}{q_\sigma(w)}.
\end{aligned}
\]
The right-hand side is integrable by the preceding step, and the nonnegative left-hand side is dominated by it. Integrating and using
\[
\eta_{\mathrm{pert}}^2
=\kappa^2\left(\frac{b}{m^2}\right)^{2\beta}
\]
proves the first inequality.
% lean: CausalSmith.Stat.RdTruesideNoiseFrontier.witnessTreatedCellConvolution_integral_bound

\item Define the centered moments for every integer \(j\ge0\) by
\[
\zeta_j
=
\int_0^b g_m(x/b)(x-b/2)^j\,dx.
\]
The change of variable \(x=bt\) gives
\[
\zeta_j
=b\int_0^1g_m(t)(bt-b/2)^j\,dt.
\]
For \(0\le j\le m-2\), the second factor is a polynomial of degree at most \(j\). Expanding it into monomials and applying the cancellation identities of \cref{obj:lem:legal-cancellation} therefore gives
\[
\zeta_j=0,\qquad 0\le j\le m-2.
\]
The same cited result and scaling give
\[
\int_0^b g_m(x/b)^2\,dx
=b\int_0^1g_m(t)^2\,dt
\le\frac b{N_m}.
\]
Also, \(|x-b/2|\le b/2\) on \([0,b]\), so
\[
\int_0^b\bigl((x-b/2)^j\bigr)^2\,dx
\le b(b/2)^{2j}.
\]
Since \(N_m=3m^2+4m+1>0\), Cauchy--Schwarz yields
\[
\begin{aligned}
\zeta_j^2
&\le
\left(\int_0^b g_m(x/b)^2\,dx\right)
\left(\int_0^b\bigl((x-b/2)^j\bigr)^2\,dx\right)\\
&\le\frac{b^2}{N_m}(b/2)^{2j}.
\end{aligned}
\]
Dividing by the positive variance and factorial factors gives
\[
0\le\frac{\zeta_j^2}{\sigma^{2j}j!}
\le
\frac{b^2}{N_m}
\begin{cases}
0,&0\le j<m-1,\\
\lambda_{b,\sigma}^{\,j}/j!,&j\ge m-1,
\end{cases}
\]
because
\[
\frac{(b/2)^{2j}}{\sigma^{2j}}
=\lambda_{b,\sigma}^{\,j}.
\]
The nonnegative exponential series satisfies
\[
\sum_{j=0}^{\infty}\frac{\lambda_{b,\sigma}^{\,j}}{j!}
=e^{\lambda_{b,\sigma}}<\infty.
\]
It follows by comparison that both the retained factorial series and the centered-moment series converge, and
\[
\sum_{j=0}^{\infty}\frac{\zeta_j^2}{\sigma^{2j}j!}
\le
\frac{b^2}{N_m}
\sum_{j=m-1}^{\infty}\frac{\lambda_{b,\sigma}^{\,j}}{j!}.
\]
% lean: CausalSmith.Stat.RdTruesideNoiseFrontier.centered_block_moment_series_bound
% lean: CausalSmith.Stat.RdTruesideNoiseFrontier.likelihoodTail_summable

\item For \(w\in\mathbb R\) and \(x\in[0,b]\), define
\[
\xi_w(x)
=
\frac{-(w-x)^2+(w-b/2)^2+b^2/12}{2\sigma^2}.
\]
Expansion of the numerator and elementary polynomial integration give
\[
\begin{aligned}
\int_0^b\xi_w(x)\,dx
&=\frac1{2\sigma^2}
  \int_0^b\left(2wx-x^2-wb+\frac{b^2}{3}\right)dx\\
&=\frac{wb^2-b^3/3-wb^2+b^3/3}{2\sigma^2}
=0.
\end{aligned}
\]
Since \(e^{\xi_w(x)}\ge1+\xi_w(x)\),
\[
\int_0^b e^{\xi_w(x)}\,dx\ge b.
\]
The Gaussian formula factors as
\[
\varphi_\sigma(w-x)
=
e^{-b^2/(24\sigma^2)}
\varphi_\sigma(w-b/2)e^{\xi_w(x)}.
\]
Integrating this identity and restricting the unit-interval convolution to \([0,b]\) consequently gives, for every real \(w\),
\[
q_\sigma(w)
\ge\int_0^b\varphi_\sigma(w-x)\,dx
\ge
b e^{-b^2/(24\sigma^2)}\varphi_\sigma(w-b/2).
\]
All denominators are positive, so
\[
\frac{u_{b,m,\sigma}(w)^2}{q_\sigma(w)}
\le
\frac{e^{b^2/(24\sigma^2)}}b
\frac{u_{b,m,\sigma}(w)^2}{\varphi_\sigma(w-b/2)}.
\]
% lean: CausalSmith.Stat.RdTruesideNoiseFrontier.marked_square_midpoint_comparison

\item Completing the square gives, for \(x,t,w\in\mathbb R\),
\[
\frac{\varphi_\sigma(w-x)\varphi_\sigma(w-t)}
     {\varphi_\sigma(w-b/2)}
=
\exp\!\left(\frac{(x-b/2)(t-b/2)}{\sigma^2}\right)
\varphi_\sigma\bigl(w-(x+t-b/2)\bigr).
\]
Gaussian normalization therefore yields
\[
\int_{\mathbb R}
\frac{\varphi_\sigma(w-x)\varphi_\sigma(w-t)}
     {\varphi_\sigma(w-b/2)}\,dw
=
\exp\!\left(\frac{(x-b/2)(t-b/2)}{\sigma^2}\right).
\]
On the support square \(x,t\in[0,b]\), \cref{obj:lem:legal-cancellation} gives
\[
|g_m(x/b)g_m(t/b)|\le1,
\qquad
\left|\frac{(x-b/2)(t-b/2)}{\sigma^2}\right|
\le\lambda_{b,\sigma}.
\]
Hence the absolute integral of the signed three-variable kernel is bounded by
\[
\begin{aligned}
&\int_0^b\int_0^b\int_{\mathbb R}
\left|
g_m(x/b)g_m(t/b)
\frac{\varphi_\sigma(w-x)\varphi_\sigma(w-t)}
     {\varphi_\sigma(w-b/2)}
\right|\,dw\,dt\,dx\\
&\quad=
\int_0^b\int_0^b
|g_m(x/b)g_m(t/b)|
\exp\!\left(\frac{(x-b/2)(t-b/2)}{\sigma^2}\right)dt\,dx\\
&\quad\le b^2e^{\lambda_{b,\sigma}}<\infty.
\end{aligned}
\]
Multiplying the two integrals defining \(u_{b,m,\sigma}(w)\) gives
\[
\frac{u_{b,m,\sigma}(w)^2}{\varphi_\sigma(w-b/2)}
=
\int_0^b\int_0^b
g_m(x/b)g_m(t/b)
\frac{\varphi_\sigma(w-x)\varphi_\sigma(w-t)}
     {\varphi_\sigma(w-b/2)}\,dt\,dx.
\]
The absolute-integrability bound justifies Fubini and proves integrability of the left-hand side over \(\mathbb R\). It follows that
\[
\begin{aligned}
&\int_{\mathbb R}
\frac{u_{b,m,\sigma}(w)^2}{\varphi_\sigma(w-b/2)}\,dw\\
&\quad=
\int_0^b\int_0^b
g_m(x/b)g_m(t/b)
\exp\!\left(\frac{(x-b/2)(t-b/2)}{\sigma^2}\right)dt\,dx.
\end{aligned}
\]
In particular, the preceding pointwise midpoint comparison can be integrated:
\[
\int_{\mathbb R}\frac{u_{b,m,\sigma}(w)^2}{q_\sigma(w)}\,dw
\le
\frac{e^{b^2/(24\sigma^2)}}b
\int_{\mathbb R}
\frac{u_{b,m,\sigma}(w)^2}{\varphi_\sigma(w-b/2)}\,dw.
\]
% lean: CausalSmith.Stat.RdTruesideNoiseFrontier.marked_square_reference_integral
% lean: CausalSmith.Stat.RdTruesideNoiseFrontier.marked_square_integral_midpoint_comparison

\item Expand the exponential in the compact double integral. For every \(j\ge0\) and \(x,t\in[0,b]\),
\[
\left|
g_m(x/b)g_m(t/b)
\frac{\bigl((x-b/2)(t-b/2)/\sigma^2\bigr)^j}{j!}
\right|
\le\frac{\lambda_{b,\sigma}^{\,j}}{j!}.
\]
The integrals of these absolute values have the summable majorant
\[
\sum_{j=0}^{\infty}
b^2\frac{\lambda_{b,\sigma}^{\,j}}{j!}
=b^2e^{\lambda_{b,\sigma}}.
\]
Thus termwise integration is justified. Each term factors into the product of two centered moments:
\[
\begin{aligned}
&\int_0^b\int_0^b
g_m(x/b)g_m(t/b)
\frac{\bigl((x-b/2)(t-b/2)/\sigma^2\bigr)^j}{j!}\,dt\,dx\\
&\quad=
\frac1{\sigma^{2j}j!}
\left(\int_0^b g_m(x/b)(x-b/2)^j\,dx\right)^2
=\frac{\zeta_j^2}{\sigma^{2j}j!}.
\end{aligned}
\]
Consequently,
\[
\int_{\mathbb R}
\frac{u_{b,m,\sigma}(w)^2}{\varphi_\sigma(w-b/2)}\,dw
=
\sum_{j=0}^{\infty}\frac{\zeta_j^2}{\sigma^{2j}j!}.
\]
Substituting this identity and the centered-moment bound into the integrated midpoint comparison gives
\[
\begin{aligned}
\int_{\mathbb R}\frac{u_{b,m,\sigma}(w)^2}{q_\sigma(w)}\,dw
&\le
\frac{e^{b^2/(24\sigma^2)}}b
\sum_{j=0}^{\infty}\frac{\zeta_j^2}{\sigma^{2j}j!}\\
&\le
\frac{e^{b^2/(24\sigma^2)}}b\,
\frac{b^2}{N_m}
\sum_{j=m-1}^{\infty}\frac{\lambda_{b,\sigma}^{\,j}}{j!}\\
&=
\frac b{N_m}e^{b^2/(24\sigma^2)}
\sum_{j=m-1}^{\infty}\frac{\lambda_{b,\sigma}^{\,j}}{j!}.
\end{aligned}
\]
Multiplication by the nonnegative factor
\(16\kappa^2(b/m^2)^{2\beta}\) proves the second displayed inequality. Absolute continuity, both integrability assertions, and convergence of the factorial series were established above. The likelihood calculation sums both binary outcome cells and integrates over every real proxy value, so it compares the full observed triples.
% lean: CausalSmith.Stat.RdTruesideNoiseFrontier.marked_square_reference_moment_series
% lean: CausalSmith.Stat.RdTruesideNoiseFrontier.marked_square_factorial_tail_bound
% lean: CausalSmith.Stat.RdTruesideNoiseFrontier.full_marked_likelihood
\end{enumerate}
\end{proof}

The polynomial in \cref{obj:lem:full-marked-likelihood} is the block defined in \cref{obj:def:legal-extension}, with properties established in \cref{obj:lem:legal-cancellation}. The factorial tail records the role of polynomial order in the observed-law comparison. Absolute continuity and integrability make the displayed likelihood comparison valid over the entire proxy support, including extreme noisy scores.

For the decision problem, the relevant proximity condition concerns the independent sample laws. Total variation \leanref{sym:TV}{\ensuremath{\operatorname{TV}}} expresses this condition directly in terms of measurable events. The following transfer gives separate conclusions for absolute estimation loss and honest interval length.

\begin{lemmav}{lem:two-point-transfer}[Two-point minimax risk transfer]
Let \(P^+_{b,m}\) and \(P^-_{b,m}\) be the alternative laws from \cref{obj:def:alternatives}, constructed using \cref{obj:lem:classical-legendre-facts}. Suppose:
\begin{itemize}
\item \textbf{(Parameter ranges.)} \(0<\beta\le1\), \(0<b\le1\), and \(0\le\sigma\le1\).
\item \textbf{(Sample size and degree.)} \(n,m\in\mathbb N\) with \(n\ge2\) and \(m\ge2\).
\item \textbf{(Sample-law proximity.)} The independent observed-sample laws satisfy
\[
\operatorname{TV}\!\left(
(P^+_{b,m})_{\mathrm{obs}}^{\otimes n},
(P^-_{b,m})_{\mathrm{obs}}^{\otimes n}
\right)\le\frac15,
\]
where total variation is the supremum, over measurable events, of the absolute difference between their probabilities.
\end{itemize}
Set \(\kappa=1/100\). Write \(R_\beta(n,\sigma)\) for the benchmark minimax expected absolute error in the cutoff treatment contrast, and \(\mathcal L_\beta(n,\sigma)\) for the benchmark minimax worst expected length of measurable connected closed intervals with uniform \(0.9\) coverage. Both criteria minimize extended nonnegative losses before converting the resulting minimax values to real numbers; their decision classes allow independent procedure randomization. Then
\[
R_\beta(n,\sigma)\ge
\frac45\,\kappa\left(\frac{b}{m^2}\right)^\beta
\qquad\text{and}\qquad
\mathcal L_\beta(n,\sigma)\ge
\frac65\,\kappa\left(\frac{b}{m^2}\right)^\beta.
\]
\end{lemmav}

\begin{proof}[Proof of \cref{obj:lem:two-point-transfer}]
All expectations below are nonnegative extended integrals. We establish the lower bounds for the extended minimax values and then use their finiteness to pass to the real-valued criteria.

\begin{enumerate}
\item
By \cref{obj:lem:legal-cancellation}, the alternatives satisfy
\[
P^+_{b,m},P^-_{b,m}\in\mathcal M_\beta(\sigma),
\qquad
\Delta_{\mathrm{wit}}
:=
\left|\theta(P^+_{b,m})-\theta(P^-_{b,m})\right|
=
2\kappa\left(\frac{b}{m^2}\right)^\beta.
\]
Include the independent procedure seed in the experiment. Define the sample and input spaces, the seed law, and the two sample and experiment laws by
\[
\begin{gathered}
\Omega_{\mathrm{sam},n}
=
\bigl(\mathbb R\times\{0,1\}\times\{0,1\}\bigr)^n,
\qquad
\Omega_n=\Omega_{\mathrm{sam},n}\times[0,1],
\qquad
\lambda_U=\operatorname{Uniform}[0,1],
\\
\mathsf S_\pm
=
\bigl((P^\pm_{b,m})_{\mathrm{obs}}\bigr)^{\otimes n},
\qquad
\mathsf P_\pm
=
\mathsf S_\pm\otimes\lambda_U
=
\mathbb P_{P^\pm_{b,m},n}.
\end{gathered}
\]
These are probability laws on their respective spaces.

For a measurable event \(\mathcal E_{\mathrm{seed}}\subseteq\Omega_n\), define its seed-section probability by
\[
\varphi_{\mathrm{seed}}(\mathbf s)
=
\lambda_U\{u\in[0,1]:(\mathbf s,u)\in\mathcal E_{\mathrm{seed}}\},
\qquad
\mathbf s\in\Omega_{\mathrm{sam},n}.
\]
This function is measurable and takes values in \([0,1]\). Product integration and the layer-cake identity give
\[
\begin{aligned}
\left|
\mathsf P_+(\mathcal E_{\mathrm{seed}})
-
\mathsf P_-(\mathcal E_{\mathrm{seed}})
\right|
&=
\left|
\int\varphi_{\mathrm{seed}}\,d\mathsf S_+
-
\int\varphi_{\mathrm{seed}}\,d\mathsf S_-
\right|
\\
&=
\left|
\int_0^1
\left(
\mathsf S_+\{\varphi_{\mathrm{seed}}\ge t\}
-
\mathsf S_-\{\varphi_{\mathrm{seed}}\ge t\}
\right)\,dt
\right|
\\
&\le
\int_0^1\operatorname{TV}(\mathsf S_+,\mathsf S_-)\,dt
=
\operatorname{TV}(\mathsf S_+,\mathsf S_-).
\end{aligned}
\]
Taking the supremum over measurable events therefore yields
\[
\operatorname{TV}(\mathsf P_+,\mathsf P_-)
\le
\operatorname{TV}(\mathsf S_+,\mathsf S_-)
\le\frac15.
\]
Finally, define the witness amplitude by
\[
\delta_{\mathrm{wit}}
:=
\kappa\left(\frac{b}{m^2}\right)^\beta
\ge0,
\qquad
\Delta_{\mathrm{wit}}=2\delta_{\mathrm{wit}}.
\]
Its nonnegativity follows from \(\kappa=1/100\) and \(b/m^2>0\).
% lean: CausalSmith.Stat.RdTruesideNoiseFrontier.experiment_tv_le

\item
For every measurable estimator \(T:\Omega_n\to\mathbb R\) and every \(P\in\mathcal M_\beta(\sigma)\), define
\[
\mathfrak r_{\mathrm{abs}}(T,P)
:=
\int_{\Omega_n}|T(z)-\theta(P)|\,d\mathbb P_{P,n}(z)
\in[0,\infty].
\]
The extended minimax absolute risk is
\[
R^{\mathrm{ext}}_\beta(n,\sigma)
:=
\inf_{\substack{T:\Omega_n\to\mathbb R\\T\ \mathrm{measurable}}}
\ \sup_{P\in\mathcal M_\beta(\sigma)}
\mathfrak r_{\mathrm{abs}}(T,P).
\]
For a law in the benchmark class, \cref{obj:ass:mean-bounds} gives
\[
\frac14\le\mu_d(P,0)\le\frac34
\quad(d\in\{0,1\}),
\qquad
-\frac12
\le
\theta(P)=\mu_1(P,0)-\mu_0(P,0)
\le\frac12.
\]
Hence the constant estimator
\[
T_{\mathrm{zero}}(z):=0
\qquad(z\in\Omega_n)
\]
satisfies
\[
\mathfrak r_{\mathrm{abs}}(T_{\mathrm{zero}},P)
=
|\theta(P)|
\le\frac12.
\]
Consequently,
\[
0\le R^{\mathrm{ext}}_\beta(n,\sigma)\le\frac12<\infty.
\]
Thus its conversion to a real number preserves its value and order.
% lean: CausalSmith.Stat.RdTruesideNoiseFrontier.risk_value_ne_top

\item
We next obtain the estimator-wise sum bound. Define the dominating measure, its two Radon--Nikodym densities, and their difference by
\[
\xi_{\mathrm{pair}}:=\mathsf P_++\mathsf P_-,
\qquad
p_\pm(z):=\frac{d\mathsf P_\pm}{d\xi_{\mathrm{pair}}}(z),
\qquad
\eta_{\mathrm{pair}}(z):=p_+(z)-p_-(z),
\quad z\in\Omega_n.
\]
Choose finite nonnegative measurable versions of the densities, assigning zero on any exceptional null set. Define the common measure by
\[
\mathsf C_{\mathrm{pair}}(\mathcal E)
:=
\int_{\mathcal E}\min\{p_+(z),p_-(z)\}\,d\xi_{\mathrm{pair}}(z)
\]
for every measurable \(\mathcal E\subseteq\Omega_n\). The pointwise bounds on the minimum imply
\[
\mathsf C_{\mathrm{pair}}\le\mathsf P_+,
\qquad
\mathsf C_{\mathrm{pair}}\le\mathsf P_-.
\]

Considering separately \(p_+\le p_-\) and \(p_-\le p_+\) gives
\[
\min\{p_+,p_-\}
=
\frac{p_++p_--|p_+-p_-|}{2}.
\]
Moreover,
\[
\int p_+\,d\xi_{\mathrm{pair}}
=
\int p_-\,d\xi_{\mathrm{pair}}
=1,
\qquad
\int\eta_{\mathrm{pair}}\,d\xi_{\mathrm{pair}}=0.
\]
To identify the total variation term, define
\[
\mathcal E_{\mathrm{pos}}
:=
\{z\in\Omega_n:\eta_{\mathrm{pair}}(z)\ge0\}.
\]
The zero integral of \(\eta_{\mathrm{pair}}\) and its signs on this event and its complement imply
\[
\begin{aligned}
\int_{\mathcal E_{\mathrm{pos}}}\eta_{\mathrm{pair}}\,d\xi_{\mathrm{pair}}
&=
-\int_{\mathcal E_{\mathrm{pos}}^c}\eta_{\mathrm{pair}}\,d\xi_{\mathrm{pair}}
\\
&=
\frac12\int|\eta_{\mathrm{pair}}|\,d\xi_{\mathrm{pair}}.
\end{aligned}
\]
For any measurable \(\mathcal E\subseteq\Omega_n\), splitting the zero total integral over the event and its complement gives
\[
\int_{\mathcal E}\eta_{\mathrm{pair}}\,d\xi_{\mathrm{pair}}
+
\int_{\mathcal E^c}\eta_{\mathrm{pair}}\,d\xi_{\mathrm{pair}}
=0,
\qquad
\int_{\mathcal E^c}\eta_{\mathrm{pair}}\,d\xi_{\mathrm{pair}}
=
-\int_{\mathcal E}\eta_{\mathrm{pair}}\,d\xi_{\mathrm{pair}}.
\]
The absolute-density integral splits as
\[
\int|\eta_{\mathrm{pair}}|\,d\xi_{\mathrm{pair}}
=
\int_{\mathcal E}|\eta_{\mathrm{pair}}|\,d\xi_{\mathrm{pair}}
+
\int_{\mathcal E^c}|\eta_{\mathrm{pair}}|\,d\xi_{\mathrm{pair}}.
\]
Integrating the pointwise bounds on the signed density and its negative over each piece yields
\[
\begin{aligned}
\int_{\mathcal E}\eta_{\mathrm{pair}}\,d\xi_{\mathrm{pair}}
&\le\int_{\mathcal E}|\eta_{\mathrm{pair}}|\,d\xi_{\mathrm{pair}},
\\
\int_{\mathcal E^c}\eta_{\mathrm{pair}}\,d\xi_{\mathrm{pair}}
&\le\int_{\mathcal E^c}|\eta_{\mathrm{pair}}|\,d\xi_{\mathrm{pair}},
\\
-\int_{\mathcal E}\eta_{\mathrm{pair}}\,d\xi_{\mathrm{pair}}
&\le\int_{\mathcal E}|\eta_{\mathrm{pair}}|\,d\xi_{\mathrm{pair}},
\\
-\int_{\mathcal E^c}\eta_{\mathrm{pair}}\,d\xi_{\mathrm{pair}}
&\le\int_{\mathcal E^c}|\eta_{\mathrm{pair}}|\,d\xi_{\mathrm{pair}}.
\end{aligned}
\]
Adding the third and second inequalities, and then the first and fourth, and using the complement identity gives
\[
\begin{aligned}
-2\int_{\mathcal E}\eta_{\mathrm{pair}}\,d\xi_{\mathrm{pair}}
&\le\int|\eta_{\mathrm{pair}}|\,d\xi_{\mathrm{pair}},
\\
2\int_{\mathcal E}\eta_{\mathrm{pair}}\,d\xi_{\mathrm{pair}}
&\le\int|\eta_{\mathrm{pair}}|\,d\xi_{\mathrm{pair}}.
\end{aligned}
\]
Therefore
\[
\left|\mathsf P_+(\mathcal E)-\mathsf P_-(\mathcal E)\right|
=
\left|\int_{\mathcal E}\eta_{\mathrm{pair}}\,d\xi_{\mathrm{pair}}\right|
\le
\frac12\int|\eta_{\mathrm{pair}}|\,d\xi_{\mathrm{pair}},
\]
with equality for \(\mathcal E=\mathcal E_{\mathrm{pos}}\). It follows that
\[
\operatorname{TV}(\mathsf P_+,\mathsf P_-)
=
\frac12\int|p_+-p_-|\,d\xi_{\mathrm{pair}},
\]
and hence
\[
\begin{aligned}
\mathsf C_{\mathrm{pair}}(\Omega_n)
&=
\int\min\{p_+,p_-\}\,d\xi_{\mathrm{pair}}
\\
&=
\frac{1+1-\int|p_+-p_-|\,d\xi_{\mathrm{pair}}}{2}
=
1-\operatorname{TV}(\mathsf P_+,\mathsf P_-).
\end{aligned}
\]

Now fix any measurable estimator \(T:\Omega_n\to\mathbb R\). The triangle inequality gives, pointwise,
\[
\Delta_{\mathrm{wit}}
\le
|T(z)-\theta(P^+_{b,m})|
+
|T(z)-\theta(P^-_{b,m})|.
\]
Integrating against the common measure and using its domination by each experiment law yields
\[
\begin{aligned}
\Delta_{\mathrm{wit}}
\bigl(1-\operatorname{TV}(\mathsf P_+,\mathsf P_-)\bigr)
&=
\int\Delta_{\mathrm{wit}}\,d\mathsf C_{\mathrm{pair}}
\\
&\le
\int
\left(
|T-\theta(P^+_{b,m})|
+
|T-\theta(P^-_{b,m})|
\right)\,d\mathsf C_{\mathrm{pair}}
\\
&=
\int|T-\theta(P^+_{b,m})|\,d\mathsf C_{\mathrm{pair}}
+
\int|T-\theta(P^-_{b,m})|\,d\mathsf C_{\mathrm{pair}}
\\
&\le
\mathfrak r_{\mathrm{abs}}(T,P^+_{b,m})
+
\mathfrak r_{\mathrm{abs}}(T,P^-_{b,m}).
\end{aligned}
\]
These are inequalities of nonnegative extended integrals, including when an individual risk is infinite.
% lean: CausalSmith.Stat.RdTruesideNoiseFrontier.twoPoint_absRisk_sum

\item
Define the candidate estimation lower bound by
\[
c_{\mathrm{est}}:=\frac45\,\delta_{\mathrm{wit}}\ge0.
\]
The experiment-law proximity and the preceding sum bound imply
\[
2c_{\mathrm{est}}
=
\frac45\,\Delta_{\mathrm{wit}}
\le
\Delta_{\mathrm{wit}}
\bigl(1-\operatorname{TV}(\mathsf P_+,\mathsf P_-)\bigr)
\le
\mathfrak r_{\mathrm{abs}}(T,P^+_{b,m})
+
\mathfrak r_{\mathrm{abs}}(T,P^-_{b,m}).
\]
If the maximum of these two risks were less than \(c_{\mathrm{est}}\), both risks would be finite and less than \(c_{\mathrm{est}}\), giving
\[
\mathfrak r_{\mathrm{abs}}(T,P^+_{b,m})
+
\mathfrak r_{\mathrm{abs}}(T,P^-_{b,m})
<
2c_{\mathrm{est}},
\]
a contradiction. Thus
\[
c_{\mathrm{est}}
\le
\max\left\{
\mathfrak r_{\mathrm{abs}}(T,P^+_{b,m}),
\mathfrak r_{\mathrm{abs}}(T,P^-_{b,m})
\right\}
\le
\sup_{P\in\mathcal M_\beta(\sigma)}
\mathfrak r_{\mathrm{abs}}(T,P).
\]
Taking the infimum over measurable estimators gives
\[
R^{\mathrm{ext}}_\beta(n,\sigma)\ge c_{\mathrm{est}}.
\]
The established finiteness permits conversion to the real-valued criterion, so
\[
R_\beta(n,\sigma)
\ge
\frac45\,\kappa\left(\frac{b}{m^2}\right)^\beta.
\]
% lean: Causalean.Stat.le_minimaxValueENNReal_of_two_point

\item
For an honest interval procedure \(I\in\mathcal I_{\beta,n,\sigma}\), with endpoints \(\ell,\upsilon\) as in \cref{obj:def:honest-class}, and a benchmark law \(P\), define
\[
\mathfrak r_{\mathrm{len}}(I,P)
:=
\int_{\Omega_n}
\max\{\upsilon(z)-\ell(z),0\}\,d\mathbb P_{P,n}(z)
\in[0,\infty].
\]
Define the extended minimax length by
\[
\mathcal L^{\mathrm{ext}}_\beta(n,\sigma)
:=
\inf_{I\in\mathcal I_{\beta,n,\sigma}}
\ \sup_{P\in\mathcal M_\beta(\sigma)}
\mathfrak r_{\mathrm{len}}(I,P).
\]
The constant interval procedure
\[
I_{\mathrm{full}}(z):=[-1,1]
\qquad(z\in\Omega_n)
\]
is measurable. Since \(|\theta(P)|\le1/2\) throughout the benchmark class,
\[
\{z:\theta(P)\in I_{\mathrm{full}}(z)\}=\Omega_n,
\qquad
\mathbb P_{P,n}\{\theta(P)\in I_{\mathrm{full}}\}=1.
\]
It is therefore honest, and
\[
\mathfrak r_{\mathrm{len}}(I_{\mathrm{full}},P)
=
\int_{\Omega_n}2\,d\mathbb P_{P,n}
=2.
\]
Consequently,
\[
0\le\mathcal L^{\mathrm{ext}}_\beta(n,\sigma)\le2<\infty,
\]
so this extended minimax value also admits an order-preserving conversion to a real number.
% lean: CausalSmith.Stat.RdTruesideNoiseFrontier.length_value_ne_top

\item
Fix any \(I\in\mathcal I_{\beta,n,\sigma}\), and define its two coverage events by
\[
\mathcal E_{\mathrm{cov},\pm}
:=
\{z\in\Omega_n:
\ell(z)\le\theta(P^\pm_{b,m})\le\upsilon(z)\}.
\]
These events are measurable. Since both alternatives are legal, honesty in \cref{obj:def:honest-class} supplies
\[
\mathsf P_+(\mathcal E_{\mathrm{cov},+})\ge\frac9{10},
\qquad
\mathsf P_-(\mathcal E_{\mathrm{cov},-})\ge\frac9{10}.
\]
The eventwise total variation bound gives
\[
\left|
\mathsf P_+(\mathcal E_{\mathrm{cov},+})
-
\mathsf P_-(\mathcal E_{\mathrm{cov},+})
\right|
\le
\operatorname{TV}(\mathsf P_+,\mathsf P_-)
\le\frac15,
\]
and therefore
\[
\mathsf P_-(\mathcal E_{\mathrm{cov},+})
\ge\frac9{10}-\frac15=\frac7{10}.
\]
Using the union--intersection identity under the probability law \(\mathsf P_-\), we obtain
\[
\begin{aligned}
\mathsf P_-\bigl(
\mathcal E_{\mathrm{cov},+}\cap\mathcal E_{\mathrm{cov},-}
\bigr)
&=
\mathsf P_-(\mathcal E_{\mathrm{cov},+})
+
\mathsf P_-(\mathcal E_{\mathrm{cov},-})
-
\mathsf P_-\bigl(
\mathcal E_{\mathrm{cov},+}\cup\mathcal E_{\mathrm{cov},-}
\bigr)
\\
&\ge
\frac7{10}+\frac9{10}-1
=
\frac35.
\end{aligned}
\]

On the intersection of the coverage events, both targets lie between \(\ell(z)\) and \(\upsilon(z)\), so
\[
\Delta_{\mathrm{wit}}
=
|\theta(P^+_{b,m})-\theta(P^-_{b,m})|
\le
\upsilon(z)-\ell(z)
\le
\max\{\upsilon(z)-\ell(z),0\}.
\]
Outside the intersection, its indicator is zero and interval length is nonnegative. Thus, on all of \(\Omega_n\),
\[
\Delta_{\mathrm{wit}}\,
\mathbf1_{\mathcal E_{\mathrm{cov},+}\cap\mathcal E_{\mathrm{cov},-}}(z)
\le
\max\{\upsilon(z)-\ell(z),0\}.
\]
Integration under \(\mathsf P_-\) now gives
\[
\begin{aligned}
\mathfrak r_{\mathrm{len}}(I,P^-_{b,m})
&\ge
\Delta_{\mathrm{wit}}\,
\mathsf P_-\bigl(
\mathcal E_{\mathrm{cov},+}\cap\mathcal E_{\mathrm{cov},-}
\bigr)
\\
&\ge
\frac35\,\Delta_{\mathrm{wit}}
=
\frac65\,\delta_{\mathrm{wit}}.
\end{aligned}
\]
% lean: CausalSmith.Stat.RdTruesideNoiseFrontier.twoPoint_interval_lower

\item
For every honest interval procedure, the risk at the minus alternative is bounded by its worst-case risk:
\[
\frac65\,\delta_{\mathrm{wit}}
\le
\mathfrak r_{\mathrm{len}}(I,P^-_{b,m})
\le
\sup_{P\in\mathcal M_\beta(\sigma)}
\mathfrak r_{\mathrm{len}}(I,P).
\]
Taking the infimum over honest procedures yields
\[
\mathcal L^{\mathrm{ext}}_\beta(n,\sigma)
\ge
\frac65\,\delta_{\mathrm{wit}}.
\]
Its finiteness then gives
\[
\mathcal L_\beta(n,\sigma)
\ge
\frac65\,\kappa\left(\frac{b}{m^2}\right)^\beta.
\]
Both arguments use the joint sample-and-seed laws, and hence apply to every allowed randomized procedure.
% lean: Causalean.Stat.le_minimaxValueENNReal
\end{enumerate}
\end{proof}

\Cref{obj:lem:two-point-transfer} connects observed-sample proximity to the target separation in \cref{obj:lem:legal-cancellation}. Its expected-length conclusion applies to measurable connected closed intervals with uniform \(0.9\) coverage, the decision requirement in \cref{obj:def:honest-class}. Both lower bounds cover procedures with arbitrary independent randomization. Consequently, the transfer applies to the full declared decision classes, with the displayed constants distinguishing estimation risk from honest expected length.

\subsection{Direct witnesses and noise-dependent support}

The direct-experiment component of \cref{obj:thm:direct-reduction} uses a localized profile whose support width sets its Hölder scale. Let \leanref{sym:h}{\ensuremath{h}} denote that width, with \(0<h\le1\). The profile and the associated conditional means are specified next.

\begin{definitionv}{synth_7}[Direct-experiment witness profile]
For \(0<h\le1\), define the direct-experiment witness profile on the full latent-score interval by
\[
v_h^{\mathrm{dir}}(x)=
\begin{cases}
1,&-1\le x<0,\\
\max\{1-x/h,0\},&0\le x\le1.
\end{cases}
\]
This profile is continuous on \([-1,1]\), takes values in \([0,1]\), satisfies \(v_h^{\mathrm{dir}}(0)=1\), and vanishes for \(h\le x\le1\). Its full-interval regularity satisfies
\[
|v_h^{\mathrm{dir}}(x)-v_h^{\mathrm{dir}}(t)|
\le\min\{1,|x-t|/h\},
\qquad x,t\in[-1,1],
\]
and hence, for every \(0<\beta\le1\),
\[
[v_h^{\mathrm{dir}}]_\beta\le h^{-\beta}.
\]
The direct latent witnesses use this profile through the conditional means
\[
\mu_0(P_h^{\mathrm{dir},\pm},x)=\frac12,
\qquad
\mu_1(P_h^{\mathrm{dir},\pm},x)
=\frac12\pm\kappa h^\beta v_h^{\mathrm{dir}}(x),
\qquad \kappa=\frac1{100}.
\]
\end{definitionv}

The direct profile \leanref{sym:Directbump}{\ensuremath{v_h^{\mathrm{dir}}}} in \cref{obj:synth_7} enters the following Bernoulli alternative laws.

\begin{definitionv}{def:direct-alternatives}[Direct alternative laws \(P_h^{\mathrm{dir},\pm}\)]
The latent laws \(P_h^{\mathrm{dir},\pm}\) are specified by:
\begin{itemize}
\item \textbf{(Latent score.)}
\[
X\sim\operatorname{Uniform}[-1,1].
\]
\item \textbf{(Potential outcomes.)} Conditional on \(X=x\), the binary potential outcomes are independent Bernoulli variables with conditional means
\[
\mu_0(P_h^{\mathrm{dir},\pm},x)=\frac12,
\qquad
\mu_1(P_h^{\mathrm{dir},\pm},x)
=\frac12\pm\kappa h^\beta v_h^{\mathrm{dir}}(x),
\]
where \(v_h^{\mathrm{dir}}\) is the direct-experiment profile.
\item \textbf{(Measurement error.)} The standardized measurement error \(\varepsilon\) is independent of the latent score and potential outcomes, with
\[
\varepsilon\sim N(0,1).
\]
\end{itemize}
\end{definitionv}

The direct alternatives \leanref{sym:DirectAlt}{\ensuremath{P_h^{\mathrm{dir},\pm}}} in \cref{obj:def:direct-alternatives} pair a common score distribution and untreated mean with a treated-mean perturbation of amplitude \(\kappa h^\beta\). Their role in \cref{obj:thm:direct-reduction} is to link the direct-experiment lower bound to full-interval Hölder regularity through the explicit profile in \cref{obj:synth_7}.

For positive measurement-error scales, the cancellation witnesses instead use a support width that depends jointly on noise and polynomial order. The following rule fixes that width and the associated endpoint resolution before they enter the noisy lower bound.

\begin{definitionv}{def:noise-support}[Support and resolution \(b_m(\sigma),h_m(\sigma)\)]
The noise-dependent support scale \(b_m(\sigma)\) and endpoint resolution \(h_m(\sigma)\) are
\[
b_m(\sigma)=\min\{1,\sigma\sqrt m/8\},
\qquad
h_m(\sigma)=\frac{b_m(\sigma)}{m^2},
\qquad \sigma>0.
\]
The corresponding alternative pair is
\[
P^\pm_{b_m(\sigma),m}.
\]
\end{definitionv}

The support scale \leanref{sym:Support}{\ensuremath{b_m(\sigma)}} in \cref{obj:def:noise-support} stays within the latent interval, while the endpoint resolution \leanref{sym:Resolution}{\ensuremath{h_m(\sigma)}} retains the degree-squared scale established in \cref{obj:lem:legal-cancellation}. This rule specifies the noise-dependent witnesses used in the lower bound of \cref{obj:thm:uniform-frontier}. Their mathematical roles remain separate: \cref{obj:lem:legal-cancellation} supplies benchmark membership and target separation, \cref{obj:lem:full-marked-likelihood} controls the complete observed-law comparison for positive noise, and \cref{obj:lem:two-point-transfer} converts the stated sample-law proximity condition into estimation-risk and honest-length lower bounds for randomized procedures.


\section{Rate comparison, branch formulas, and verification scope}

The support rule in \cref{obj:def:noise-support} completes the specification of the noisy alternatives. This appendix collects the comparisons that connect those alternatives and the public variance certificate to the resolution equation in \cref{obj:def:uniform-rate}. Three auxiliary results control changes in resolution, observable second moments, and cancellation information. Their roles in \cref{obj:thm:uniform-frontier} are followed by the branch interpretation and the larger-class converse.

\subsection{Noise cost and the two envelopes}

The noise cost summarizes the dependence of Gaussian inversion on endpoint resolution. Comparisons across resolutions must remain valid when the direct, intermediate, and compact branches meet. The following bounds provide that common control.

\begin{lemmav}{lem:noise-cost-scaling}[Noise cost scaling bounds]
For \(h\in(0,1]\) and \(\sigma\in[0,1]\), define the noise cost
\[
\Psi(h,\sigma)=
\begin{cases}
0, & \sigma\le h,\\
(\sigma/h)^{2/3}-1, & h<\sigma\ \text{and}\ \sigma^4\le h,\\
h^{-1/2}\bigl[1+\log(\sigma^2/\sqrt h)\bigr]-1,
& h<\sigma\ \text{and}\ h<\sigma^4.
\end{cases}
\]
For every \(\sigma\in[0,1]\), the function \(h\mapsto\Psi(h,\sigma)\) is continuous, nonnegative, and nonincreasing on \((0,1]\).

For all real \(h,t,\sigma\) satisfying \(0<h\le t\le\sigma\le1\),
\[
(t/h)^{1/2}
\le
\frac{1+\Psi(h,\sigma)}{1+\Psi(t,\sigma)}
\le t/h.
\]
For every \(h\in(0,1]\), \(\sigma\in[0,1]\), and real \(a\ge1\) satisfying \(ah\le1\),
\[
1+\Psi(ah,\sigma)
\le
\max\bigl\{1,\ a^{-1/2}[1+\Psi(h,\sigma)]\bigr\}.
\]
\end{lemmav}

\begin{proof}[Proof of \cref{obj:lem:noise-cost-scaling}]
For \(0\le \sigma\le1\), write
\[
F_\sigma(h)=1+\Psi(h,\sigma),\qquad 0<h\le1.
\]
\begin{enumerate}
\item \emph{Continuity.}
Fix \(\sigma\in[0,1]\). If \(\sigma=0\), then
\[
\Psi(h,0)=0,\qquad 0<h\le1,
\]
so continuity follows. Suppose \(\sigma>0\). The power and logarithmic expressions defining the cost are continuous for positive \(h\): their denominators are positive, and the logarithm has positive argument. At their interface,
\[
\sqrt{\sigma^4}=\sigma^2,\qquad
\left(\frac{\sigma}{\sigma^4}\right)^{2/3}
=\sigma^{-2}
=(\sigma^4)^{-1/2}
 \left[1+\log\left(\frac{\sigma^2}{\sqrt{\sigma^4}}\right)\right].
\]
Thus the two expressions for the cost agree at \(h=\sigma^4\).
% lean: CausalSmith.Stat.RdTruesideNoiseFrontier.noiseCost_compact_interface
At the direct-branch interface, \(\sigma^4\le\sigma\), since \(0<\sigma\le1\), and
\[
\left(\frac{\sigma}{\sigma}\right)^{2/3}-1=0.
\]
Gluing the continuous expressions at these interfaces proves continuity on \((0,1]\), including when \(\sigma=1\) and the interfaces coincide.
% lean: CausalSmith.Stat.RdTruesideNoiseFrontier.noiseCost_continuousOn

\item \emph{Nonnegativity.}
Fix \(0<h\le1\) and \(0\le\sigma\le1\). If \(\sigma\le h\), the cost is zero. If \(h<\sigma\) and \(\sigma^4\le h\), then
\[
\frac{\sigma}{h}\ge1,\qquad
\left(\frac{\sigma}{h}\right)^{2/3}\ge1,
\]
so the power-branch cost is nonnegative. In the remaining case, \(h<\sigma^4\), positivity of \(h\) gives
\[
\sqrt h\le\sigma^2,\qquad
\log\left(\frac{\sigma^2}{\sqrt h}\right)\ge0,
\qquad h^{-1/2}\ge1.
\]
Consequently,
\[
h^{-1/2}\left[1+\log\left(\frac{\sigma^2}{\sqrt h}\right)\right]\ge1,
\]
which proves nonnegativity on the logarithmic branch as well.
% lean: CausalSmith.Stat.RdTruesideNoiseFrontier.noiseCost_nonneg

\item \emph{Monotonicity.}
Fix \(0<h\le t\le1\). If \(\sigma\le h\), both costs are zero. If \(h<\sigma\le t\), nonnegativity gives
\[
\Psi(t,\sigma)=0\le\Psi(h,\sigma).
\]
It remains to consider \(h\le t<\sigma\), so \(\sigma>0\).

If \(\sigma^4\le h\), both resolutions lie on the power branch, and
\[
0\le\frac{\sigma}{t}\le\frac{\sigma}{h}
\quad\Longrightarrow\quad
\left(\frac{\sigma}{t}\right)^{2/3}-1
\le
\left(\frac{\sigma}{h}\right)^{2/3}-1.
\]
% lean: CausalSmith.Stat.RdTruesideNoiseFrontier.noiseCost_power_antitone

If \(h<\sigma^4\le t\), compare through the interface:
\[
\begin{aligned}
\Psi(t,\sigma)
&=\left(\frac{\sigma}{t}\right)^{2/3}-1\\
&\le\left(\frac{\sigma}{\sigma^4}\right)^{2/3}-1\\
&=(\sigma^4)^{-1/2}
 \left[1+\log\left(\frac{\sigma^2}{\sqrt{\sigma^4}}\right)\right]-1\\
&\le h^{-1/2}
 \left[1+\log\left(\frac{\sigma^2}{\sqrt h}\right)\right]-1
=\Psi(h,\sigma).
\end{aligned}
\]
The first inequality is the power comparison just established, and the equality is the interface identity. For the last inequality,
\[
(\sigma^4)^{-1/2}\le h^{-1/2},
\qquad
0=\log\left(\frac{\sigma^2}{\sqrt{\sigma^4}}\right)
\le\log\left(\frac{\sigma^2}{\sqrt h}\right);
\]
multiplying the corresponding nonnegative factors gives the required bound.
% lean: CausalSmith.Stat.RdTruesideNoiseFrontier.noiseCost_compact_antitone

Finally, if \(h\le t<\sigma^4\), then
\[
\sqrt t\le\sigma^2,\qquad
0\le\log\left(\frac{\sigma^2}{\sqrt t}\right)
\le\log\left(\frac{\sigma^2}{\sqrt h}\right),
\qquad
t^{-1/2}\le h^{-1/2}.
\]
Here the logarithmic comparison follows from \(\sqrt h\le\sqrt t\) and monotonicity of the logarithm on positive arguments. Multiplication yields
\[
t^{-1/2}\left[1+\log\left(\frac{\sigma^2}{\sqrt t}\right)\right]-1
\le
h^{-1/2}\left[1+\log\left(\frac{\sigma^2}{\sqrt h}\right)\right]-1.
\]
These cases establish that the cost is nonincreasing.
% lean: CausalSmith.Stat.RdTruesideNoiseFrontier.noiseCost_antitoneOn

\item \emph{Ratio bounds.}
Assume \(0<h\le t\le\sigma\le1\). Then \(t,\sigma>0\), and nonnegativity gives
\[
F_\sigma(t)\ge1>0.
\]
We first prove the multiplicative comparisons
\[
\left(\frac th\right)^{1/2}F_\sigma(t)
\le F_\sigma(h)
\le\frac th F_\sigma(t).
\]

If \(\sigma^4\le h\), the power formula holds at both resolutions:
\[
F_\sigma(h)=\left(\frac{\sigma}{h}\right)^{2/3},
\qquad
F_\sigma(t)=\left(\frac{\sigma}{t}\right)^{2/3}.
\]
These identities include \(t=\sigma\), since the direct-branch value there is one. Hence
\[
F_\sigma(h)=\left(\frac th\right)^{2/3}F_\sigma(t).
\]
Because \(t/h\ge1\),
\[
\left(\frac th\right)^{1/2}
\le\left(\frac th\right)^{2/3}
\le\frac th.
\]
Multiplying by \(F_\sigma(t)>0\) proves both comparisons in this case.
% lean: CausalSmith.Stat.RdTruesideNoiseFrontier.noiseCost_power_comparison

Suppose next that \(h<\sigma^4\) and \(t\le\sigma^4\). The logarithmic formula, extended to the interface by the identity in the first step, gives
\[
F_\sigma(h)=h^{-1/2}
 \left[1+\log\left(\frac{\sigma^2}{\sqrt h}\right)\right],
\qquad
F_\sigma(t)=t^{-1/2}
 \left[1+\log\left(\frac{\sigma^2}{\sqrt t}\right)\right].
\]
The same formulas hold whenever \(0<h\le t\le\sigma^4\).
For such resolutions define
\[
\eta_{\mathrm{cmp}}=\frac{\sqrt t}{\sqrt h},
\qquad
\Lambda_{\mathrm{cmp}}=\log\left(\frac{\sigma^2}{\sqrt t}\right).
\]
Then
\[
\eta_{\mathrm{cmp}}\ge1,\qquad
\Lambda_{\mathrm{cmp}}\ge0,\qquad
\log\left(\frac{\sigma^2}{\sqrt h}\right)
=\Lambda_{\mathrm{cmp}}+\log\eta_{\mathrm{cmp}}.
\]
Using \(0\le\log\eta_{\mathrm{cmp}}\le\eta_{\mathrm{cmp}}-1\), we obtain
\[
\begin{aligned}
1+\Lambda_{\mathrm{cmp}}
&\le1+\Lambda_{\mathrm{cmp}}+\log\eta_{\mathrm{cmp}},\\
1+\Lambda_{\mathrm{cmp}}+\log\eta_{\mathrm{cmp}}
&\le\Lambda_{\mathrm{cmp}}+\eta_{\mathrm{cmp}}
\le\eta_{\mathrm{cmp}}(1+\Lambda_{\mathrm{cmp}}).
\end{aligned}
\]
The final inequality uses
\((\eta_{\mathrm{cmp}}-1)\Lambda_{\mathrm{cmp}}\ge0\).
Since
\[
h^{-1/2}=\frac1{\sqrt h},\qquad
t^{-1/2}=\frac1{\sqrt t},\qquad
\eta_{\mathrm{cmp}}=\left(\frac th\right)^{1/2},
\qquad
\eta_{\mathrm{cmp}}^2=\frac th,
\]
multiplication by \(1/\sqrt h\) gives
\[
\left(\frac th\right)^{1/2}F_\sigma(t)
\le F_\sigma(h)
\le\frac th F_\sigma(t).
\]
% lean: CausalSmith.Stat.RdTruesideNoiseFrontier.noiseCost_compact_comparison

The remaining case is \(h<\sigma^4<t\). Apply the logarithmic comparison from \(h\) to \(\sigma^4\), and the power comparison from \(\sigma^4\) to \(t\). As \(\sigma^4>0\), their lower bounds compose as
\[
\begin{aligned}
\left(\frac th\right)^{1/2}F_\sigma(t)
&=\left(\frac{\sigma^4}{h}\right)^{1/2}
  \left(\frac{t}{\sigma^4}\right)^{1/2}F_\sigma(t)\\
&\le\left(\frac{\sigma^4}{h}\right)^{1/2}F_\sigma(\sigma^4)
\le F_\sigma(h),
\end{aligned}
\]
and their upper bounds compose as
\[
F_\sigma(h)
\le\frac{\sigma^4}{h}F_\sigma(\sigma^4)
\le\frac{\sigma^4}{h}\frac{t}{\sigma^4}F_\sigma(t)
=\frac th F_\sigma(t).
\]
% lean: CausalSmith.Stat.RdTruesideNoiseFrontier.noiseCost_comparison_trans
Dividing the comparisons by the positive quantity \(F_\sigma(t)\) proves
\[
\left(\frac th\right)^{1/2}
\le\frac{F_\sigma(h)}{F_\sigma(t)}
\le\frac th.
\]
% lean: CausalSmith.Stat.RdTruesideNoiseFrontier.noiseCost_ratio_bounds

\item \emph{Dilation bound.}
Let \(0<h\le1\), \(0\le\sigma\le1\), \(a\ge1\), and \(ah\le1\). If \(\sigma\le ah\), then
\[
F_\sigma(ah)=1
\le\max\{1,a^{-1/2}F_\sigma(h)\}.
\]
This case includes \(\sigma=0\). Otherwise, \(ah<\sigma\). Since \(a>0\),
\[
0<h\le ah\le1,\qquad F_\sigma(ah)>0.
\]
Applying the lower ratio bound with \(t=ah\) gives
\[
a^{1/2}
\le\frac{F_\sigma(h)}{F_\sigma(ah)},
\qquad
a^{1/2}F_\sigma(ah)\le F_\sigma(h).
\]
Multiplication by \(a^{-1/2}\ge0\), together with
\(a^{-1/2}a^{1/2}=1\), yields
\[
F_\sigma(ah)
\le a^{-1/2}F_\sigma(h)
\le\max\{1,a^{-1/2}F_\sigma(h)\}.
\]
Substituting \(F_\sigma(h)=1+\Psi(h,\sigma)\) proves the last assertion.
% lean: CausalSmith.Stat.RdTruesideNoiseFrontier.noiseCost_dilation_of_lower_ratio
\end{enumerate}
\end{proof}

Continuity keeps the cost consistent at the branch interfaces, while monotonicity orders finer and coarser endpoint resolutions. The ratio bounds quantify this ordering within the noisy region. The final inequality also accommodates a coarsening that reaches the direct region, where the cost vanishes. These properties support comparisons between the continuous resolution and the discrete polynomial choices used by the procedures and witnesses.

For the upper comparison, the relevant quantity is the integrated derivative sum entering the finite-sample radius. The next envelope expresses its dependence on degree and noise through the same cost function.

\begin{lemmav}{lem:variance-cost-envelope}[Uniform variance cost envelope]
Let \(m_L\) be the degree of the endpoint polynomial \(K_L\), and define its integrated derivative cost by
\[
V_L(\sigma)
=\frac34\sum_{j=0}^{m_L}\frac{\sigma^{2j}}{j!}
\int_0^1\bigl(K_L^{(j)}(x)\bigr)^2\,dx.
\]
For \(h>0\), define the noise cost by
\[
\Psi(h,\sigma)=
\begin{cases}
0, & \sigma\le h,\\
(\sigma/h)^{2/3}-1, & \sigma>h\ \text{and}\ \sigma^4\le h,\\
h^{-1/2}\bigl[1+\log(\sigma^2/\sqrt h)\bigr]-1,
& \sigma>h\ \text{and}\ \sigma^4>h.
\end{cases}
\]
There exists an absolute constant \(C>0\) such that, for every integer \(L\ge1\) and every \(\sigma\in[0,1]\),
\[
V_L(\sigma)\le C L^2
\exp\!\left\{C\bigl[1+\Psi(L^{-2},\sigma)\bigr]\right\}.
\]
\end{lemmav}

\begin{proof}[Proof of \cref{obj:lem:variance-cost-envelope}]
\begin{enumerate}
\item We use the classical polynomial identities supplied by
\cref{obj:lem:classical-legendre-facts}. The factorial-series estimate in
\cref{obj:lem:exact-covariance} supplies an absolute constant
\[
C_{\mathrm{fac}}\in(0,\infty)
\]
such that, for every integer \(L\ge1\) and every \(\sigma\in[0,1]\),
\[
m_L=3(L-1),\qquad
V_L(\sigma)\le C_{\mathrm{fac}}L^2
\sum_{j=0}^{m_L}
\frac{(C_{\mathrm{fac}}\sigma^2L^4)^j}{(j!)^3}.
\]
Throughout the proof, write
\[
h_L=L^{-2}>0.
\]
% lean: CausalSmith.Stat.RdTruesideNoiseFrontier.exact_covariance

\item First construct an envelope for the first two noise regimes. Define
\[
C_{\mathrm{nc}}
=\max\{C_{\mathrm{fac}},\,3(C_{\mathrm{fac}}+1)\}>0,
\qquad
\eta_L(\sigma)=\sigma L^2\ge0,
\qquad
\zeta_L(\sigma)=(C_{\mathrm{fac}}+1)\eta_L(\sigma)^{2/3}\ge0.
\]
These definitions apply to every \(L\ge1\) and \(\sigma\in[0,1]\). Since
\[
\bigl(\eta_L(\sigma)^{2/3}\bigr)^3=\eta_L(\sigma)^2,
\qquad
C_{\mathrm{fac}}\le(C_{\mathrm{fac}}+1)^3,
\]
we have
\[
C_{\mathrm{fac}}\sigma^2L^4
\le \zeta_L(\sigma)^3.
\]

For nonnegative summands, the sum of their cubes is bounded by the cube
of their sum. This follows by induction from
\[
(\xi+\omega)^3
=\xi^3+\omega^3+3\xi\omega(\xi+\omega)
\ge \xi^3+\omega^3,
\qquad \xi,\omega\in[0,\infty).
\]
Applying this inequality and then the nonnegative exponential series gives
\[
\begin{aligned}
\sum_{j=0}^{m_L}
\frac{(C_{\mathrm{fac}}\sigma^2L^4)^j}{(j!)^3}
&\le
\sum_{j=0}^{m_L}
\left(\frac{\zeta_L(\sigma)^j}{j!}\right)^3\\
&\le
\left(\sum_{j=0}^{m_L}\frac{\zeta_L(\sigma)^j}{j!}\right)^3\\
&\le \exp\{3\zeta_L(\sigma)\}\\
&\le \exp\{C_{\mathrm{nc}}\eta_L(\sigma)^{2/3}\}.
\end{aligned}
\]
Consequently,
\[
V_L(\sigma)
\le C_{\mathrm{nc}}L^2
\exp\{C_{\mathrm{nc}}\eta_L(\sigma)^{2/3}\}.
\]

If \(\sigma\le h_L\), then
\[
0\le\eta_L(\sigma)\le1,
\qquad
\eta_L(\sigma)^{2/3}\le1
=1+\Psi(h_L,\sigma).
\]
If \(\sigma>h_L\) and \(\sigma^4\le h_L\), the intermediate branch instead gives
\[
1+\Psi(h_L,\sigma)
=(\sigma/h_L)^{2/3}
=\eta_L(\sigma)^{2/3}.
\]
Thus the same constant satisfies
\[
V_L(\sigma)\le C_{\mathrm{nc}}L^2
\exp\{C_{\mathrm{nc}}[1+\Psi(h_L,\sigma)]\}
\quad\text{whenever}\quad
\sigma\le h_L\ \text{or}\ \sigma^4\le h_L.
\]
% lean: CausalSmith.Stat.RdTruesideNoiseFrontier.kernelVariance_noncompact_noise_envelope

\item Next construct the compact-branch envelope. Define the absolute constant
\[
C_{\mathrm{comp}}
=\max\{C_{\mathrm{fac}},\,7+3\log(C_{\mathrm{fac}}+1)\}>0.
\]
Suppose that \(\sigma>h_L\) and \(\sigma^4>h_L\), and define
\[
\eta_{\mathrm{comp}}=\sigma^2L,
\qquad
\alpha_{\mathrm{comp}}=(C_{\mathrm{fac}}+1)\eta_{\mathrm{comp}}.
\]
Multiplication of \(\sigma^4>L^{-2}\) by \(L^2>0\) yields
\[
\eta_{\mathrm{comp}}^2=\sigma^4L^2>1.
\]
Since \(\eta_{\mathrm{comp}}\ge0\), it follows that
\[
\eta_{\mathrm{comp}}\ge1,
\qquad
\alpha_{\mathrm{comp}}\ge1.
\]

We first establish the finite-series estimate used here. For
\[
\alpha\in[1,\infty),
\qquad
N_{\mathrm{cut}}\in\mathbb Z_{\ge0},
\qquad
N_{\mathrm{cut}}\le3L+1,
\]
every integer \(0\le j<N_{\mathrm{cut}}\) satisfies \(j\le3L\).
The nonnegative exponential series gives \(L^j/j!\le e^L\), and therefore
\[
\frac{(\alpha L^3)^j}{(j!)^3}
=\alpha^j\left(\frac{L^j}{j!}\right)^3
\le \alpha^{3L}e^{3L}.
\]
The number of summands satisfies
\[
N_{\mathrm{cut}}\le3L+1\le4L+1\le e^{4L},
\]
where the last inequality is \(1+u\le e^u\) at \(u=4L\). Hence
\[
\sum_{j=0}^{N_{\mathrm{cut}}-1}
\frac{(\alpha L^3)^j}{(j!)^3}
\le N_{\mathrm{cut}}\alpha^{3L}e^{3L}
\le \exp\{L(7+3\log\alpha)\}.
\]
For \(N_{\mathrm{cut}}=0\), the sum is empty and equals zero, so this estimate
also covers that case.

Now \(m_L+1\le3L+1\), and
\[
C_{\mathrm{fac}}\sigma^2L^4
\le \alpha_{\mathrm{comp}}L^3.
\]
Applying the preceding estimate gives
\[
\sum_{j=0}^{m_L}
\frac{(C_{\mathrm{fac}}\sigma^2L^4)^j}{(j!)^3}
\le
\exp\{L(7+3\log\alpha_{\mathrm{comp}})\}.
\]
On this branch,
\[
\sqrt{h_L}=L^{-1},
\qquad
h_L^{-1/2}=L,
\qquad
1+\Psi(h_L,\sigma)
=L\bigl(1+\log\eta_{\mathrm{comp}}\bigr).
\]
Moreover,
\[
\log\alpha_{\mathrm{comp}}
=\log(C_{\mathrm{fac}}+1)+\log\eta_{\mathrm{comp}},
\qquad
\log\eta_{\mathrm{comp}}\ge0.
\]
Since \(\log(C_{\mathrm{fac}}+1)\ge0\), the definition of
\(C_{\mathrm{comp}}\) implies both
\[
7+3\log(C_{\mathrm{fac}}+1)\le C_{\mathrm{comp}},
\qquad
3\le C_{\mathrm{comp}}.
\]
It follows that
\[
\begin{aligned}
L(7+3\log\alpha_{\mathrm{comp}})
&=L\bigl(7+3\log(C_{\mathrm{fac}}+1)
            +3\log\eta_{\mathrm{comp}}\bigr)\\
&\le C_{\mathrm{comp}}L(1+\log\eta_{\mathrm{comp}})\\
&=C_{\mathrm{comp}}[1+\Psi(h_L,\sigma)].
\end{aligned}
\]
Combining these inequalities with the factorial-series estimate and
\(C_{\mathrm{fac}}\le C_{\mathrm{comp}}\) proves
\[
V_L(\sigma)\le C_{\mathrm{comp}}L^2
\exp\{C_{\mathrm{comp}}[1+\Psi(h_L,\sigma)]\}
\quad\text{if}\quad
\sigma>h_L,\ \sigma^4>h_L.
\]
% lean: CausalSmith.Stat.RdTruesideNoiseFrontier.kernelVariance_compact_noise_envelope

\item Choose the common absolute constant
\[
C=\max\{C_{\mathrm{nc}},C_{\mathrm{comp}}\}>0.
\]
For every \(L\ge1\) and \(\sigma\in[0,1]\), the branch identities above show
\[
1+\Psi(h_L,\sigma)
=
\begin{cases}
1,&\sigma\le h_L,\\
(\sigma L^2)^{2/3},&\sigma>h_L,\ \sigma^4\le h_L,\\
L[1+\log(\sigma^2L)],&\sigma>h_L,\ \sigma^4>h_L
\end{cases}
\ \ge0.
\]
In the last branch, nonnegativity follows from \(\sigma^2L\ge1\).
Thus, for every auxiliary constant
\[
C_{\mathrm{aux}}\in(0,C],
\]
we have
\[
C_{\mathrm{aux}}L^2\le CL^2,
\qquad
C_{\mathrm{aux}}[1+\Psi(h_L,\sigma)]
\le C[1+\Psi(h_L,\sigma)].
\]
Monotonicity and positivity of the exponential therefore yield
\[
C_{\mathrm{aux}}L^2
\exp\{C_{\mathrm{aux}}[1+\Psi(h_L,\sigma)]\}
\le
CL^2\exp\{C[1+\Psi(h_L,\sigma)]\}.
\]
% lean: CausalSmith.Stat.RdTruesideNoiseFrontier.endpoint_noiseCost_factor_nonneg

\item Fix \(L\ge1\) and \(\sigma\in[0,1]\). If \(\sigma\le h_L\), apply the
first envelope with \(C_{\mathrm{aux}}=C_{\mathrm{nc}}\).
If \(\sigma>h_L\), split further according to whether
\(\sigma^4\le h_L\). In that case, apply the same envelope; when
\(\sigma^4>h_L\), apply the compact-branch envelope with
\(C_{\mathrm{aux}}=C_{\mathrm{comp}}\). The enlargement inequality in each
case gives
\[
V_L(\sigma)\le CL^2
\exp\{C[1+\Psi(L^{-2},\sigma)]\}.
\]
The first case includes \(\sigma=0\) and all \(\sigma\in[0,1]\) when
\(L=1\); the boundary \(\sigma^4=h_L\) belongs to the second case.
% lean: CausalSmith.Stat.RdTruesideNoiseFrontier.variance_cost_envelope
\end{enumerate}
\end{proof}

The envelope in \cref{obj:lem:variance-cost-envelope} places the derivative cost from \cref{obj:lem:exact-covariance} on the endpoint scale \(L^{-2}\). Together with the approximation bound in \cref{obj:lem:positive-endpoint}, it supplies the two components of the public radius in \cref{obj:synth_6}. Its range includes \(L=1\), whose polynomial degree is zero, and \(\sigma=0\), where the derivative sum retains its order-zero term.

The lower comparison uses the complete observed-law likelihood bound in \cref{obj:lem:full-marked-likelihood}. At the support and resolution specified by \cref{obj:def:noise-support}, its information bound takes a corresponding exponential form.

\begin{lemmav}{lem:cancellation-information-envelope}[Cancellation information envelope]
Let the parameters satisfy:
\begin{itemize}
\item \textbf{(Smoothness.)} \(0<\beta\le1\).
\item \textbf{(Noise scale.)} \(0<\sigma\le1\).
\item \textbf{(Polynomial degree.)} \(m\in\mathbb N\) and \(m\ge2\).
\end{itemize}
Write \(b=b_m(\sigma)\) and \(h=h_m(\sigma)\) as in
\cref{obj:def:noise-support}, and let \(P^+_{b,m}\) and \(P^-_{b,m}\)
be the alternative laws of \cref{obj:def:alternatives}, constructed using
\cref{obj:lem:classical-legendre-facts}. Set the perturbation constant
\(\kappa=1/100\), and define the noise cost by
\[
\Psi(h,\sigma)=
\begin{cases}
0, & \sigma\le h,\\
(\sigma/h)^{2/3}-1, & \sigma>h\ \text{and}\ \sigma^4\le h,\\
h^{-1/2}\bigl(1+\log(\sigma^2/\sqrt h)\bigr)-1,
& \sigma>h\ \text{and}\ \sigma^4>h.
\end{cases}
\]
For probability laws \(\mu,\nu\), use the extended chi-squared divergence
\[
\chi^2(\mu,\nu)=
\begin{cases}
\displaystyle\int\left(\frac{d\mu}{d\nu}-1\right)^2\,d\nu,
& \mu\ll\nu,\\
+\infty, & \mu\not\ll\nu.
\end{cases}
\]
Then the observed laws induced by the two alternatives at noise scale
\(\sigma\) satisfy
\[
\chi^2\bigl((P^+_{b,m})_{\mathrm{obs}},
           (P^-_{b,m})_{\mathrm{obs}}\bigr)
\le
\frac{32}{3}\kappa^2 h^{2\beta+1}
\exp\!\left(-\frac{1+\Psi(h,\sigma)}{32}\right).
\]
\end{lemmav}

\begin{proof}[Proof of \cref{obj:lem:cancellation-information-envelope}]
\begin{enumerate}
\item Set
\[
b=b_m(\sigma)=\min\{1,\sigma\sqrt m/8\},
\qquad
h=h_m(\sigma)=\frac{b}{m^2},
\qquad
\lambda_{b,\sigma}=\frac{b^2}{4\sigma^2},
\qquad
N_m=(2m+1)^2-m^2.
\]
Since \(m\ge2\) and \(\sigma>0\), the support and resolution satisfy
\[
0<b\le1,\qquad h>0.
\]
Moreover, \(b\le\sigma\sqrt m/8\), so
\[
0<\lambda_{b,\sigma}
=\frac{b^2}{4\sigma^2}
\le
\frac{(\sigma\sqrt m/8)^2}{4\sigma^2}
=\frac{m}{256}.
\]
Thus the parameter conditions of \cref{obj:lem:full-marked-likelihood} hold for the alternatives constructed using \cref{obj:lem:classical-legendre-facts}. Its two likelihood bounds give
\[
\chi^2\!\left((P^+_{b,m})_{\mathrm{obs}},
             (P^-_{b,m})_{\mathrm{obs}}\right)
\le
16\kappa^2h^{2\beta}\frac{b}{N_m}
\exp\!\left(\frac{\lambda_{b,\sigma}}6\right)
\sum_{j=m-1}^{\infty}\frac{\lambda_{b,\sigma}^{\,j}}{j!}.
\]
% lean: CausalSmith.Stat.RdTruesideNoiseFrontier.noiseSupport_mem
% lean: CausalSmith.Stat.RdTruesideNoiseFrontier.noiseSupport_likelihoodLambda_le
% lean: CausalSmith.Stat.RdTruesideNoiseFrontier.full_marked_likelihood

\item Define the first retained factorial index by
\[
J_m:=m-1\in\mathbb N,
\qquad J_m\ge1,\qquad J_m\ge m/2.
\]
For every integer \(k\ge0\),
\[
\frac{\lambda_{b,\sigma}^{\,J_m+k+1}}{(J_m+k+1)!}
=
\frac{\lambda_{b,\sigma}^{\,J_m+k}}{(J_m+k)!}
\frac{\lambda_{b,\sigma}}{J_m+k+1},
\qquad
0\le\frac{\lambda_{b,\sigma}}{J_m+k+1}
\le\frac12,
\]
because \(J_m+1=m\) and \(\lambda_{b,\sigma}\le m/256\le m/2\).
Induction on \(k\), followed by summation of the geometric series, therefore yields
\[
\frac{\lambda_{b,\sigma}^{\,J_m+k}}{(J_m+k)!}
\le
\frac{\lambda_{b,\sigma}^{\,J_m}}{J_m!}\,2^{-k},
\qquad
\sum_{j=J_m}^{\infty}\frac{\lambda_{b,\sigma}^{\,j}}{j!}
\le
2\frac{\lambda_{b,\sigma}^{\,J_m}}{J_m!}.
\]
% lean: CausalSmith.Stat.RdTruesideNoiseFrontier.factorial_tail_term_le_geometric
% lean: CausalSmith.Stat.RdTruesideNoiseFrontier.noiseSupport_likelihoodTail_le

For every integer \(j\ge1\), the inequality
\(\log t\le t-1\), applied at \(t=(j+1)/j\), gives
\[
j\{\log(j+1)-\log j\}\le1.
\]
Starting from \(\log(1!)=0\ge-1\), induction gives
\[
\begin{aligned}
\log((j+1)!)
&=\log(j!)+\log(j+1)\\
&\ge j\log j-j+\log(j+1)\\
&\ge (j+1)\log(j+1)-(j+1).
\end{aligned}
\]
Consequently,
\[
\log(J_m!)\ge J_m\log J_m-J_m,
\]
and taking logarithms of the positive factorial term gives
\[
\frac{\lambda_{b,\sigma}^{\,J_m}}{J_m!}
\le
\exp\!\left[
-J_m\log\!\left(\frac{J_m}{\mathrm e\,\lambda_{b,\sigma}}\right)
\right].
\]
% lean: CausalSmith.Stat.RdTruesideNoiseFrontier.log_factorial_lower_bound
% lean: CausalSmith.Stat.RdTruesideNoiseFrontier.factorial_term_le_log_exp

To calibrate this exponent, use
\[
\frac{m}{\lambda_{b,\sigma}}\ge256,
\qquad
\log\!\left(\frac{m}{\lambda_{b,\sigma}}\right)\ge8\log2,
\qquad
\log2\ge1-\frac12=\frac12.
\]
Monotonicity of the logarithm and \(J_m\ge m/2\) imply
\[
\log\!\left(\frac{J_m}{\mathrm e\,\lambda_{b,\sigma}}\right)
\ge
\log\!\left(\frac{m}{\lambda_{b,\sigma}}\right)-\log2-1
\ge7\log2-1\ge0.
\]
Hence
\[
\begin{aligned}
&J_m\log\!\left(\frac{J_m}{\mathrm e\,\lambda_{b,\sigma}}\right)
-\frac{\lambda_{b,\sigma}}6\\
&\qquad\ge
\frac m2\left\{
\log\!\left(\frac{m}{\lambda_{b,\sigma}}\right)-\log2-1
\right\}-\frac m{1536}\\
&\qquad\ge
\frac m8\left\{1+\log\!\left(\frac{m}{\lambda_{b,\sigma}}\right)\right\}.
\end{aligned}
\]
For the last inequality, the difference between its two sides, divided by \(m\), is
\[
\begin{aligned}
&\frac38\log\!\left(\frac{m}{\lambda_{b,\sigma}}\right)
-\frac12\log2-\frac58-\frac1{1536}\\
&\qquad\ge
\frac52\log2-\frac58-\frac1{1536}
\ge\frac58-\frac1{1536}>0.
\end{aligned}
\]
Combining the preceding estimates gives
\[
\begin{aligned}
\exp\!\left(\frac{\lambda_{b,\sigma}}6\right)
\sum_{j=m-1}^{\infty}\frac{\lambda_{b,\sigma}^{\,j}}{j!}
&\le
2\exp\!\left[
\frac{\lambda_{b,\sigma}}6
-J_m\log\!\left(\frac{J_m}{\mathrm e\,\lambda_{b,\sigma}}\right)
\right]\\
&\le
2\exp\!\left[
-\frac m8\left\{1+\log\!\left(\frac{m}{\lambda_{b,\sigma}}\right)\right\}
\right].
\end{aligned}
\]
% lean: CausalSmith.Stat.RdTruesideNoiseFrontier.factorial_exponent_lower_bound
% lean: CausalSmith.Stat.RdTruesideNoiseFrontier.noiseSupport_likelihoodTail_exp_bound

\item The normalization satisfies
\[
N_m=3m^2+4m+1\ge3m^2,
\qquad
\frac{b}{N_m}\le\frac{b}{3m^2}=\frac h3.
\]
Substituting this and the factorial-tail estimate into the likelihood bound gives
\[
\begin{aligned}
\chi^2\!\left((P^+_{b,m})_{\mathrm{obs}},
             (P^-_{b,m})_{\mathrm{obs}}\right)
&\le
16\kappa^2h^{2\beta}\frac h3\,
2\exp\!\left[
-\frac m8\left\{1+\log\!\left(\frac{m}{\lambda_{b,\sigma}}\right)\right\}
\right]\\
&=
\frac{32}{3}\kappa^2h^{2\beta+1}
\exp\!\left[
-\frac m8\left\{1+\log\!\left(\frac{m}{\lambda_{b,\sigma}}\right)\right\}
\right].
\end{aligned}
\]
Since the prefactor is nonnegative and the exponential is increasing, it suffices to establish
\[
1+\Psi(h,\sigma)
\le
4m\left\{1+\log\!\left(\frac{m}{\lambda_{b,\sigma}}\right)\right\}.
\]
We split according to whether \(\sigma\sqrt m/8\ge1\).
% lean: CausalSmith.Stat.RdTruesideNoiseFrontier.noiseSupport_div_blockNorm_le
% lean: CausalSmith.Stat.RdTruesideNoiseFrontier.cancellation_information_degree_bound
% lean: CausalSmith.Stat.RdTruesideNoiseFrontier.cancellation_information_envelope

\item Suppose first that \(1\le\sigma\sqrt m/8\). Then
\[
b=1,\qquad
\sigma^2m\ge64,\qquad
h=\frac1{m^2},\qquad
\lambda_{b,\sigma}=\frac1{4\sigma^2}.
\]
Squaring \(\sigma^2m\ge64\) shows that
\[
\sigma^4m^2\ge4096>1,
\qquad
h<\sigma^4\le\sigma,
\]
where the final inequality follows from \(0<\sigma\le1\).
Thus the compact branch of the cost applies. Using
\[
\sqrt h=\frac1m,\qquad h^{-1/2}=m,
\]
we obtain
\[
\begin{aligned}
1+\Psi(h,\sigma)
&=m\{1+\log(\sigma^2m)\}\\
&\le m\{1+\log(4\sigma^2m)\}\\
&=m\left\{1+\log\!\left(\frac{m}{\lambda_{b,\sigma}}\right)\right\}.
\end{aligned}
\]
The logarithmic comparison follows from
\(0<\sigma^2m\le4\sigma^2m\).
Also, \(\lambda_{b,\sigma}\le m/256\) implies
\[
0\le
m\left\{1+\log\!\left(\frac{m}{\lambda_{b,\sigma}}\right)\right\}
\le
4m\left\{1+\log\!\left(\frac{m}{\lambda_{b,\sigma}}\right)\right\}.
\]
This proves the required cost comparison in the saturated case, including the equality \(\sigma\sqrt m/8=1\).
% lean: CausalSmith.Stat.RdTruesideNoiseFrontier.noiseSupport_saturated_threshold
% lean: CausalSmith.Stat.RdTruesideNoiseFrontier.noiseCost_saturated_eq
% lean: CausalSmith.Stat.RdTruesideNoiseFrontier.noiseCost_saturated_le_degree

\item Suppose now that \(\sigma\sqrt m/8<1\). Direct substitution in the support and likelihood scales gives
\[
b=\frac{\sigma\sqrt m}{8},
\qquad
\sigma^2m=64b^2,
\qquad
\lambda_{b,\sigma}=\frac m{256},
\qquad
\frac{m}{\lambda_{b,\sigma}}=256.
\]
We follow the three branches of the cost definition.

If \(\sigma\le h\), then
\[
1+\Psi(h,\sigma)=1
\le4m(1+\log256),
\]
because \(m\ge2\) and \(\log256\ge0\).
% lean: CausalSmith.Stat.RdTruesideNoiseFrontier.noiseSupport_unsaturated_identities

If \(\sigma>h\) and \(\sigma^4\le h\), then
\[
1+\Psi(h,\sigma)=(\sigma/h)^{2/3}.
\]
Indeed, \(m\ge2\) gives \(\sqrt m\le m<8m^2\), and hence
\[
h=\frac{\sigma\sqrt m}{8m^2}<\sigma.
\]
The scale identity \(\sigma^2m=64b^2\) yields
\[
\left\{(\sigma/h)^{2/3}\right\}^{3}
=(\sigma/h)^2
=\frac{\sigma^2m^4}{b^2}
=64m^3=(4m)^3.
\]
Both quantities whose cubes are equal are positive, so
\[
1+\Psi(h,\sigma)=4m\le4m(1+\log256).
\]
% lean: CausalSmith.Stat.RdTruesideNoiseFrontier.noiseCost_unsaturated_intermediate_eq

Finally, suppose that \(\sigma>h\) and \(h<\sigma^4\). The scale identity gives
\[
\sigma^4m^2=4096b^4,
\qquad
b<\sigma^4m^2=4096b^4.
\]
These inequalities force \(b\ge1/16\): if \(b<1/16\), then
\[
b^3\le\frac1{4096},
\qquad
4096b^4\le b,
\]
contradicting the strict inequality above. Therefore
\[
\sqrt h=\frac{\sqrt b}{m}\ge\frac1{4m},
\qquad
h^{-1/2}=\frac1{\sqrt h}\le4m.
\]
On the other hand, \(0<b\le1\) implies
\[
b^3\le1,\qquad b^4\le b,\qquad b^2\le\sqrt b.
\]
Consequently, the logarithm's argument satisfies
\[
0<
\frac{\sigma^2}{\sqrt h}
=\frac{\sigma^2m}{\sqrt b}
=\frac{64b^2}{\sqrt b}
\le64\le256.
\]
Monotonicity of the logarithm and \(1+\log256\ge0\) now give
\[
\begin{aligned}
1+\Psi(h,\sigma)
&=h^{-1/2}\left\{1+\log\!\left(\frac{\sigma^2}{\sqrt h}\right)\right\}\\
&\le h^{-1/2}(1+\log256)\\
&\le4m(1+\log256).
\end{aligned}
\]
Thus every unsaturated branch satisfies
\[
1+\Psi(h,\sigma)
\le
4m\left\{1+\log\!\left(\frac{m}{\lambda_{b,\sigma}}\right)\right\}.
\]
The weak inequalities in the first two branches include their boundary values.
% lean: CausalSmith.Stat.RdTruesideNoiseFrontier.noiseCost_compact_support_bound
% lean: CausalSmith.Stat.RdTruesideNoiseFrontier.noiseCost_unsaturated_le_degree

\item The cost comparison in both support cases implies
\[
-\frac m8\left\{1+\log\!\left(\frac{m}{\lambda_{b,\sigma}}\right)\right\}
\le-\frac{1+\Psi(h,\sigma)}{32}.
\]
Applying monotonicity of the exponential to the degree bound therefore proves
\[
\chi^2\!\left((P^+_{b,m})_{\mathrm{obs}},
             (P^-_{b,m})_{\mathrm{obs}}\right)
\le
\frac{32}{3}\kappa^2h^{2\beta+1}
\exp\!\left(-\frac{1+\Psi(h,\sigma)}{32}\right).
\]
% lean: CausalSmith.Stat.RdTruesideNoiseFrontier.cancellation_information_envelope
\end{enumerate}
\end{proof}

The bound in \cref{obj:lem:cancellation-information-envelope} concerns the full observed triple, including the assignment and outcome marks. Its prefactor combines the squared perturbation amplitude with the endpoint resolution, while its exponential factor describes the information reduction associated with cancellation. Membership and target separation are supplied separately by \cref{obj:lem:legal-cancellation}. The degree range starts at \(m=2\), preserving the lowest admissible cancellation witness, and the positive-noise condition identifies the domain of this envelope.

\subsection{Uniform comparison and boundary cases}

The upper and lower comparisons in \cref{obj:thm:uniform-frontier} connect these envelopes to the balance equation. For the upper comparison, \cref{obj:thm:finite-certificate} already supplies measurable estimation, uniform coverage, and expected-length control by the selected public radius. The remaining rate comparison concerns that deterministic radius, using the approximation and variance bounds at integer polynomial indices.

The finite candidate set in \cref{obj:def:public-selection} is part of this comparison. Positive indices are restricted by the public cap, ties are resolved by the smallest minimizer, and index zero contributes the radius \(1/2\) specified in \cref{obj:synth_6}. Thus the upper bound in \cref{obj:thm:uniform-frontier} applies to the stated finite selection rule. The scaling control in \cref{obj:lem:noise-cost-scaling} provides the resolution comparisons needed for discrete indices, while the constant candidate supplies a valid certificate at coarse resolutions.

For the lower comparison at positive noise, \cref{obj:lem:cancellation-information-envelope} supplies observed-law proximity for the integer-degree alternatives of \cref{obj:def:noise-support}. The target separation and decision transfer in \cref{obj:lem:legal-cancellation,obj:lem:two-point-transfer} give the corresponding estimation-risk and honest-length bounds. The direct alternatives in \cref{obj:def:direct-alternatives} supply the direct-experiment component recorded in \cref{obj:thm:direct-reduction}. Both decision comparisons include the independent procedure seed.

Zero noise belongs to the exact direct branch of \cref{obj:thm:uniform-frontier}, with attainment and lower bounds also stated in \cref{obj:thm:direct-reduction}. At the direct interface, the resolution equals both the noise scale and \(n^{-1/(2\beta+1)}\). At the intermediate--compact interface, it equals \(\sigma^4\). Continuity in \cref{obj:lem:noise-cost-scaling} and the common constants in \cref{obj:thm:uniform-frontier} keep these boundary cases within the simultaneous comparison.

\subsection{Scalar branches and converse transfer}

The scalar coordinates in \cref{obj:thm:uniform-frontier} express the noisy resolution in two ways. On the intermediate branch, \(z=(\sigma/h_\star)^{2/3}\) measures the noise-to-resolution ratio. On the compact branch, \(\tau=h_\star^{-1/2}\) measures inverse square-root resolution. Their equations and admissible ranges specify the corresponding portions of the same balance equation. At the common boundary, both coordinates equal \(\sigma^{-2}\).

The sequence specializations in \cref{obj:thm:uniform-frontier} retain their stated quantifiers. For polynomially shrinking noise with a fixed exponent below \(1/(2\beta+1)\), the intermediate asymptotic order combines the noise power with a logarithmic improvement. At or above that exponent, the exact direct branch applies for every \(n\ge2\). Fixed positive noise eventually selects the compact branch. Uniformity of the original comparison also covers sequences that pass between branches, with constants independent of sample size and noise scale.

The larger-class implication is the target- and experiment-preserving embedding in \cref{obj:prop:published-class-converse-transfer}. The benchmark density and conditional-mean restrictions imply the local conditions of that class, and independence from the latent score and both potential outcomes implies the required joint independence from the score and observed outcome. Preservation of the entire randomized experiment makes the benchmark risk and honest-length lower bounds applicable to the larger decision problem. For positive noise, the published conditions invoked by this interpretation are Assumptions~1Y, 2C, and 7 of \citet{PeiShen2017DevilTails}.

\subsection{Direct-experiment specialization}

The direct alternatives and the zero-noise certificate identify the ordinary endpoint rate within the same benchmark experiment. The following result records a lower bound for the minimax absolute risk \leanref{sym:Risk}{\ensuremath{R_\beta(n,\sigma)}} and minimax honest expected length \leanref{sym:Length}{\ensuremath{\mathcal L_\beta(n,\sigma)}} at every admissible noise scale, together with attainment by the selected procedures at zero noise.

\begin{theoremv}{thm:direct-reduction}[Direct-experiment rates and attainment]
Fix \(0<\beta\le1\), where \(\beta\) is the Hölder exponent. Write \(R_\beta(n,\sigma)\) and \(\mathcal L_\beta(n,\sigma)\) for the real conversions of the extended minimax absolute risk and extended minimax honest expected length, respectively. Their extended objectives are
\[
\inf_{\widehat\theta}\sup_{P\in\mathcal M_\beta(\sigma)}
\mathbb E_{P,n}\!\left[|\widehat\theta-\theta(P)|\right],
\qquad
\inf_{I\in\mathcal I_{\beta,n,\sigma}}
\sup_{P\in\mathcal M_\beta(\sigma)}
\mathbb E_{P,n}\!\left[\operatorname{length}(I)\right].
\]
Here the first infimum ranges over admissible estimators, \(\mathcal M_\beta(\sigma)\) is the benchmark model class, \(\theta(P)\) is the cutoff treatment contrast, and \(\mathbb E_{P,n}\) denotes expectation under the joint law of the observed sample and independent procedure seed. The honest interval class \(\mathcal I_{\beta,n,\sigma}\) is defined in \cref{obj:def:honest-class}. Both minimizations are performed in the extended nonnegative reals before conversion to real numbers.

There exist constants \(c,C>0\), depending only on \(\beta\), such that, for every integer sample size \(n\ge2\) and every noise scale \(\sigma\in[0,1]\),
\[
R_\beta(n,\sigma)\ge c\,n^{-\beta/(2\beta+1)},
\qquad
\mathcal L_\beta(n,\sigma)\ge c\,n^{-\beta/(2\beta+1)}.
\]
At \(\sigma=0\),
\[
R_\beta(n,0)\le C\,n^{-\beta/(2\beta+1)},
\qquad
\mathcal L_\beta(n,0)\le C\,n^{-\beta/(2\beta+1)}.
\]
The selected estimator \(\widehat\theta_{L_\star}\) and selected interval \(I_{L_\star}\), evaluated at \((\beta,n,0)\), satisfy the extended worst-case bounds
\[
\sup_{P\in\mathcal M_\beta(0)}
\mathbb E_{P,n}\!\left[
|\widehat\theta_{L_\star}-\theta(P)|
\right]
\le C\,n^{-\beta/(2\beta+1)},
\qquad
\sup_{P\in\mathcal M_\beta(0)}
\mathbb E_{P,n}\!\left[
\operatorname{length}(I_{L_\star})
\right]
\le C\,n^{-\beta/(2\beta+1)}.
\]
Moreover, \(I_{L_\star}\in\mathcal I_{\beta,n,0}\), the honest interval class of \cref{obj:def:honest-class}.
\end{theoremv}

At zero noise, the endpoint approximation and the sampling variation in the public certificate recover the Hölder endpoint rate. The matching lower bounds identify that rate as the benchmark difficulty for both estimation and honest connected intervals. Their validity throughout \(\sigma\in[0,1]\) supplies the direct-experiment baseline for the uniform comparison, while the zero-noise attainment clauses connect that baseline to the observable procedures already specified in \cref{obj:def:public-selection}.

\subsection{Precision summary}

The established relationships can now be collected around the causal target, the public certificate, and the two families of alternatives. The Lean component definitions are total on their ambient natural-number and real domains, including the index-zero fallback and values outside the statistical domain. The guarantees below apply only for \(n\ge2\) and \(\sigma\in[0,1]\), while positive kernel indices and cancellation degrees use the ranges stated in their defining environments. The following object records these components with those domain distinctions and their defining references.

\begin{remarkv}{def:frontier-handle}[Endpoint precision handle]
The endpoint precision handle comprises the following objects, for \(0<\beta\le1\), \(n\in\mathbb N\), and \(\sigma\in\mathbb R\), given the classical Legendre identities in \cref{obj:lem:classical-legendre-facts}.
\begin{itemize}
\item \textbf{(Target.)} The cutoff causal contrast
\[
\theta(P)=\mu_1(P,0)-\mu_0(P,0).
\]
\item \textbf{(Variance.)} For every \(L\in\mathbb N\), the finite derivative sum
\[
m_L=\deg K_L,\qquad
V_L(\sigma)=\frac34\sum_{j=0}^{m_L}
\frac{\sigma^{2j}}{j!}
\int_0^1\bigl(K_L^{(j)}(x)\bigr)^2\,dx,
\]
where \(K_L\) is defined in \cref{obj:def:endpoint-kernel}.
\item \textbf{(Certificate.)} The selected public radius
\[
E_\beta(n,\sigma)=\rho_{L_\star},
\]
with the radius and selected index specified in \cref{obj:def:public-selection}.
\item \textbf{(Cancellation witnesses.)} For every integer \(m\ge2\), when \(0<\sigma\le1\), and for either sign, the alternative laws
\[
b_m(\sigma)=\min\left\{1,\frac{\sigma\sqrt m}{8}\right\},
\qquad
P^\pm_{b_m(\sigma),m},
\]
defined in \cref{obj:def:alternatives,obj:def:noise-support}.
\item \textbf{(Direct witnesses.)} For every \(0<h\le1\) and either sign, the alternative laws \(P_h^{\mathrm{dir},\pm}\) defined in \cref{obj:def:direct-alternatives}.
\item \textbf{(Rate and procedures.)} The resolution rate, selected measurable estimator, and selected interval procedure
\[
r_\beta(n,\sigma),\qquad
\widetilde\theta_{\beta,n,\sigma},\qquad
\widetilde I_{\beta,n,\sigma},
\]
defined in \cref{obj:def:uniform-rate}.
\end{itemize}
\end{remarkv}

The precision handle in \cref{obj:def:frontier-handle} collects quantities already defined and used above. The next proposition consolidates their relationships, including the attainable public radius, the minimax comparisons, and the scalar branch equations of \cref{obj:thm:uniform-frontier}.

\begin{propositionv}{thm:uniform-frontier-resolution}[Uniform precision frontier]
Fix \(\beta\in(0,1]\), and write \(\nu=2\beta+1\). For \(n\ge2\) and \(\sigma\in[0,1]\), put \(h_0=n^{-1/\nu}\) and define
\[
\Psi(h,\sigma)=
\begin{cases}
0, & \sigma\le h,\\
(\sigma/h)^{2/3}-1, & \sigma^4\le h<\sigma,\\
h^{-1/2}\bigl\{1+\log(\sigma^2h^{-1/2})\bigr\}-1,
    & 0<h<\sigma^4.
\end{cases}
\]
There is a unique \(h_\star(n,\sigma)\in[h_0,1]\) satisfying
\[
\nu\log\!\left(\frac1{h_\star(n,\sigma)}\right)
+\Psi\bigl(h_\star(n,\sigma),\sigma\bigr)=\log n.
\]
Set
\[
r_\beta(n,\sigma)=h_\star(n,\sigma)^\beta,\qquad
\widetilde\theta_{\beta,n,\sigma}=\widehat\theta_{L_\star},
\qquad
\widetilde I_{\beta,n,\sigma}=I_{L_\star}.
\]

There exist constants \(c,C>0\), depending on \(\beta\) and independent of \(n\) and \(\sigma\), such that, simultaneously for every integer \(n\ge2\) and every \(\sigma\in[0,1]\),
\[
\begin{aligned}
c\,r_\beta(n,\sigma)
&\le R_\beta(n,\sigma)
\le \sup_{P\in\mathcal M_\beta(\sigma)}
   \mathbb E_{P,n}
   \left|\widetilde\theta_{\beta,n,\sigma}-\theta(P)\right|
\le E_\beta(n,\sigma)
\le C\,r_\beta(n,\sigma),\\
c\,r_\beta(n,\sigma)
&\le \mathcal L_\beta(n,\sigma)
\le \sup_{P\in\mathcal M_\beta(\sigma)}
   \mathbb E_{P,n}
   \operatorname{length}\bigl(\widetilde I_{\beta,n,\sigma}\bigr)
\le 2E_\beta(n,\sigma)
\le C\,r_\beta(n,\sigma),\\
\widetilde I_{\beta,n,\sigma}
&\in\mathcal I_{\beta,n,\sigma}.
\end{aligned}
\]
The estimation and interval comparisons range over the declared measurable estimator and honest connected-interval decision classes, including procedures with an independent uniform random seed.

The resolution has the following branches:
\begin{itemize}
\item \textbf{(Direct branch.)}
If \(\sigma\le n^{-1/\nu}\), then
\[
h_\star(n,\sigma)=n^{-1/\nu}.
\]

\item \textbf{(Intermediate branch.)}
If \(\sigma>n^{-1/\nu}\) and
\[
\log(n\sigma^{4\nu})\le \sigma^{-2}-1,
\]
then
\[
h_\star(n,\sigma)=\sigma z^{-3/2},
\qquad
z+\frac{3\nu}{2}\log z=1+\log(n\sigma^\nu),
\qquad
1<z\le\sigma^{-2},
\]
where the displayed equation uniquely determines \(z\).

\item \textbf{(Compact branch.)}
If \(\sigma>n^{-1/\nu}\) and
\[
\log(n\sigma^{4\nu})\ge \sigma^{-2}-1,
\]
then
\[
h_\star(n,\sigma)=\tau^{-2},
\qquad
2\nu\log\tau+
\tau\bigl\{1+\log(\sigma^2\tau)\bigr\}=\log n+1,
\qquad
\tau\ge\sigma^{-2},
\]
where the displayed equation uniquely determines \(\tau\).
\end{itemize}
The intermediate and compact equations agree at
\(h_\star=\sigma^4\); the direct interface is
\(h_\star=\sigma=n^{-1/\nu}\). The bounds use the same constants along every sequence \(\sigma=\sigma_n\in[0,1]\), including shrinking-noise sequences.

In particular, for \(\beta=1\) and \(\sigma=n^{-1/6}\), the intermediate branch applies for every \(n\ge2\), and
\[
r_1(n,n^{-1/6})
\asymp
\frac{n^{-1/6}}{\bigl[1+\tfrac12\log n\bigr]^{3/2}},
\]
with comparison constants independent of \(n\). Throughout, the observed side label is \(D=\mathbf1\{X\ge0\}\).
\end{propositionv}

The comparisons in \cref{obj:thm:uniform-frontier-resolution} place minimax risk, honest expected length, and the selected public radius on the same resolution scale. The variance and cancellation envelopes connect attainable precision and statistical indistinguishability through the common noise cost. Because the constants apply simultaneously over the stated sample sizes and noise scales, the branch formulas describe both fixed-noise and shrinking-noise experiments with the same comparisons.

\subsection{Verification scope}

The statements and proofs are machine-checked in Lean 4, including the finite-sample certificate, the auxiliary bounds, the uniform and direct-experiment comparisons, and the converse transfer. Their statistical scope uses the Gaussian calibration, independence, bounded-support density restrictions, conditional-mean bounds, and full-interval Hölder restrictions in \cref{obj:def:model}, together with the declared sampling and decision classes. Lean definitions are total on their ambient types; statistical guarantees use the stated domains \(n\ge2\), \(\sigma\in[0,1]\), and the specified positive-index ranges, with index-zero procedures defined separately. The representation uses finite polynomial inverse-heat sums, induced observed and product laws, and extended nonnegative minimax objectives before real conversion. The classical Legendre and Hermite identities are established by verified declarations, and the supplied theorem-local verification metadata records empty external proof-dependency lists. The citation to \citet[Assumptions~1Y, 2C, and 7]{PeiShen2017DevilTails} supplies the published conditions used for the larger-class interpretation.


\section{Proofs of the main results}\label{sec:deferred-proofs}

\begin{proof}[Proof of \cref{obj:thm:finite-certificate}]
Fix the public parameters \((\beta,n,\sigma)\) satisfying the stated hypotheses. We first establish the three fixed-index guarantees simultaneously for every integer \(J\ge1\) and every \(P'\in\mathcal M_\beta(\sigma)\). For such \(J\) and \(P'\), define
\[
\begin{aligned}
A_{d,J}(P')
&=\int_0^1 K_J(x)f(P',s_dx)\mu_d(P',s_dx)\,dx,\\
B_{d,J}(P')
&=\int_0^1 K_J(x)f(P',s_dx)\,dx,\\
\eta_{d,J}(P')
&=\frac{A_{d,J}(P')}{B_{d,J}(P')},
\qquad d\in\{0,1\}.
\end{aligned}
\]
Here division by zero is assigned value zero; the positive denominator bound established below ensures that these are ordinary ratios throughout the model class. The signs \(s_1=1\) and \(s_0=-1\) are those of \cref{obj:synth_5}.

\begin{enumerate}
\item \textbf{Population identities and empirical moments.}
Define the single-observation summands, as functions on
\(\{0,1\}^2\times\mathbb R\), by
\[
\begin{aligned}
\xi^{\mathrm A}_{d,J}(D,Y,W)
&=\mathbf1\{D=d\}YQ_{J,\sigma}(s_dW),\\
\xi^{\mathrm B}_{d,J}(D,Y,W)
&=\mathbf1\{D=d\}Q_{J,\sigma}(s_dW).
\end{aligned}
\]
The inverse-heat polynomial is defined in \cref{obj:def:inverse-heat}.
Its evaluation at the proxy is square integrable: the latent score is bounded, Gaussian variables have every polynomial moment, and the indicator and binary outcome are bounded. Consequently,
\[
\xi^{\mathrm A}_{d,J},\xi^{\mathrm B}_{d,J}
\in L^2(P'_{\mathrm{obs}}).
\]

By \cref{obj:ass:gaussian,obj:ass:error-independence}, the reflected error \(s_d\varepsilon\) is standard Gaussian and independent of \((X,Y^0,Y^1)\). The first identity of \cref{obj:lem:exact-covariance} therefore gives, for every real \(x\),
\[
\int_{\mathbb R}Q_{J,\sigma}(s_dx+\sigma e)\,N(0,1)(de)
=K_J(s_dx).
\]
Integrating the independent Gaussian coordinate first, and using
\(\mathbf1\{D=d\}Y=\mathbf1\{D=d\}Y^d\), yields
\[
\begin{aligned}
\int \xi^{\mathrm B}_{d,J}\,dP'_{\mathrm{obs}}
&=\int \mathbf1\{D=d\}K_J(s_dX)\,dP',\\
\int \xi^{\mathrm A}_{d,J}\,dP'_{\mathrm{obs}}
&=\int \mathbf1\{D=d\}Y^dK_J(s_dX)\,dP'.
\end{aligned}
\]
The score marginal has density \(f(P',\cdot)\), so
\[
P'\{X=0\}=\int_{\{0\}}f(P',x)\,dx=0.
\]
Thus the arm indicator agrees almost surely with
\(\mathbf1\{s_dX\in[0,1]\}\). The defining conditional-mean identity gives the score marginal weighted by \(Y^d\) the density
\(\mu_d(P',x)f(P',x)\). Integrating the bounded arm weight against these two marginal identities gives
\[
\begin{aligned}
\int \xi^{\mathrm B}_{d,J}\,dP'_{\mathrm{obs}}
&=\int_{\{x:s_dx\in[0,1]\}}K_J(s_dx)f(P',x)\,dx
=B_{d,J}(P'),\\
\int \xi^{\mathrm A}_{d,J}\,dP'_{\mathrm{obs}}
&=\int_{\{x:s_dx\in[0,1]\}}
K_J(s_dx)\mu_d(P',x)f(P',x)\,dx
=A_{d,J}(P').
\end{aligned}
\]
For \(d=1\) the integration interval is \([0,1]\); for \(d=0\) it is \([-1,0]\), and substitution \(x\mapsto-x\) gives the displayed endpoint integrals.

For the second moments, the second identity of
\cref{obj:lem:exact-covariance}, integrated in the same order, gives
\[
\begin{aligned}
\int (\xi^{\mathrm B}_{d,J})^2\,dP'_{\mathrm{obs}}
&=\int_0^1 f(P',s_dx)
\sum_{j=0}^{m_J}\frac{\sigma^{2j}}{j!}
[K_J^{(j)}(x)]^2\,dx\\
&\le\frac34\int_0^1
\sum_{j=0}^{m_J}\frac{\sigma^{2j}}{j!}
[K_J^{(j)}(x)]^2\,dx
=V_J(\sigma),
\end{aligned}
\]
where the inequality uses \cref{obj:ass:density-bounds} and the nonnegativity of every summand. Since \(Y\in\{0,1\}\),
\[
(\xi^{\mathrm A}_{d,J})^2
=Y^2(\xi^{\mathrm B}_{d,J})^2
\le(\xi^{\mathrm B}_{d,J})^2,
\qquad
\int(\xi^{\mathrm A}_{d,J})^2\,dP'_{\mathrm{obs}}
\le V_J(\sigma).
\]

Define the centered empirical errors on the sample and seed space by
\[
\begin{aligned}
\Delta^{\mathrm A}_{d,J}(S_n,U)
&=\widehat A_{d,J}(S_n)-A_{d,J}(P'),\\
\Delta^{\mathrm B}_{d,J}(S_n,U)
&=\widehat B_{d,J}(S_n)-B_{d,J}(P').
\end{aligned}
\]
The empirical averages are those of \cref{obj:def:ratio}. For either
\(\diamond\in\{\mathrm A,\mathrm B\}\), the population identities show that
\[
\Delta^\diamond_{d,J}
=\frac1n\sum_{i=1}^n
\left[
\xi^\diamond_{d,J}(D_i,Y_i,W_i)
-\int\xi^\diamond_{d,J}\,dP'_{\mathrm{obs}}
\right].
\]
Independence makes the cross terms between distinct centered observations vanish. Expanding the square therefore gives
\[
\begin{aligned}
\mathbb E_{P',n}\Delta^\diamond_{d,J}&=0,\\
\mathbb E_{P',n}(\Delta^\diamond_{d,J})^2
&=\frac1n\int
\left(\xi^\diamond_{d,J}
-\int\xi^\diamond_{d,J}\,dP'_{\mathrm{obs}}\right)^2
dP'_{\mathrm{obs}}\\
&=\frac1n\left[
\int(\xi^\diamond_{d,J})^2\,dP'_{\mathrm{obs}}
-\left(\int\xi^\diamond_{d,J}\,dP'_{\mathrm{obs}}\right)^2
\right]\\
&\le\frac{V_J(\sigma)}n.
\end{aligned}
\]
These errors are square integrable. Integrating the independent seed contributes a factor one, since the averages depend on the sample. The same mean calculation proves
\[
\mathbb E_{P',n}\widehat B_{d,J}=B_{d,J}(P').
\]
% lean: CausalSmith.Stat.RdTruesideNoiseFrontier.marked_observed_integral_eq_endpoint
% lean: CausalSmith.Stat.RdTruesideNoiseFrontier.unmarked_observed_integral_eq_endpoint
% lean: CausalSmith.Stat.RdTruesideNoiseFrontier.marked_observed_sq_le_kernelVariance
% lean: CausalSmith.Stat.RdTruesideNoiseFrontier.unmarked_observed_sq_le_kernelVariance
% lean: CausalSmith.Stat.RdTruesideNoiseFrontier.seeded_iid_centered_moments
% lean: CausalSmith.Stat.RdTruesideNoiseFrontier.unmarkedMean_eq_single_observation

\item \textbf{Positive denominators and population bias.}
By \cref{obj:lem:positive-endpoint}, the kernel is nonnegative on \([0,1]\) and integrates to one. Hence \cref{obj:ass:density-bounds} gives
\[
\frac14
=\frac14\int_0^1K_J(x)\,dx
\le B_{d,J}(P')
\le\frac34\int_0^1K_J(x)\,dx
=\frac34.
\]
In particular, the denominator is positive and
\[
A_{d,J}(P')=\eta_{d,J}(P')B_{d,J}(P').
\]
The nonnegative weight \(K_J(x)f(P',s_dx)\), together with
\cref{obj:ass:mean-bounds}, gives
\[
\frac14B_{d,J}(P')
\le A_{d,J}(P')
\le\frac34B_{d,J}(P'),
\qquad
\eta_{d,J}(P')\in[1/4,3/4].
\]

For \(x\in[0,1]\), reflection preserves the distance from zero. Thus
\cref{obj:ass:mean-holder} gives
\[
|\mu_d(P',s_dx)-\mu_d(P',0)|\le2x^\beta.
\]
Subtracting the cutoff mean inside the weighted integral yields
\[
\eta_{d,J}(P')-\mu_d(P',0)
=
\frac{
\int_0^1K_J(x)f(P',s_dx)
[\mu_d(P',s_dx)-\mu_d(P',0)]\,dx
}{B_{d,J}(P')}.
\]
The numerator satisfies
\[
\begin{aligned}
\left|
\int_0^1K_J(x)f(P',s_dx)
[\mu_d(P',s_dx)-\mu_d(P',0)]\,dx
\right|
&\le\int_0^1K_J(x)f(P',s_dx)\,2x^\beta\,dx\\
&\le\frac32M_{J,\beta}.
\end{aligned}
\]
Since \(M_{J,\beta}\ge0\) and \(B_{d,J}(P')\ge1/4\), it follows that
\[
|\eta_{d,J}(P')-\mu_d(P',0)|\le6M_{J,\beta}.
\]
Together with the expectation identity above, this also proves the required lower bound on each unmarked population expectation.
% lean: CausalSmith.Stat.RdTruesideNoiseFrontier.endpointKernel_weighted_integral_bounds
% lean: CausalSmith.Stat.RdTruesideNoiseFrontier.arm_kernel_mean_bounds_and_bias

\item \textbf{Absolute risk.}
For every sample, define the floored denominator
\[
\widehat B^{\mathrm{floor}}_{d,J}
=\max\{\widehat B_{d,J},1/8\}.
\]
If \(\widehat B_{d,J}\ge1/8\), flooring leaves it unchanged. If
\(\widehat B_{d,J}<1/8\), then \(B_{d,J}(P')\ge1/4\) gives
\[
\left|\widehat B^{\mathrm{floor}}_{d,J}-B_{d,J}(P')\right|
=B_{d,J}(P')-\frac18
\le B_{d,J}(P')-\widehat B_{d,J}
=\left|\widehat B_{d,J}-B_{d,J}(P')\right|.
\]
Consequently, in both cases,
\[
\left|\widehat B^{\mathrm{floor}}_{d,J}-B_{d,J}(P')\right|
\le|\Delta^{\mathrm B}_{d,J}|,
\qquad
\widehat B^{\mathrm{floor}}_{d,J}\ge\frac18.
\]

Projection onto \([1/4,3/4]\) decreases distance to
\(\eta_{d,J}(P')\), which belongs to that interval. Using the population product identity,
\[
\frac{\widehat A_{d,J}}{\widehat B^{\mathrm{floor}}_{d,J}}
-\eta_{d,J}(P')
=
\frac{
\Delta^{\mathrm A}_{d,J}
+\eta_{d,J}(P')
[B_{d,J}(P')-\widehat B^{\mathrm{floor}}_{d,J}]
}{\widehat B^{\mathrm{floor}}_{d,J}}.
\]
Since \(|\eta_{d,J}(P')|\le3/4\), the projected estimate of
\cref{obj:def:ratio} satisfies
\[
|\widehat\mu_{d,J}-\eta_{d,J}(P')|
\le8|\Delta^{\mathrm A}_{d,J}|
+6|\Delta^{\mathrm B}_{d,J}|.
\]
This inequality applies to every sample, including those with an empty arm.

The triangle inequality for the two arm estimates and their population biases now gives
\[
\begin{aligned}
|\widehat\theta_J-\theta(P')|
\le{}&12M_{J,\beta}\\
&+8\bigl(|\Delta^{\mathrm A}_{1,J}|
+|\Delta^{\mathrm A}_{0,J}|\bigr)
+6\bigl(|\Delta^{\mathrm B}_{1,J}|
+|\Delta^{\mathrm B}_{0,J}|\bigr).
\end{aligned}
\]
For each of the four errors, Cauchy--Schwarz on the probability experiment gives
\[
\mathbb E_{P',n}|\Delta^\diamond_{d,J}|
\le\sqrt{\mathbb E_{P',n}(\Delta^\diamond_{d,J})^2}
\le\sqrt{\frac{V_J(\sigma)}n}.
\]
The displayed upper bound on the effect error is integrable, so integration yields
\[
\mathbb E_{P',n}|\widehat\theta_J-\theta(P')|
\le12M_{J,\beta}
+28\sqrt{\frac{V_J(\sigma)}n}.
\]
% lean: CausalSmith.Stat.RdTruesideNoiseFrontier.denominator_floor_error_le
% lean: CausalSmith.Stat.RdTruesideNoiseFrontier.clipped_ratio_error_le
% lean: CausalSmith.Stat.RdTruesideNoiseFrontier.thetaHat_absRisk_le_of_second_moments

\item \textbf{Fixed-index coverage.}
Kernel normalization implies
\[
\int_0^1K_J(x)^2\,dx>0.
\]
Indeed, if this nonnegative integral were zero, then \(K_J=0\) almost everywhere on \([0,1]\), contradicting
\(\int_0^1K_J(x)\,dx=1\). Every derivative term in \(V_J(\sigma)\) is nonnegative, and its order-zero coefficient is one, including at \(\sigma=0\). Hence
\[
V_J(\sigma)\ge\frac34\int_0^1K_J(x)^2\,dx>0.
\]
Define the positive threshold
\[
\zeta_J=\sqrt{\frac{40V_J(\sigma)}n}.
\]
The centered-moment identities imply, for every \(d\) and
\(\diamond\in\{\mathrm A,\mathrm B\}\),
\[
\operatorname{Var}_{P',n}(\Delta^\diamond_{d,J})
=\mathbb E_{P',n}(\Delta^\diamond_{d,J})^2
\le\frac{V_J(\sigma)}n.
\]
Chebyshev's inequality therefore gives
\[
\begin{aligned}
\mathbb P_{P',n}\{|\Delta^\diamond_{d,J}|>\zeta_J\}
&\le\mathbb P_{P',n}\{|\Delta^\diamond_{d,J}|\ge\zeta_J\}\\
&\le\frac{\operatorname{Var}_{P',n}(\Delta^\diamond_{d,J})}
{\zeta_J^2}
\le\frac1{40}.
\end{aligned}
\]
Define the simultaneous-control event on the sample and seed space by
\[
\Gamma_J
=\bigcap_{d\in\{0,1\}}
\left(
\{|\Delta^{\mathrm A}_{d,J}|\le\zeta_J\}
\cap
\{|\Delta^{\mathrm B}_{d,J}|\le\zeta_J\}
\right).
\]
The errors are measurable, and the union bound gives
\[
\mathbb P_{P',n}(\Gamma_J)
\ge1-\sum_{d\in\{0,1\}}
\sum_{\diamond\in\{\mathrm A,\mathrm B\}}
\mathbb P_{P',n}\{|\Delta^\diamond_{d,J}|>\zeta_J\}
\ge1-\frac4{40}=\frac9{10}.
\]
On this event the two arm errors relative to their weighted population means are each at most \(14\zeta_J\). Adding the two population biases gives
\[
|\widehat\theta_J-\theta(P')|
\le12M_{J,\beta}+28\zeta_J
=\rho_J,
\]
where \(\rho_J\) is the public radius of \cref{obj:synth_6}.

The cutoff mean bounds imply
\[
-\frac12\le\theta(P')\le\frac12.
\]
Thus, on \(\Gamma_J\),
\[
\max\{-1,\widehat\theta_J-\rho_J\}
\le\theta(P')
\le\min\{1,\widehat\theta_J+\rho_J\}.
\]
These are precisely the endpoints of \(I_J\) in
\cref{obj:synth_6}. Therefore
\[
\mathbb P_{P',n}\{\theta(P')\in I_J\}\ge\frac9{10}.
\]
Since \(J\ge1\) and \(P'\) were arbitrary, substituting \(J=L\) and \(P'=P\) proves all three fixed-index assertions.
% lean: CausalSmith.Stat.RdTruesideNoiseFrontier.kernelVariance_pos
% lean: CausalSmith.Stat.RdTruesideNoiseFrontier.centered_error_failure_le_one_fortieth
% lean: CausalSmith.Stat.RdTruesideNoiseFrontier.polyInterval_coverage_of_average_failure_bounds

\item \textbf{Honesty after public selection.}
By \cref{obj:def:public-selection}, \(L_\star\) is a deterministic function of the public parameters. Consider separately its zero and positive values.

If \(L_\star=0\), the candidate specified in \cref{obj:synth_6} is
\[
\widehat\theta_0=0,
\qquad
\rho_0=\frac12,
\qquad
I_0=[-1/2,1/2].
\]
The target bounds above give
\[
\mathbb P_{P',n}\{\theta(P')\in I_0\}=1
\qquad\text{for every }P'\in\mathcal M_\beta(\sigma).
\]
If \(L_\star\ge1\), the fixed-index coverage result gives
\[
\mathbb P_{P',n}\{\theta(P')\in I_{L_\star}\}
\ge\frac9{10}
\qquad\text{for every }P'\in\mathcal M_\beta(\sigma).
\]

For positive indices, projection of both arm estimates gives
\[
-\frac12\le\widehat\theta_J\le\frac12,
\qquad
\rho_J\ge0.
\]
Thus the centered interval and \([-1,1]\) both contain
\(\widehat\theta_J\); their intersection is a nonempty connected closed interval contained in \([-1,1]\). Its endpoints are measurable, being formed from polynomial averages, a positive denominator floor, projection, and finite maxima and minima. The zero-index interval has these properties as well. Since selection is deterministic, the selected interval has these same properties. Together with the uniform coverage bound, this proves
\[
I_{L_\star}\in\mathcal I_{\beta,n,\sigma}
\]
by \cref{obj:def:honest-class}.
% lean: CausalSmith.Stat.RdTruesideNoiseFrontier.polyInterval_properties
% lean: CausalSmith.Stat.RdTruesideNoiseFrontier.honestAttainer_of_positive_coverage

\item \textbf{Selected risk.}
If \(L_\star=0\), then for every admissible \(P'\),
\[
|\widehat\theta_0-\theta(P')|
=|\theta(P')|\le\frac12=\rho_0,
\qquad
\mathbb E_{P',n}|\widehat\theta_0-\theta(P')|\le\rho_0.
\]
If \(L_\star\ge1\), nonnegativity of \(V_{L_\star}(\sigma)/n\) and monotonicity of the square root give
\[
\begin{aligned}
\mathbb E_{P',n}|\widehat\theta_{L_\star}-\theta(P')|
&\le12M_{L_\star,\beta}
+28\sqrt{\frac{V_{L_\star}(\sigma)}n}\\
&\le12M_{L_\star,\beta}
+28\sqrt{\frac{40V_{L_\star}(\sigma)}n}
=\rho_{L_\star}.
\end{aligned}
\]
In both cases the bound is uniform over \(P'\). Since
\(E_\beta(n,\sigma)=\rho_{L_\star}\) by
\cref{obj:def:public-selection}, taking the supremum gives
\[
\sup_{P'\in\mathcal M_\beta(\sigma)}
\mathbb E_{P',n}|\widehat\theta_{L_\star}-\theta(P')|
\le E_\beta(n,\sigma).
\]
% lean: CausalSmith.Stat.RdTruesideNoiseFrontier.attainerRisk_of_positive_risk

\item \textbf{Selected expected length.}
For every positive candidate index \(J\), the endpoint formula gives samplewise
\[
\begin{aligned}
\operatorname{length}(I_J)
&=\max\!\left\{0,\,
\min\{1,\widehat\theta_J+\rho_J\}
-\max\{-1,\widehat\theta_J-\rho_J\}\right\}\\
&\le2\rho_J.
\end{aligned}
\]
Indeed, the upper endpoint is at most
\(\widehat\theta_J+\rho_J\), the lower endpoint is at least
\(\widehat\theta_J-\rho_J\), and \(2\rho_J\ge0\).
For the zero candidate,
\[
\operatorname{length}(I_0)=1=2\rho_0.
\]
Consequently the bound holds at the selected index for every sample. Integrating under the probability experiment and then taking the supremum yields
\[
\sup_{P'\in\mathcal M_\beta(\sigma)}
\mathbb E_{P',n}\operatorname{length}(I_{L_\star})
\le2\rho_{L_\star}
=2E_\beta(n,\sigma).
\]
% lean: CausalSmith.Stat.RdTruesideNoiseFrontier.attainerLength_le_certificate
\end{enumerate}
\end{proof}

\begin{proof}[Proof of \cref{obj:thm:uniform-frontier}]
Fix \(0<\beta\le1\). We use the Legendre identities of
\cref{obj:lem:classical-legendre-facts} and the Hermite generating function and orthogonality of
\cref{obj:lem:classical-hermite-facts}. All constants introduced below depend at most on \(\beta\), unless explicitly indexed by an additional fixed parameter.

\begin{enumerate}
\item \textbf{Root facts used in the comparisons.}
For \(n\ge2\) and \(0\le\sigma\le1\), use the notation of
\cref{obj:def:uniform-rate}:
\[
\nu=2\beta+1,\qquad h_0=n^{-1/\nu}.
\]
Define, on \(0<h\le1\),
\[
\gamma_{\mathrm{bal}}(h)
=\log n+\nu\log h-\Psi(h,\sigma).
\]
We have
\[
0<h_0\le1,\qquad
\log n+\nu\log h_0=0.
\]
By \cref{obj:lem:noise-cost-scaling}, the function
\(\gamma_{\mathrm{bal}}\) is continuous. For \(0<h<t\le1\),
\[
\gamma_{\mathrm{bal}}(t)-\gamma_{\mathrm{bal}}(h)
=\nu\log(t/h)+\Psi(h,\sigma)-\Psi(t,\sigma)>0.
\]
Its endpoint values satisfy
\[
\gamma_{\mathrm{bal}}(h_0)=-\Psi(h_0,\sigma)\le0,
\qquad
\gamma_{\mathrm{bal}}(1)=\log n>0.
\]
The intermediate value theorem therefore gives a unique zero
\(\xi_{\mathrm{root}}\in[h_0,1]\). Strict increase also gives
\[
\{h\in[h_0,1]:\gamma_{\mathrm{bal}}(h)\le0\}
=[h_0,\xi_{\mathrm{root}}].
\]
Consequently, the supremum of this sublevel set equals
\(\xi_{\mathrm{root}}\), and the resolution in
\cref{obj:def:uniform-rate} satisfies
\[
h_\star(n,\sigma)=\xi_{\mathrm{root}},
\qquad
\log n+\nu\log h_\star(n,\sigma)
=\Psi(h_\star(n,\sigma),\sigma).
\]
Exponentiating yields the balance identity
\[
\exp\{\Psi(h_\star(n,\sigma),\sigma)\}
=n\,h_\star(n,\sigma)^\nu.
\]

If \(\sigma\le h_0\), then \(\Psi(h_0,\sigma)=0\), so uniqueness gives
\[
h_\star(n,\sigma)=h_0,\qquad
r_\beta(n,\sigma)=n^{-\beta/\nu}.
\]
This includes \(\sigma=0\).

If \(\sigma>h_0\), then \(\Psi(h_0,\sigma)>0\). Indeed, when
\(\sigma^4\le h_0\),
\[
(\sigma/h_0)^{2/3}>1.
\]
When \(h_0<\sigma^4\), we have \(h_0<1\) and
\(\sqrt{h_0}<\sigma^2\), whence
\[
h_0^{-1/2}>1,\qquad
1+\log(\sigma^2/\sqrt{h_0})\ge1,
\]
and their product exceeds one. Thus
\[
\gamma_{\mathrm{bal}}(h_0)<0,
\qquad
\gamma_{\mathrm{bal}}(\sigma)
=\log n+\nu\log\sigma
>\log n+\nu\log h_0=0.
\]
Strict increase now gives
\[
h_0<h_\star(n,\sigma)<\sigma.
\]
% lean: CausalSmith.Stat.RdTruesideNoiseFrontier.unique_rate_root
% lean: CausalSmith.Stat.RdTruesideNoiseFrontier.rateRoot_exp_balance
% lean: CausalSmith.Stat.RdTruesideNoiseFrontier.rateResolution_eq_direct_of_le
% lean: CausalSmith.Stat.RdTruesideNoiseFrontier.noisy_rateResolution_interior

\item \textbf{Lower comparisons.}
By \cref{obj:thm:direct-reduction}, there is \(c_{\mathrm{dir}}>0\) such that
\[
R_\beta(n,\sigma)\ge c_{\mathrm{dir}}n^{-\beta/\nu},
\qquad
\mathcal L_\beta(n,\sigma)\ge c_{\mathrm{dir}}n^{-\beta/\nu}
\]
for every admissible \(n,\sigma\). These are the required lower comparisons when
\(\sigma\le h_0\).

Consider \(\sigma>h_0\). For each integer \(m\ge2\), the scales of
\cref{obj:def:noise-support} satisfy
\[
b_m(\sigma)=\min\{1,\sigma\sqrt m/8\},
\qquad
h_m(\sigma)=b_m(\sigma)/m^2,
\qquad
0<b_m(\sigma)\le1.
\]
The support scale is nondecreasing in \(m\). Since
\(m^2\le4(m-1)^2\) for \(m\ge2\), it follows that
\[
h_{m-1}(\sigma)
=\frac{b_{m-1}(\sigma)}{(m-1)^2}
\le\frac{b_m(\sigma)}{(m-1)^2}
\le4h_m(\sigma).
\]
Also \(h_m(\sigma)\le m^{-2}\). In particular, any integer
\[
m>\sqrt{\frac{4096}{h_\star(n,\sigma)}}+2
\]
satisfies \(m\ge2\) and \(h_m(\sigma)\le h_\star(n,\sigma)/4096\).
Choose the smallest integer \(m\ge2\) with this latter property. At index two,
\[
b_2(\sigma)=\frac{\sigma\sqrt2}{8},
\qquad
h_2(\sigma)=\frac{\sigma\sqrt2}{32}
>\frac{h_\star(n,\sigma)}{4096},
\]
because \(h_\star(n,\sigma)<\sigma\) and
\(1\le\sqrt2<2\). Hence \(m\ge3\), and minimality gives
\[
h_{m-1}(\sigma)>\frac{h_\star(n,\sigma)}{4096}.
\]
Combining this with the predecessor bound gives
\[
\frac{h_\star(n,\sigma)}{16384}
<h_m(\sigma)
\le\frac{h_\star(n,\sigma)}{4096}.
\]

For this chosen index, define
\[
h_{\mathrm{low}}=h_m(\sigma),\qquad
b_{\mathrm{low}}=b_m(\sigma),\qquad
F_{\mathrm{cost}}(h)=1+\Psi(h,\sigma)\quad(0<h\le1).
\]
Thus \(0<h_{\mathrm{low}}\le h_\star(n,\sigma)<\sigma\le1\).
The ratio inequality of \cref{obj:lem:noise-cost-scaling} gives
\[
\frac{F_{\mathrm{cost}}(h_{\mathrm{low}})}
     {F_{\mathrm{cost}}(h_\star(n,\sigma))}
\ge
\left(\frac{h_\star(n,\sigma)}{h_{\mathrm{low}}}\right)^{1/2}
\ge64.
\]
Using \(h_{\mathrm{low}}^\nu\le h_\star(n,\sigma)^\nu\), the balance identity, and
\(\Psi(h_\star(n,\sigma),\sigma)\ge0\), we obtain
\[
\begin{aligned}
n h_{\mathrm{low}}^\nu
 \exp\{-F_{\mathrm{cost}}(h_{\mathrm{low}})/32\}
&\le
n h_\star(n,\sigma)^\nu
 \exp\{-2F_{\mathrm{cost}}(h_\star(n,\sigma))\}\\
&=\exp\{-2-\Psi(h_\star(n,\sigma),\sigma)\}\\
&\le e^{-2}.
\end{aligned}
\]

Define the observed alternative laws and their information budget by
\[
\mathsf P_\pm=(P^\pm_{b_{\mathrm{low}},m})_{\mathrm{obs}},
\qquad
\Delta_{\mathrm{info}}
=\frac{32}{3}\kappa^2h_{\mathrm{low}}^\nu
 \exp\{-F_{\mathrm{cost}}(h_{\mathrm{low}})/32\},
\qquad
\kappa=\frac1{100}.
\]
By \cref{obj:lem:cancellation-information-envelope},
\[
\chi^2(\mathsf P_+,\mathsf P_-)\le\Delta_{\mathrm{info}}<\infty,
\]
and the preceding calculation gives
\[
0\le n\Delta_{\mathrm{info}}
\le\frac{32}{3}\kappa^2e^{-2}
\le\frac{32}{30000}
\le\frac1{100}.
\]
Finiteness of this divergence supplies absolute continuity and an integrable
squared likelihood-ratio deviation. To spell out the product comparison, put
\[
\chi_{\mathrm{obs}}=\chi^2(\mathsf P_+,\mathsf P_-),
\qquad
\chi_{\mathrm{sample}}
=\chi^2(\mathsf P_+^{\otimes n},\mathsf P_-^{\otimes n}).
\]
The product likelihood identity and \(1+\chi_{\mathrm{obs}}\le e^{\chi_{\mathrm{obs}}}\) give
\[
1+\chi_{\mathrm{sample}}
=(1+\chi_{\mathrm{obs}})^n
\le e^{n\chi_{\mathrm{obs}}}
\le e^{1/100}
\le\frac1{1-1/100}
=\frac{100}{99}.
\]
Therefore
\[
\chi_{\mathrm{sample}}\le\frac1{99}\le\frac4{25},
\qquad
\operatorname{TV}(\mathsf P_+^{\otimes n},\mathsf P_-^{\otimes n})
\le\frac12\sqrt{\chi_{\mathrm{sample}}}
\le\frac15.
\]
Here the total-variation inequality follows by applying Cauchy--Schwarz to
the likelihood-ratio deviation; the positive and negative parts of that
deviation have equal mass.

The support and order conditions required by
\cref{obj:lem:two-point-transfer} hold. That result, which includes independent
procedure randomization, yields
\[
R_\beta(n,\sigma)\ge\frac45\kappa h_{\mathrm{low}}^\beta,
\qquad
\mathcal L_\beta(n,\sigma)\ge\frac65\kappa h_{\mathrm{low}}^\beta.
\]
Consequently, with
\[
c_{\mathrm{noise}}=\frac45\kappa\,16384^{-\beta},
\qquad
c_{\mathrm{mat}}=\min\{c_{\mathrm{dir}},c_{\mathrm{noise}}\}>0,
\]
the direct and noisy cases together give
\[
c_{\mathrm{mat}}r_\beta(n,\sigma)\le R_\beta(n,\sigma),
\qquad
c_{\mathrm{mat}}r_\beta(n,\sigma)\le\mathcal L_\beta(n,\sigma)
\]
uniformly over \(n\ge2\) and \(0\le\sigma\le1\).
% lean: CausalSmith.Stat.RdTruesideNoiseFrontier.exists_noiseResolution_bracket
% lean: CausalSmith.Stat.RdTruesideNoiseFrontier.rounded_noise_information_factor
% lean: CausalSmith.Stat.RdTruesideNoiseFrontier.sampleLaw_tv_le_one_fifth_of_chiSqExt
% lean: CausalSmith.Stat.RdTruesideNoiseFrontier.frontierRate_lower_bounds

\item \textbf{Upper comparison for the public radius.}
Choose \(K_{\mathrm{bias}}>0\) and \(D_{\mathrm{env}}>0\) from
\cref{obj:lem:positive-endpoint,obj:lem:variance-cost-envelope}, so that
\[
M_{L,\beta}\le K_{\mathrm{bias}}L^{-2\beta},
\qquad
V_L(\sigma)\le D_{\mathrm{env}}L^2
 \exp\{D_{\mathrm{env}}[1+\Psi(L^{-2},\sigma)]\}
\]
for every \(L\ge1\). Define
\[
\begin{aligned}
A_{\mathrm{dil}}&=4(D_{\mathrm{env}}+1)^2,\\
Q_{\mathrm{var}}&=\sqrt{160D_{\mathrm{env}}e^{D_{\mathrm{env}}+1}},\\
C_{\mathrm{up}}&=\max\left\{
 \frac12A_{\mathrm{dil}}^\beta,\,
 12K_{\mathrm{bias}}A_{\mathrm{dil}}^\beta+28Q_{\mathrm{var}}
 \right\}.
\end{aligned}
\]
These constants are positive, and \(A_{\mathrm{dil}}\ge4\).

The radii in \cref{obj:synth_6} and the public minimization in
\cref{obj:def:public-selection} give
\[
\rho_0=\frac12,\qquad
\rho_L=12M_{L,\beta}+28\sqrt{\frac{40V_L(\sigma)}n}\quad(L\ge1),
\qquad
E_\beta(n,\sigma)\le\rho_L\quad(0\le L\le L_{\max}),
\]
where the finite cap is
\[
L_{\max}=\left\lceil n^{1/(4\beta+2)}\right\rceil.
\]

If \(1\le A_{\mathrm{dil}}h_\star(n,\sigma)\), then
\[
E_\beta(n,\sigma)\le\frac12
\le\frac12A_{\mathrm{dil}}^\beta h_\star(n,\sigma)^\beta
\le C_{\mathrm{up}}r_\beta(n,\sigma).
\]
Suppose instead that \(A_{\mathrm{dil}}h_\star(n,\sigma)<1\), and define
\[
x_{\mathrm{up}}
=(A_{\mathrm{dil}}h_\star(n,\sigma))^{-1/2},
\qquad
L_{\mathrm{up}}=\lceil x_{\mathrm{up}}\rceil,
\qquad
h_{\mathrm{up}}=L_{\mathrm{up}}^{-2}.
\]
Since \(x_{\mathrm{up}}>1\), ceiling bounds give
\[
1\le L_{\mathrm{up}},
\qquad
x_{\mathrm{up}}\le L_{\mathrm{up}}
<x_{\mathrm{up}}+1\le2x_{\mathrm{up}}.
\]
Taking inverse squares yields
\[
\frac{A_{\mathrm{dil}}h_\star(n,\sigma)}4
\le h_{\mathrm{up}}
\le A_{\mathrm{dil}}h_\star(n,\sigma)<1.
\]
Moreover,
\[
x_{\mathrm{up}}
\le h_\star(n,\sigma)^{-1/2}
\le h_0^{-1/2}
=n^{1/(4\beta+2)},
\]
so monotonicity of the ceiling gives \(L_{\mathrm{up}}\le L_{\max}\).

Define the dilation multiplier by
\[
\alpha_{\mathrm{dil}}=\frac{h_{\mathrm{up}}}{h_\star(n,\sigma)}.
\]
The width bounds imply
\[
(D_{\mathrm{env}}+1)^2\le\alpha_{\mathrm{dil}}\le A_{\mathrm{dil}},
\qquad
\alpha_{\mathrm{dil}}\ge1,
\qquad
\alpha_{\mathrm{dil}}h_\star(n,\sigma)=h_{\mathrm{up}}\le1.
\]
Thus \cref{obj:lem:noise-cost-scaling} applies and gives
\[
1+\Psi(h_{\mathrm{up}},\sigma)
\le\max\left\{
1,\,
\alpha_{\mathrm{dil}}^{-1/2}
[1+\Psi(h_\star(n,\sigma),\sigma)]
\right\}.
\]
Because \(D_{\mathrm{env}}/\sqrt{\alpha_{\mathrm{dil}}}\le1\), multiplication by
\(D_{\mathrm{env}}\) gives
\[
\begin{aligned}
D_{\mathrm{env}}[1+\Psi(h_{\mathrm{up}},\sigma)]
&\le
\max\left\{
D_{\mathrm{env}},\,
1+\Psi(h_\star(n,\sigma),\sigma)
\right\}\\
&\le D_{\mathrm{env}}+1+\Psi(h_\star(n,\sigma),\sigma).
\end{aligned}
\]
Hence
\[
V_{L_{\mathrm{up}}}(\sigma)
\le D_{\mathrm{env}}L_{\mathrm{up}}^2
 e^{D_{\mathrm{env}}+1}
 \exp\{\Psi(h_\star(n,\sigma),\sigma)\}.
\]
The rounding bound also gives
\[
L_{\mathrm{up}}^2
\le\frac4{A_{\mathrm{dil}}h_\star(n,\sigma)}
\le\frac4{h_\star(n,\sigma)}.
\]
Using the root balance and \(\nu=2\beta+1\), we obtain
\[
\frac{40V_{L_{\mathrm{up}}}(\sigma)}n
\le160D_{\mathrm{env}}e^{D_{\mathrm{env}}+1}
 h_\star(n,\sigma)^{2\beta},
\]
and therefore
\[
\sqrt{\frac{40V_{L_{\mathrm{up}}}(\sigma)}n}
\le Q_{\mathrm{var}}h_\star(n,\sigma)^\beta.
\]
The bias bound gives
\[
M_{L_{\mathrm{up}},\beta}
\le K_{\mathrm{bias}}h_{\mathrm{up}}^\beta
\le K_{\mathrm{bias}}A_{\mathrm{dil}}^\beta
 h_\star(n,\sigma)^\beta.
\]
Since \(L_{\mathrm{up}}\) is an eligible candidate,
\[
\begin{aligned}
E_\beta(n,\sigma)
&\le\rho_{L_{\mathrm{up}}}\\
&\le
(12K_{\mathrm{bias}}A_{\mathrm{dil}}^\beta+28Q_{\mathrm{var}})
 h_\star(n,\sigma)^\beta\\
&\le C_{\mathrm{up}}r_\beta(n,\sigma).
\end{aligned}
\]
Both cases establish this comparison over the full parameter range.
% lean: CausalSmith.Stat.RdTruesideNoiseFrontier.upperDegree_width_bounds
% lean: CausalSmith.Stat.RdTruesideNoiseFrontier.upperDegree_le_Lmax
% lean: CausalSmith.Stat.RdTruesideNoiseFrontier.dilation_cost_absorption
% lean: CausalSmith.Stat.RdTruesideNoiseFrontier.certificate_le_frontierRate

\item \textbf{Matched estimation and interval chains.}
By \cref{obj:thm:finite-certificate}, the selected procedures satisfy
\[
\begin{aligned}
\sup_{P\in\mathcal M_\beta(\sigma)}
\mathbb E_{P,n}
|\widetilde\theta_{\beta,n,\sigma}-\theta(P)|
&\le E_\beta(n,\sigma),\\
\sup_{P\in\mathcal M_\beta(\sigma)}
\mathbb E_{P,n}\operatorname{length}(\widetilde I_{\beta,n,\sigma})
&\le2E_\beta(n,\sigma),\\
\widetilde I_{\beta,n,\sigma}&\in\mathcal I_{\beta,n,\sigma}.
\end{aligned}
\]
The minimax infima are bounded above by the objectives of these eligible procedures.

For completeness, the mean bounds imply
\[
|\theta(P)|=|\mu_1(P,0)-\mu_0(P,0)|\le\frac12.
\]
Thus the constant estimator zero has worst expected absolute error at most
\(1/2\), and the constant interval \([-1,1]\) has coverage one and length two.
Both extended minimax objectives are consequently finite. Conversion to real
numbers preserves their comparisons with the selected procedures.

Put
\[
C_{\mathrm{mat}}=2C_{\mathrm{up}}>0.
\]
Combining the lower comparisons, the finite certificate, and the upper comparison gives
\[
\begin{aligned}
c_{\mathrm{mat}}r_\beta(n,\sigma)
&\le R_\beta(n,\sigma)
\le\sup_{P\in\mathcal M_\beta(\sigma)}
 \mathbb E_{P,n}|\widetilde\theta_{\beta,n,\sigma}-\theta(P)|\\
&\le E_\beta(n,\sigma)
\le C_{\mathrm{mat}}r_\beta(n,\sigma),\\
c_{\mathrm{mat}}r_\beta(n,\sigma)
&\le\mathcal L_\beta(n,\sigma)
\le\sup_{P\in\mathcal M_\beta(\sigma)}
 \mathbb E_{P,n}\operatorname{length}(\widetilde I_{\beta,n,\sigma})\\
&\le2E_\beta(n,\sigma)
\le C_{\mathrm{mat}}r_\beta(n,\sigma).
\end{aligned}
\]
The lower transfer ranges over the stated randomized decision classes, with
arbitrary procedure degree.
% lean: CausalSmith.Stat.RdTruesideNoiseFrontier.matchedFrontier_of_scalar_bounds
% lean: CausalSmith.Stat.RdTruesideNoiseFrontier.matched_frontier

\item \textbf{Finite-sample branches and interfaces.}
Define the branch comparison constants by
\[
\begin{gathered}
A_{\mathrm{int}}=(1+3\nu/2)^{3/2},\qquad
c_\tau=\frac1{2(1+2\nu)},\qquad C_\tau=e+6,\\
c_{\mathrm{pow}}=C_\tau^{-2\beta},\qquad
C_{\mathrm{pow}}=c_\tau^{-2\beta},\\
c_{\mathrm{br}}=\min\{1,c_\tau,c_{\mathrm{pow}}\},
\qquad
C_{\mathrm{br}}=\max\{A_{\mathrm{int}},C_\tau,C_{\mathrm{pow}}\}.
\end{gathered}
\]
All are positive.

The direct branch was established above. In the noisy region,
\(0<\sigma\le1\) and \(h_0<h_\star(n,\sigma)<\sigma\). Evaluation at the interface gives
\[
\Psi(\sigma^4,\sigma)=\sigma^{-2}-1,
\qquad
\gamma_{\mathrm{bal}}(\sigma^4)
=\log(n\sigma^{4\nu})-(\sigma^{-2}-1).
\]
This identity also holds at \(\sigma=1\), where both quantities on the right
of the noise-cost identity are zero. Strict increase on \((0,1]\) gives
\[
\sigma^4\le h_\star(n,\sigma)
\quad\Longleftrightarrow\quad
\log(n\sigma^{4\nu})\le\sigma^{-2}-1.
\]
The comparison is valid throughout the noisy region, including when
\(\sigma^4<h_0\).

On this intermediate branch, define
\[
z=\left(\frac{\sigma}{h_\star(n,\sigma)}\right)^{2/3},
\qquad
S_{\mathrm{int}}=1+\log(n\sigma^\nu).
\]
The inequalities \(h_\star(n,\sigma)<\sigma\) and
\(\sigma^4\le h_\star(n,\sigma)\) imply
\[
1<z\le\sigma^{-2}.
\]
The intermediate noise-cost formula and the root identity give
\[
1+\Psi(h_\star(n,\sigma),\sigma)=z,
\qquad
\log z=\frac23\{\log\sigma-\log h_\star(n,\sigma)\}.
\]
Substitution yields
\[
z+\frac{3\nu}{2}\log z=S_{\mathrm{int}},
\qquad
h_\star(n,\sigma)=\sigma z^{-3/2}.
\]
The left side of the scalar equation is strictly increasing for \(z>0\),
so the indicated root is unique. Since \(z>1\),
\[
0\le\log z\le z,
\qquad
z\le S_{\mathrm{int}}\le(1+3\nu/2)z.
\]
Taking inverse powers gives
\[
\frac{\sigma}{S_{\mathrm{int}}^{3/2}}
\le h_\star(n,\sigma)
\le A_{\mathrm{int}}\frac{\sigma}{S_{\mathrm{int}}^{3/2}}.
\]

If the branch-test inequality is strict in the other direction, then
\(h_\star(n,\sigma)<\sigma^4\). Define
\[
\tau=h_\star(n,\sigma)^{-1/2}.
\]
We have
\[
\tau>\sigma^{-2}\ge1,\qquad
h_\star(n,\sigma)=\tau^{-2},
\]
and the compact cost formula gives
\[
1+\Psi(h_\star(n,\sigma),\sigma)
=\tau\{1+\log(\sigma^2\tau)\}.
\]
Since \(\log\tau=-\tfrac12\log h_\star(n,\sigma)\), the root equation becomes
\[
2\nu\log\tau+\tau\{1+\log(\sigma^2\tau)\}
=1+\log n.
\]
For \(\sigma^{-2}<\tau_1<\tau_2\), both
\(\log\tau\) and \(1+\log(\sigma^2\tau)\) increase strictly, and the latter
factor is positive. In particular,
\[
\begin{aligned}
\tau_1\{1+\log(\sigma^2\tau_1)\}
&<\tau_2\{1+\log(\sigma^2\tau_1)\}\\
&<\tau_2\{1+\log(\sigma^2\tau_2)\}.
\end{aligned}
\]
Thus the compact equation has a unique root on the indicated domain.

To compare that root with an explicit expression, define
\[
\Lambda_n=1+\log n,\qquad
\Omega_{n,\sigma}=\log(e+\sigma^2\Lambda_n),\qquad
y_{\mathrm{cmp}}=\sigma^2\tau,\qquad
q_{\mathrm{cmp}}=\frac{\Lambda_n}{\tau},\qquad
J_{\mathrm{cmp}}=1+\log y_{\mathrm{cmp}}.
\]
Here \(\Lambda_n>1\), \(\Omega_{n,\sigma}>0\),
\(y_{\mathrm{cmp}}>1\), and \(J_{\mathrm{cmp}}>1\). Dividing the compact
equation by \(\tau\) gives
\[
q_{\mathrm{cmp}}
=2\nu\frac{\log\tau}{\tau}+J_{\mathrm{cmp}}.
\]
Using \(0\le\log\tau\le\tau\) and \(\nu\le3\), we obtain
\[
J_{\mathrm{cmp}}\le q_{\mathrm{cmp}}
\le J_{\mathrm{cmp}}+2\nu
\le J_{\mathrm{cmp}}+6,
\qquad
q_{\mathrm{cmp}}\le(1+2\nu)J_{\mathrm{cmp}}.
\]
Also \(\sigma^2\Lambda_n=y_{\mathrm{cmp}}q_{\mathrm{cmp}}\), so
\[
\Omega_{n,\sigma}\ge1,\qquad
\Omega_{n,\sigma}\ge\log y_{\mathrm{cmp}},
\qquad
J_{\mathrm{cmp}}\le2\Omega_{n,\sigma}.
\]
For the reverse comparison,
\[
e+\sigma^2\Lambda_n
=e+y_{\mathrm{cmp}}q_{\mathrm{cmp}}
\le y_{\mathrm{cmp}}(e+J_{\mathrm{cmp}}+6).
\]
The elementary bound \(\log t\le t-1\) for \(t>0\) therefore gives
\[
\begin{aligned}
\Omega_{n,\sigma}
&\le\log y_{\mathrm{cmp}}+\log(e+J_{\mathrm{cmp}}+6)\\
&\le(J_{\mathrm{cmp}}-1)+(e+J_{\mathrm{cmp}}+5)\\
&=2J_{\mathrm{cmp}}+e+4\\
&\le(e+6)J_{\mathrm{cmp}}.
\end{aligned}
\]
It follows that
\[
q_{\mathrm{cmp}}\le2(1+2\nu)\Omega_{n,\sigma},
\qquad
\Omega_{n,\sigma}\le(e+6)q_{\mathrm{cmp}}.
\]
All quantities divided below are positive, and these inequalities yield
\[
c_\tau\frac{\Lambda_n}{\Omega_{n,\sigma}}
\le\tau
\le C_\tau\frac{\Lambda_n}{\Omega_{n,\sigma}}.
\]
Raising this comparison to the negative power \(-2\beta\) reverses the
inequalities. Since
\[
r_\beta(n,\sigma)=\tau^{-2\beta},
\qquad
\left(\frac{\Lambda_n}{\Omega_{n,\sigma}}\right)^{-2\beta}
=\left(\frac{\Omega_{n,\sigma}}{\Lambda_n}\right)^{2\beta},
\]
we obtain
\[
c_{\mathrm{pow}}
\left(\frac{\Omega_{n,\sigma}}{\Lambda_n}\right)^{2\beta}
\le r_\beta(n,\sigma)
\le C_{\mathrm{pow}}
\left(\frac{\Omega_{n,\sigma}}{\Lambda_n}\right)^{2\beta}.
\]
These are the asserted compact comparisons because
\(\Lambda_n=\log(en)\).

At the intermediate--compact equality,
\[
\gamma_{\mathrm{bal}}(\sigma^4)=0.
\]
If \(\sigma^4<h_0\), strict increase would give
\(0=\gamma_{\mathrm{bal}}(\sigma^4)<\gamma_{\mathrm{bal}}(h_0)\le0\).
Hence \(\sigma^4\in[h_0,1]\), and root uniqueness implies
\[
h_\star(n,\sigma)=\sigma^4.
\]
The equality condition can be written as
\[
\log n+4\nu\log\sigma=\sigma^{-2}-1.
\]
Using \(\log(\sigma^{-2})=-2\log\sigma\) and
\(\sigma^2\sigma^{-2}=1\), it gives both scalar identities
\[
\begin{aligned}
\sigma^{-2}+\frac{3\nu}{2}\log(\sigma^{-2})
&=1+\log(n\sigma^\nu),\\
2\nu\log(\sigma^{-2})
+\sigma^{-2}\{1+\log(\sigma^2\sigma^{-2})\}
&=1+\log n.
\end{aligned}
\]
Thus \(z=\tau=\sigma^{-2}\), and both resolution formulas equal \(\sigma^4\).

At the direct--noisy equality \(\sigma=h_0\), the direct result gives
\[
h_\star(n,\sigma)=\sigma=h_0,\qquad
\log(n\sigma^\nu)=0.
\]
The intermediate scalar equation then holds at \(z=1\), and
\(\sigma z^{-3/2}=\sigma\).

Finally, \(c_{\mathrm{br}}\le1,c_\tau,c_{\mathrm{pow}}\) and
\(C_{\mathrm{br}}\ge A_{\mathrm{int}},C_\tau,C_{\mathrm{pow}}\).
All comparison factors above are positive. Replacing their lower constants
by \(c_{\mathrm{br}}\) and their upper constants by \(C_{\mathrm{br}}\)
therefore gives simultaneous branch comparisons, with the exact equations
and interfaces unchanged.
% lean: CausalSmith.Stat.RdTruesideNoiseFrontier.intermediateCondition_iff
% lean: CausalSmith.Stat.RdTruesideNoiseFrontier.intermediate_branch_representation
% lean: CausalSmith.Stat.RdTruesideNoiseFrontier.intermediate_branch_bounds
% lean: CausalSmith.Stat.RdTruesideNoiseFrontier.compact_branch_representation
% lean: CausalSmith.Stat.RdTruesideNoiseFrontier.compactScalar_bounds
% lean: CausalSmith.Stat.RdTruesideNoiseFrontier.compactPower_bounds
% lean: CausalSmith.Stat.RdTruesideNoiseFrontier.intermediate_compact_interface
% lean: CausalSmith.Stat.RdTruesideNoiseFrontier.direct_interface
% lean: CausalSmith.Stat.RdTruesideNoiseFrontier.frontier_branches

\item \textbf{Shrinking-noise sequences.}
Fix \(0<a<1/\nu\) and set \(\sigma=n^{-a}\). Define
\[
\eta_a=1-4a\nu,\qquad
\epsilon_a=\frac1{2(|\eta_a|+1)}.
\]
Since \(2a>0\), the relation \(\log n=o(n^{2a})\) and divergence of \(n^{2a}\)
give an integer \(N\ge2\) such that, for all \(n\ge N\),
\[
\log n\le\epsilon_a n^{2a},
\qquad n^{2a}\ge2.
\]
Consequently,
\[
\eta_a\log n
\le|\eta_a|\epsilon_a n^{2a}
\le\frac12n^{2a}
\le n^{2a}-1.
\]
The relevant quantities are exactly
\[
\log(n\sigma^{4\nu})=(1-4a\nu)\log n,
\qquad
\sigma^{-2}=n^{2a}.
\]
Thus the intermediate branch holds for \(n\ge N\); moreover,
\(0<\sigma\le1\) and \(\sigma>h_0\) for every \(n\ge2\).

Define
\[
d_{\mathrm{shrink}}=1-a\nu>0,\qquad
K_{\mathrm{shrink}}=d_{\mathrm{shrink}}+\frac1{\log2}>0.
\]
For \(n\ge2\),
\[
1+\log(n\sigma^\nu)=1+d_{\mathrm{shrink}}\log n,
\]
and \(\log n\ge\log2\) implies
\[
d_{\mathrm{shrink}}\log n
\le1+\log(n\sigma^\nu)
\le K_{\mathrm{shrink}}\log n.
\]
Applying the intermediate comparison with \(c_{\mathrm{br}},C_{\mathrm{br}}\)
therefore gives, for \(n\ge N\),
\[
\frac{c_{\mathrm{br}}}{K_{\mathrm{shrink}}^{3/2}}
\,n^{-a}(\log n)^{-3/2}
\le h_\star(n,n^{-a})
\le
\frac{C_{\mathrm{br}}}{d_{\mathrm{shrink}}^{3/2}}
\,n^{-a}(\log n)^{-3/2}.
\]
Taking the positive power \(\beta\) yields
\[
\begin{aligned}
\left(\frac{c_{\mathrm{br}}}{K_{\mathrm{shrink}}^{3/2}}\right)^\beta
n^{-a\beta}(\log n)^{-3\beta/2}
&\le r_\beta(n,n^{-a})\\
&\le
\left(\frac{C_{\mathrm{br}}}{d_{\mathrm{shrink}}^{3/2}}\right)^\beta
n^{-a\beta}(\log n)^{-3\beta/2}.
\end{aligned}
\]
Both constants are positive and may depend on \(a\).

For \(a\ge1/\nu\), monotonicity in the exponent for \(n>1\) gives
\[
0<n^{-a}\le n^{-1/\nu}=h_0\le1.
\]
Hence the exact direct branch applies for every \(n\ge2\), with
\[
h_\star(n,n^{-a})=h_0,
\qquad
r_\beta(n,n^{-a})=n^{-\beta/\nu}.
\]
% lean: CausalSmith.Stat.RdTruesideNoiseFrontier.eventually_shrinking_intermediate
% lean: CausalSmith.Stat.RdTruesideNoiseFrontier.shrinking_intermediate_rate_bounds
% lean: CausalSmith.Stat.RdTruesideNoiseFrontier.frontier_sequences

\item \textbf{Fixed positive noise.}
Fix \(0<\sigma\le1\), and define
\[
\begin{aligned}
\ell_{\mathrm{noise}}&=\log(\sigma^2),\\
M_{\mathrm{noise}}&=\max\left\{
1,\,-2\ell_{\mathrm{noise}},\,|\log(e+\sigma^2)|
\right\},\\
K_{\mathrm{noise}}&=1+|\log(e+\sigma^2)|,\\
T_{\mathrm{noise}}&=\sigma^{-2}-1-4\nu\log\sigma.
\end{aligned}
\]
Since \(\log n\to\infty\) and \(h_0\to0\), there is \(N\ge2\) such that, for
\(n\ge N\),
\[
e^{M_{\mathrm{noise}}}\le\log n,\qquad
T_{\mathrm{noise}}+1\le\log n,\qquad
h_0<\sigma.
\]
The middle inequality gives
\[
\log(n\sigma^{4\nu})
=\log n+4\nu\log\sigma
>\sigma^{-2}-1,
\]
so the compact branch applies.

Retain
\[
\Lambda_n=1+\log n,\qquad
\Omega_{n,\sigma}=\log(e+\sigma^2\Lambda_n),
\]
and define
\[
L_{\mathrm{log}}=\log\Lambda_n.
\]
The first threshold implies
\[
L_{\mathrm{log}}\ge M_{\mathrm{noise}}\ge1,
\qquad
L_{\mathrm{log}}\ge-2\ell_{\mathrm{noise}}.
\]
Monotonicity of the logarithm gives
\[
\Omega_{n,\sigma}
\ge\log(\sigma^2\Lambda_n)
=\ell_{\mathrm{noise}}+L_{\mathrm{log}}
\ge\frac12L_{\mathrm{log}}.
\]
Since \(\Lambda_n\ge1\),
\[
e+\sigma^2\Lambda_n\le(e+\sigma^2)\Lambda_n,
\]
and hence
\[
\begin{aligned}
\Omega_{n,\sigma}
&\le\log(e+\sigma^2)+L_{\mathrm{log}}\\
&\le|\log(e+\sigma^2)|L_{\mathrm{log}}+L_{\mathrm{log}}\\
&=K_{\mathrm{noise}}L_{\mathrm{log}}.
\end{aligned}
\]
Applying the compact rate comparison and taking the power \(2\beta>0\) gives
\[
\begin{aligned}
c_{\mathrm{br}}2^{-2\beta}
\left(\frac{L_{\mathrm{log}}}{\Lambda_n}\right)^{2\beta}
&\le r_\beta(n,\sigma)\\
&\le
C_{\mathrm{br}}K_{\mathrm{noise}}^{2\beta}
\left(\frac{L_{\mathrm{log}}}{\Lambda_n}\right)^{2\beta}.
\end{aligned}
\]
Since \(\Lambda_n=\log(en)\) and \(L_{\mathrm{log}}=\log\log(en)\), this is the
fixed-noise assertion with the permitted dependence on \(\sigma\).
% lean: CausalSmith.Stat.RdTruesideNoiseFrontier.fixed_noise_rate_bounds
% lean: CausalSmith.Stat.RdTruesideNoiseFrontier.frontier_sequences

\item \textbf{The case \(\beta=1\), \(\sigma=n^{-1/6}\).}
For every \(n\ge2\),
\[
\nu=3,\qquad h_0=n^{-1/3},\qquad
\sigma=n^{-1/6}>h_0,\qquad 0<\sigma\le1.
\]
Power identities give
\[
\sigma^{12}=n^{-2},\qquad
n\sigma^{12}=n^{-1},\qquad
\sigma^{-2}=n^{1/3}.
\]
Thus
\[
\log(n\sigma^{12})=-\log n
\le n^{1/3}-1=\sigma^{-2}-1,
\]
so the intermediate branch holds at every such sample size. Also
\[
n\sigma^3=n^{1/2},
\qquad
\log(n\sigma^3)=\frac12\log n,
\qquad
r_1(n,\sigma)=h_\star(n,\sigma).
\]
The intermediate comparison at \(\beta=1\) therefore yields
\[
r_1(n,n^{-1/6})
=h_\star(n,n^{-1/6})
\asymp
\frac{n^{-1/6}}{[1+\tfrac12\log n]^{3/2}},
\]
with constants independent of \(n\).
% lean: CausalSmith.Stat.RdTruesideNoiseFrontier.beta_one_frontier

\item \textbf{A simultaneous choice of constants.}
Finally, define
\[
c=\min\{c_{\mathrm{mat}},c_{\mathrm{br}}\},
\qquad
C=\max\{C_{\mathrm{mat}},C_{\mathrm{br}}\}.
\]
Then
\[
0<c,\qquad 0<C,\qquad
c\le c_{\mathrm{mat}},c_{\mathrm{br}},
\qquad
C\ge C_{\mathrm{mat}},C_{\mathrm{br}}.
\]
The root bounds give
\[
h_\star(n,\sigma)\ge h_0>0,
\qquad
r_\beta(n,\sigma)=h_\star(n,\sigma)^\beta\ge0.
\]
Thus, in the matched chains,
\[
c\,r_\beta(n,\sigma)
\le c_{\mathrm{mat}}r_\beta(n,\sigma),
\qquad
C_{\mathrm{mat}}r_\beta(n,\sigma)
\le C\,r_\beta(n,\sigma).
\]
The branch comparisons admit the same replacement: their factors
\[
\frac{\sigma}{[1+\log(n\sigma^\nu)]^{3/2}},
\qquad
\frac{\Lambda_n}{\Omega_{n,\sigma}},
\qquad
\left(\frac{\Omega_{n,\sigma}}{\Lambda_n}\right)^{2\beta}
\]
are positive on their respective noisy branches. Multiplication therefore
preserves the inequalities when their lower constants are reduced to \(c\)
and their upper constants are increased to \(C\). The root, branch equations,
uniqueness assertions, and interfaces retain their exact values. Together
with the sequence comparisons and the case \(\beta=1\), this proves every
assertion simultaneously.
% lean: CausalSmith.Stat.RdTruesideNoiseFrontier.matched_frontier_weaken
% lean: CausalSmith.Stat.RdTruesideNoiseFrontier.frontier_branches_weaken
% lean: CausalSmith.Stat.RdTruesideNoiseFrontier.uniform_frontier
\end{enumerate}
\end{proof}

\begin{proof}[Proof of \cref{obj:prop:published-class-converse-transfer}]
Fix \(\beta\in(0,1]\).

\begin{enumerate}
\item By \cref{obj:thm:uniform-frontier}, there is a constant \(c>0\), depending only on \(\beta\), such that, simultaneously for every \(n\ge2\) and \(\sigma\in[0,1]\),
\[
c\,r_\beta(n,\sigma)\le R_\beta(n,\sigma),
\qquad
c\,r_\beta(n,\sigma)\le\mathcal L_\beta(n,\sigma).
\]
We use this same constant for both transferred lower bounds.
% lean: CausalSmith.Stat.RdTruesideNoiseFrontier.uniform_frontier

\item Fix \(\sigma\in[0,1]\) and \(P\in\mathcal M_\beta(\sigma)\). Define the embedded law and its accompanying density and mean versions by
\[
Q=P,\qquad
(Z,V^0,V^1,E)=(X,Y^0,Y^1,\varepsilon),
\qquad
a(z)=f(P,z),\qquad v_d(z)=\mu_d(P,z)
\quad(d\in\{0,1\},\ z\in\mathbb R).
\]
The density identity is inherited from \cref{obj:synth_1}. For every measurable set \(F_{\mathrm{score}}\subseteq\mathbb R\), the conditional-mean identity becomes
\[
\begin{aligned}
\int_{\{Z\in F_{\mathrm{score}}\}}V^d\,dQ
&=\int_{F_{\mathrm{score}}\cap[-1,1]}
       \mu_d(P,z)f(P,z)\,dz\\
&=\int_{F_{\mathrm{score}}}\mu_d(P,z)f(P,z)\,dz\\
&=\int_{F_{\mathrm{score}}}v_d(z)a(z)\,dz.
\end{aligned}
\]
The first equality is supplied by \cref{obj:synth_1}; the second follows because \(f(P,z)=0\) outside \([-1,1]\). Thus the chosen \(v_d\) remain conditional-mean versions on the real score space.
% lean: CausalSmith.Stat.RdTruesideNoiseFrontier.benchmark_version_on_real

Choose
\[
r=\frac12.
\]
Since \((-r,r)\subseteq[-1,1]\), the full-interval continuity in \cref{obj:synth_1} gives continuity of \(a\) and both \(v_d\) on \((-r,r)\). Moreover, \cref{obj:ass:density-bounds,obj:ass:mean-bounds} give
\[
a(z)\ge\frac14>0\quad(z\in(-r,r)),
\qquad
0\le\frac14\le v_d(0)\le\frac34\le1
\quad(d\in\{0,1\}).
\]
The Gaussian condition follows from \cref{obj:ass:gaussian}:
\[
E\sim N(0,1).
\]
To verify the required joint independence, define the measurable map
\[
\Phi_{\mathrm{obs}}(z,y_0,y_1)
=
\left(
z,\,
\mathbf1\{z\ge0\}y_1+
\bigl(1-\mathbf1\{z\ge0\}\bigr)y_0
\right),
\qquad
(z,y_0,y_1)\in\mathbb R\times\{0,1\}^2.
\]
Then
\[
(Z,B)=\Phi_{\mathrm{obs}}(X,Y^0,Y^1).
\]
By \cref{obj:ass:error-independence}, independence is preserved under this measurable transformation, so
\[
E\perp(Z,B).
\]
All defining conditions of \(\mathcal G(\sigma)\) therefore hold.
% lean: CausalSmith.Stat.RdTruesideNoiseFrontier.benchmarkEmbedding_mem

The embedding leaves the coordinates used by the observation map unchanged. Consequently,
\[
\vartheta(Q)=v_1(0)-v_0(0)=\theta(P),
\qquad
Q\circ(R,A,B)^{-1}=P_{\mathrm{obs}},
\]
and hence, for every integer \(n\ge0\),
\[
\bigl(Q\circ(R,A,B)^{-1}\bigr)^{\otimes n}
\otimes\operatorname{Unif}[0,1]
=
P_{\mathrm{obs}}^{\otimes n}
\otimes\operatorname{Unif}[0,1]
=
\mathbb P_{P,n}.
\]
At \(n=0\), the sample factor is the empty product probability law, so the identity also covers that case.
% lean: CausalSmith.Stat.RdTruesideNoiseFrontier.benchmarkEmbedding_preserves

\item Fix \(n\ge2\) and \(\sigma\in[0,1]\). For every measurable estimate \(T\), target and experiment preservation give
\[
\sup_{P\in\mathcal M_\beta(\sigma)}
\mathbb E_{P,n}|T-\theta(P)|
\le
\sup_{Q\in\mathcal G(\sigma)}
\mathbb E_Q|T-\vartheta(Q)|.
\]
Taking infima over the common estimator class gives the corresponding comparison of extended nonnegative minimax objectives.

To justify the comparison after conversion to real values, note that every \(Q\in\mathcal G(\sigma)\) satisfies
\[
-1\le v_1(0)-v_0(0)=\vartheta(Q)\le1.
\]
Define the constant estimate by
\[
T_{\mathrm{zero}}(S_n,U)=0.
\]
Since the randomized experiments are probability laws,
\[
\sup_{Q\in\mathcal G(\sigma)}
\mathbb E_Q|T_{\mathrm{zero}}-\vartheta(Q)|
\le1.
\]
Thus the larger-class extended minimax risk is finite, and conversion to real values preserves the preceding inequality:
\[
R_\beta(n,\sigma)\le\mathcal R_{\mathcal G}(n,\sigma).
\]
% lean: CausalSmith.Stat.RdTruesideNoiseFrontier.benchmark_risk_le_gaussian

For interval length, define the constant interval procedure by
\[
I_{\mathrm{full}}(S_n,U)=[-1,1].
\]
The same target bounds imply, for every \(Q\in\mathcal G(\sigma)\),
\[
Q\{\vartheta(Q)\in I_{\mathrm{full}}\}=1,
\qquad
\mathbb E_Q\operatorname{length}(I_{\mathrm{full}})=2.
\]
Hence the larger-class honest interval class is nonempty and its extended minimax expected length is at most \(2\).

Now let \(I=[\ell,\upsilon]\) be any interval procedure honest on \(\mathcal G(\sigma)\). For every benchmark law \(P\), its embedded law \(Q\) satisfies
\[
\mathbb P_{P,n}\{\theta(P)\in I\}
=
Q\{\vartheta(Q)\in I\}
\ge0.9.
\]
Thus \(I\) belongs to the benchmark honest interval class of \cref{obj:def:honest-class}. Moreover,
\[
\mathcal L_\beta(n,\sigma)
\le
\sup_{P\in\mathcal M_\beta(\sigma)}
\mathbb E_{P,n}[\upsilon-\ell]
\le
\sup_{Q\in\mathcal G(\sigma)}
\mathbb E_Q[\upsilon-\ell].
\]
The second inequality follows from experiment preservation. Taking the infimum over intervals honest on \(\mathcal G(\sigma)\) proves the extended minimax comparison; the finite bound \(2\) permits its conversion to real values. Therefore
\[
\mathcal L_\beta(n,\sigma)\le\mathcal L_{\mathcal G}(n,\sigma).
\]
Combining these comparisons with the lower bounds already obtained yields
\[
c\,r_\beta(n,\sigma)
\le R_\beta(n,\sigma)
\le\mathcal R_{\mathcal G}(n,\sigma),
\qquad
c\,r_\beta(n,\sigma)
\le\mathcal L_\beta(n,\sigma)
\le\mathcal L_{\mathcal G}(n,\sigma).
\]
% lean: CausalSmith.Stat.RdTruesideNoiseFrontier.benchmark_length_le_gaussian

\item Finally, fix \(\sigma\in(0,1]\) and \(Q\in\mathcal G(\sigma)\). Its joint independence condition, followed by the measurable scaling map \(e\mapsto\sigma e\), gives
\[
E\perp(Z,B),
\qquad
\sigma E\perp(Z,B).
\]
These supply the standardized and scaled error independence in the stated Assumption 1Y condition. The local density condition supplies a radius \(r>0\) for which
\[
a(z)>0\qquad(z\in(-r,r)),
\]
which is Assumption 2C. Finally, scaling the standard Gaussian error gives
\[
N(0,1)\circ(e\mapsto\sigma e)^{-1}=N(0,\sigma^2),
\qquad
\sigma E\sim N(0,\sigma^2).
\]
This is Assumption 7, completing the three asserted conditions.
% lean: CausalSmith.Stat.RdTruesideNoiseFrontier.gaussian_published_hypotheses
\end{enumerate}
\end{proof}

\begin{proof}[Proof of \cref{obj:lem:classical-legendre-facts}]
Write the shifted Legendre polynomials in the normalization
\[
R_j(x):=\sum_{m=0}^{j}(-1)^m
  \binom{j}{m}\binom{j+m}{m}x^m,
\qquad
P_j^{\mathrm{Leg}}(t)=R_j\!\left(\frac{1-t}{2}\right),
\qquad j\in\mathbb N,\quad x,t\in\mathbb R.
\]

\begin{enumerate}
\item \emph{Orthogonality.}
The affine substitution \(t=1-2x\), including reversal of the integration endpoints, gives
\[
\int_0^1R_j(x)R_k(x)\,dx
=-\frac12\int_1^{-1}
  P_j^{\mathrm{Leg}}(t)P_k^{\mathrm{Leg}}(t)\,dt
=\frac12\int_{-1}^1
  P_j^{\mathrm{Leg}}(t)P_k^{\mathrm{Leg}}(t)\,dt.
\]
We therefore compute the inner product on the unit interval.

For \(k\in\mathbb N\), define
\[
\mathcal B_k(x):=x^k(1-x)^k,
\qquad x\in\mathbb R.
\]
Expanding \((1-x)^k\) and differentiating \(k\) times gives Rodrigues' identity:
\[
\mathcal B_k^{(k)}(x)
=\sum_{m=0}^k(-1)^m\binom{k}{m}
  \frac{(k+m)!}{m!}x^m
=k!R_k(x).
\]
For every \(0\le i<k\), the factors \(x^k\) and \((1-x)^k\), respectively, ensure that
\[
\mathcal B_k^{(i)}(0)=\mathcal B_k^{(i)}(1)=0.
\]
Indeed, in each term of the Leibniz expansion, fewer than \(k\) derivatives fall on the factor vanishing at the relevant endpoint. Consequently, repeated integration by parts gives, for every polynomial \(R_j\),
\[
\int_0^1R_j(x)\mathcal B_k^{(k)}(x)\,dx
=(-1)^k\int_0^1R_j^{(k)}(x)\mathcal B_k(x)\,dx.
\]
For \(k=0\), this identity involves zero integrations by parts and has no boundary conditions.
% lean: Causalean.Mathlib.Analysis.SpecialFunctions.Jacobi.integral_mul_iterated_deriv_eq

First consider \(j=k\). The coefficient of \(x^k\) in \(R_k\) gives
\[
R_k^{(k)}(x)=k!(-1)^k\binom{2k}{k}.
\]
Combining this identity with Rodrigues' formula and integration by parts yields
\[
\begin{aligned}
k!\int_0^1R_k(x)^2\,dx
&=\int_0^1R_k(x)\mathcal B_k^{(k)}(x)\,dx\\
&=(-1)^k\int_0^1R_k^{(k)}(x)\mathcal B_k(x)\,dx\\
&=k!\binom{2k}{k}\int_0^1x^k(1-x)^k\,dx.
\end{aligned}
\]
The beta integral evaluates to
\[
\int_0^1x^k(1-x)^k\,dx
=\frac{k!}{(k+1)\prod_{i=0}^{k-1}(k+2+i)}
=\frac{(k!)^2}{(2k+1)!},
\]
where the product is \(1\) when \(k=0\). Since \(k!>0\), cancellation and the factorial formula for the binomial coefficient give
\[
\int_0^1R_k(x)^2\,dx
=\binom{2k}{k}\frac{(k!)^2}{(2k+1)!}
=\frac{1}{2k+1}.
\]
% lean: Causalean.Mathlib.Analysis.SpecialFunctions.Jacobi.shiftedLegendre_sq_integral_unit

If \(j<k\), the degree of \(R_j\) is \(j\), so \(R_j^{(k)}=0\). Thus
\[
k!\int_0^1R_j(x)R_k(x)\,dx
=(-1)^k\int_0^1R_j^{(k)}(x)\mathcal B_k(x)\,dx
=0,
\]
and the inner product vanishes. If \(k<j\), interchange the two factors and apply the same calculation. Hence
\[
\int_0^1R_j(x)R_k(x)\,dx
=
\begin{cases}
\dfrac{1}{2j+1},&j=k,\\
0,&j\ne k.
\end{cases}
\]
Multiplying the affine-substitution identity by \(2\) proves the asserted inner product on \([-1,1]\).
% lean: CausalSmith.Stat.RdTruesideNoiseFrontier.legendreP_orthogonal_exact

\item \emph{Reflection symmetry.}
The shifted Legendre reflection identity is
\[
R_j(x)=(-1)^jR_j(1-x),
\qquad x\in\mathbb R.
\]
Applying it at \(x=(1+t)/2\) gives
\[
P_j^{\mathrm{Leg}}(-t)
=R_j\!\left(\frac{1+t}{2}\right)
=(-1)^jR_j\!\left(\frac{1-t}{2}\right)
=(-1)^jP_j^{\mathrm{Leg}}(t),
\qquad t\in\mathbb R.
\]
% lean: CausalSmith.Stat.RdTruesideNoiseFrontier.legendreP_reflection_exact

\item \emph{The bound on the interval.}
The finite sum defining \(R_j\) is also the normalized shifted Jacobi polynomial with both shape parameters equal to zero: its coefficient identity is
\[
(-1)^m\binom jm
\frac{(j+m)!/j!}{m!}
=(-1)^m\binom jm\binom{j+m}{m},
\qquad 0\le m\le j.
\]
For \(\xi\in\mathbb R\) and \(r\in\mathbb N\), define the rising and falling factorial products by
\[
(\xi)_{\uparrow r}:=\prod_{i=0}^{r-1}(\xi+i),
\qquad
(\xi)_{\downarrow r}:=\prod_{i=0}^{r-1}(\xi-i),
\]
with both empty products equal to \(1\). Their definitions give
\[
(-\xi)_{\downarrow r}=(-1)^r(\xi)_{\uparrow r},
\qquad
(\xi)_{\uparrow m}(\xi+m)_{\uparrow(j-m)}
=(\xi)_{\uparrow j},
\quad 0\le m\le j.
\]
The falling-factorial convolution is
\[
(\xi+\eta)_{\downarrow j}
=\sum_{m=0}^j\binom jm
  (\xi)_{\downarrow m}(\eta)_{\downarrow(j-m)},
\qquad \xi,\eta\in\mathbb R.
\]
For nonnegative integer arguments, comparing coefficients of degree \(j\) in
\[
(1+\zeta)^{\xi+\eta}
=(1+\zeta)^\xi(1+\zeta)^\eta,
\qquad \xi,\eta\in\mathbb N,\quad \zeta\in\mathbb R,
\]
and multiplying by \(j!\) proves the convolution. Both sides are polynomials in the two arguments, so the identity extends to all real arguments.

Specialize the convolution to \(\xi=-(j+1)\) and \(\eta=j\). The product identities yield, for \(0\le m\le j\),
\[
\begin{aligned}
(-1)_{\downarrow j}&=(-1)^j j!,\\
(-(j+1))_{\downarrow m}
&=(-1)^m(j+1)_{\uparrow m}
=(-1)^m\frac{(j+m)!}{j!},\\
(j)_{\downarrow(j-m)}
&=(m+1)_{\uparrow(j-m)}
=\frac{j!}{m!}.
\end{aligned}
\]
Consequently the convolution becomes
\[
(-1)^j j!
=\sum_{m=0}^j\binom jm(-1)^m
  \frac{(j+m)!}{j!}\frac{j!}{m!}
=\sum_{m=0}^j(-1)^m\binom jm\frac{(j+m)!}{m!}.
\]
Dividing by \(j!>0\) supplies the normalization:
\[
R_j(1)
=\sum_{m=0}^j(-1)^m\binom jm\binom{j+m}{m}
=\frac{1}{j!}\sum_{m=0}^j(-1)^m\binom jm\frac{(j+m)!}{m!}
=(-1)^j,
\qquad |R_j(1)|=1.
\]
These identities include \(j=0\), when the products are empty and \(0!=1\).
At the other endpoint, direct evaluation of the defining sum gives
\[
R_j(0)=1,
\qquad |R_j(0)|=1.
\]
% lean: Causalean.Mathlib.Analysis.SpecialFunctions.Jacobi.jacobiShifted_one

When \(j=0\), the defining sum gives \(R_0(x)=1\) for every real \(x\), proving the required bound. For \(j\ge1\), define
\[
\lambda_j:=j(j+1)>0,
\qquad
\mathcal E_j(x):=
R_j(x)^2+\frac{x(1-x)R_j'(x)^2}{\lambda_j},
\qquad x\in\mathbb R.
\]
We derive the differential identity governing this energy. Define the coefficients over their full range by
\[
c_{j,m}:=
\begin{cases}
(-1)^m\binom jm\binom{j+m}{m},&0\le m\le j,\\
0,&m>j,
\end{cases}
\qquad m\in\mathbb N.
\]
The binomial coefficient formulas imply
\[
(m+1)^2c_{j,m+1}
=-(j-m)(j+m+1)c_{j,m},
\qquad 0\le m<j.
\]
The coefficient of \(x^m\) in the polynomial
\[
x(1-x)R_j''(x)+(1-2x)R_j'(x)+\lambda_jR_j(x)
\]
is
\[
(m+1)^2c_{j,m+1}
+\bigl(\lambda_j-m(m+1)\bigr)c_{j,m}.
\]
It vanishes for \(m<j\) by the recurrence, for \(m=j\) because \(c_{j,j+1}=0\) and \(\lambda_j=j(j+1)\), and for \(m>j\) because both coefficients are zero. Therefore
\[
x(1-x)R_j''(x)+(1-2x)R_j'(x)+\lambda_jR_j(x)=0,
\qquad x\in\mathbb R.
\]
% lean: Causalean.Mathlib.Analysis.SpecialFunctions.Jacobi.jacobiShifted_ode

Differentiation of the energy, followed by substitution of this differential equation, yields
\[
\begin{aligned}
\mathcal E_j'(x)
&=2R_j(x)R_j'(x)
+\frac{(1-2x)R_j'(x)^2
       +2x(1-x)R_j'(x)R_j''(x)}{\lambda_j}\\
&=(2x-1)\frac{R_j'(x)^2}{\lambda_j}.
\end{aligned}
\]
Thus the energy is nonincreasing on \([0,1/2]\) and nondecreasing on \([1/2,1]\). Its endpoint values are
\[
\mathcal E_j(0)=R_j(0)^2=1,
\qquad
\mathcal E_j(1)=R_j(1)^2=1.
\]
For \(x\le1/2\), compare with the left endpoint; for \(x>1/2\), compare with the right endpoint. These comparisons give
\[
\mathcal E_j(x)\le
\begin{cases}
\mathcal E_j(0),&0\le x\le1/2,\\
\mathcal E_j(1),&1/2<x\le1,
\end{cases}
\qquad
\mathcal E_j(x)\le1.
\]
Since the additional term in the energy is nonnegative on the unit interval,
\[
R_j(x)^2
\le R_j(x)^2+\frac{x(1-x)R_j'(x)^2}{\lambda_j}
=\mathcal E_j(x)\le1,
\qquad 0\le x\le1.
\]
Hence \(|R_j(x)|\le1\) throughout that interval. Finally,
\[
-1\le t\le1
\quad\Longrightarrow\quad
0\le\frac{1-t}{2}\le1
\quad\Longrightarrow\quad
|P_j^{\mathrm{Leg}}(t)|
=\left|R_j\!\left(\frac{1-t}{2}\right)\right|\le1.
\]
% lean: CausalSmith.Stat.RdTruesideNoiseFrontier.legendreP_abs_le_one_exact

\item \emph{The endpoint value.}
Evaluation at the positive endpoint gives
\[
P_j^{\mathrm{Leg}}(1)=R_j(0)=1.
\]
Applying reflection symmetry at \(t=1\) therefore yields
\[
P_j^{\mathrm{Leg}}(-1)
=(-1)^jP_j^{\mathrm{Leg}}(1)
=(-1)^j.
\]
% lean: CausalSmith.Stat.RdTruesideNoiseFrontier.legendreP_neg_one_exact
\end{enumerate}
\end{proof}

\begin{proof}[Proof of \cref{obj:thm:direct-reduction}]
Fix \(0<\beta\le1\). For integers \(n\ge2\), define
\[
r_{\mathrm{dir},n}:=n^{-\beta/(2\beta+1)}.
\]
Expectations of nonnegative losses and their suprema below are taken in the extended nonnegative reals.

\begin{enumerate}
\item \textbf{The selected radius at zero noise.}
By \cref{obj:lem:positive-endpoint}, there is a constant \(K_{\mathrm{end}}>0\), depending only on \(\beta\), such that, for every integer \(L\ge1\),
\[
M_{L,\beta}\le K_{\mathrm{end}}L^{-2\beta},
\qquad
\int_0^1 K_L(x)^2\,dx\le K_{\mathrm{end}}L^2.
\]
Define
\[
K_{\mathrm{cert}}
:=48K_{\mathrm{end}}+28\sqrt{30K_{\mathrm{end}}}>0.
\]
For a fixed \(n\ge2\), put
\[
a_{\mathrm{deg}}:=n^{1/(4\beta+2)},
\qquad
L_{\mathrm{dir}}:=\lfloor a_{\mathrm{deg}}\rfloor,
\qquad
L_{\max}:=\lceil a_{\mathrm{deg}}\rceil.
\]
Since \(a_{\mathrm{deg}}\ge1\), the floor inequalities give
\[
1\le L_{\mathrm{dir}},
\qquad
a_{\mathrm{deg}}<L_{\mathrm{dir}}+1\le2L_{\mathrm{dir}},
\qquad
L_{\mathrm{dir}}\le a_{\mathrm{deg}},
\qquad
L_{\mathrm{dir}}\le L_{\max}.
\]
These inequalities apply throughout \(a_{\mathrm{deg}}\ge1\), including \(1\le a_{\mathrm{deg}}<2\). Consequently,
\[
\begin{aligned}
M_{L_{\mathrm{dir}},\beta}
&\le K_{\mathrm{end}}L_{\mathrm{dir}}^{-2\beta}\\
&\le K_{\mathrm{end}}(a_{\mathrm{deg}}/2)^{-2\beta}\\
&=K_{\mathrm{end}}a_{\mathrm{deg}}^{-2\beta}2^{2\beta}
\le4K_{\mathrm{end}}a_{\mathrm{deg}}^{-2\beta},
\end{aligned}
\]
where the last inequality uses \(\beta\le1\).

In the variance formula of \cref{obj:synth_6}, the coefficient of the zeroth derivative is one at \(\sigma=0\), and every positive-order coefficient vanishes. Thus
\[
V_L(0)=\frac34\int_0^1K_L(x)^2\,dx,
\]
and hence
\[
\frac{40V_{L_{\mathrm{dir}}}(0)}{n}
\le\frac{30K_{\mathrm{end}}L_{\mathrm{dir}}^2}{n}
\le\frac{30K_{\mathrm{end}}a_{\mathrm{deg}}^2}{n}.
\]
Taking square roots yields
\[
\sqrt{\frac{40V_{L_{\mathrm{dir}}}(0)}{n}}
\le\sqrt{30K_{\mathrm{end}}}\,
\frac{a_{\mathrm{deg}}}{\sqrt n}.
\]
The two balancing identities are
\[
a_{\mathrm{deg}}^{-2\beta}
=n^{-2\beta/(4\beta+2)}
=r_{\mathrm{dir},n},
\qquad
\frac{a_{\mathrm{deg}}}{\sqrt n}
=n^{1/(4\beta+2)-1/2}
=r_{\mathrm{dir},n}.
\]
The positive-index radius in \cref{obj:synth_6} therefore satisfies
\[
\rho_{L_{\mathrm{dir}}}
=12M_{L_{\mathrm{dir}},\beta}
+28\sqrt{\frac{40V_{L_{\mathrm{dir}}}(0)}{n}}
\le K_{\mathrm{cert}}r_{\mathrm{dir},n}.
\]
Since \(L_{\mathrm{dir}}\) belongs to the candidate set, the minimization in \cref{obj:def:public-selection} gives
\[
E_\beta(n,0)=\rho_{L_\star}
\le\rho_{L_{\mathrm{dir}}}
\le K_{\mathrm{cert}}r_{\mathrm{dir},n}.
\]
% lean: CausalSmith.Stat.RdTruesideNoiseFrontier.certificate_zero_rate_bound

\item \textbf{Direct witnesses and their resolution.}
Fix \(n\ge2\) and \(\sigma\in[0,1]\), and define
\[
h_0:=n^{-1/(2\beta+1)}.
\]
Because \(n\ge1\) and \(2\beta+1>0\),
\[
0<h_0\le1,
\qquad
nh_0^{2\beta+1}=1,
\qquad
h_0^\beta=r_{\mathrm{dir},n}.
\]

For \(0<h\le1\), use the laws \(P_h^{\mathrm{dir},\pm}\) of \cref{obj:def:direct-alternatives}. Define their mean perturbation by
\[
\delta_h(x):=\kappa h^\beta v_h^{\mathrm{dir}}(x),
\qquad x\in[-1,1],
\qquad
\kappa=\frac1{100}.
\]
The profile bounds in \cref{obj:synth_7} imply
\[
0\le\delta_h(x)\le\kappa,
\qquad
|\delta_h(x)-\delta_h(t)|
\le\kappa|x-t|^\beta
\quad(x,t\in[-1,1]).
\]
Together with the conditional means in \cref{obj:def:direct-alternatives}, these give
\[
\frac14
\le\mu_d(P_h^{\mathrm{dir},\pm},x)
\le\frac34,
\qquad
[\mu_0(P_h^{\mathrm{dir},\pm},\cdot)]_\beta=0,
\qquad
[\mu_1(P_h^{\mathrm{dir},\pm},\cdot)]_\beta\le\kappa\le2.
\]
Their score density is continuous and satisfies
\[
f(P_h^{\mathrm{dir},\pm},x)=\frac12
\quad(x\in[-1,1]),
\qquad
\int_{-1}^1f(P_h^{\mathrm{dir},\pm},x)\,dx=1.
\]
Their Gaussian error and independence properties are part of their construction. Thus all the conditions of \cref{obj:def:model} hold:
\[
P_h^{\mathrm{dir},\pm}\in\mathcal M_\beta(\sigma)
\quad(0<h\le1,\ 0\le\sigma\le1).
\]
Since \(v_h^{\mathrm{dir}}(0)=1\), their targets satisfy
\[
\theta(P_h^{\mathrm{dir},+})=\kappa h^\beta,
\qquad
\theta(P_h^{\mathrm{dir},-})=-\kappa h^\beta,
\qquad
\left|\theta(P_h^{\mathrm{dir},+})
-\theta(P_h^{\mathrm{dir},-})\right|
=2\kappa h^\beta.
\]
% lean: CausalSmith.Stat.RdTruesideNoiseFrontier.directAltLaw_model
% lean: CausalSmith.Stat.RdTruesideNoiseFrontier.directResolution_balance
% lean: CausalSmith.Stat.RdTruesideNoiseFrontier.directAltLaw_separation

\item \textbf{Noiseless information bound.}
For \(0<h\le1\), define the one-observation laws
\[
\mathsf P_{\sigma,h}^{\pm}
:=\operatorname{Law}_{P_h^{\mathrm{dir},\pm}}
(X+\sigma\varepsilon,D,Y).
\]
For probability measures \(\mathsf Q_1,\mathsf Q_2\) on a common measurable space, use
\[
\operatorname{TV}(\mathsf Q_1,\mathsf Q_2)
:=\sup_{\mathsf F\ \mathrm{measurable}}
|\mathsf Q_1(\mathsf F)-\mathsf Q_2(\mathsf F)|.
\]
In the absolutely continuous setting used here, the chi-squared divergence is
\[
\chi^2(\mathsf Q_1,\mathsf Q_2)
:=\int
\left(\frac{d\mathsf Q_1}{d\mathsf Q_2}-1\right)^2
\,d\mathsf Q_2.
\]

Let the reference measure on \(\mathbb R\times\{0,1\}^2\) be defined, for every measurable \(\mathsf F\), by
\[
\eta_{\mathrm{obs}}(\mathsf F)
:=\sum_{d=0}^1\sum_{y=0}^1
\int_{\mathbb R}\mathbf1_{\mathsf F}(x,d,y)\,dx.
\]
The densities
\[
p_h^\pm:=\frac{d\mathsf P_{0,h}^\pm}{d\eta_{\mathrm{obs}}}
\]
are given by
\[
\begin{aligned}
p_h^\pm(x,1,1)&=\frac14\pm\frac{\delta_h(x)}2,
&&0\le x\le1,\\
p_h^\pm(x,1,0)&=\frac14\mp\frac{\delta_h(x)}2,
&&0\le x\le1,\\
p_h^\pm(x,0,y)&=\frac14,
&&-1\le x<0,\quad y\in\{0,1\},
\end{aligned}
\]
and are zero elsewhere. Define a likelihood ratio on the entire observation space by
\[
\Lambda_h(x,d,y):=
\begin{cases}
p_h^+(x,1,y)/p_h^-(x,1,y),
&d=1,\ 0\le x\le1,\\
1,&\text{otherwise}.
\end{cases}
\]
The treated-cell densities lie between \(1/8\) and \(3/8\), so
\[
d\mathsf P_{0,h}^+
=\Lambda_h\,d\mathsf P_{0,h}^-,
\qquad
0\le\Lambda_h\le3,
\qquad
(\Lambda_h-1)^2\le4.
\]
In particular,
\[
\mathsf P_{0,h}^+\ll\mathsf P_{0,h}^-,
\qquad
\int(\Lambda_h-1)^2\,d\mathsf P_{0,h}^-<\infty.
\]

For each treated cell,
\[
p_h^+(x,1,y)-p_h^-(x,1,y)
=(2y-1)\delta_h(x),
\qquad
p_h^-(x,1,y)\ge\frac18.
\]
Consequently,
\[
0\le
p_h^-(x,1,y)(\Lambda_h(x,1,y)-1)^2
=\frac{\delta_h(x)^2}{p_h^-(x,1,y)}
\le8\delta_h(x)^2.
\]
Outside the treated cells, the corresponding integrand is zero. Summing over the two outcome cells gives
\[
\chi^2(\mathsf P_{0,h}^+,\mathsf P_{0,h}^-)
\le16\kappa^2h^{2\beta}
\int_0^1[v_h^{\mathrm{dir}}(x)]^2\,dx.
\]
The profile is \(1-x/h\) on \([0,h]\) and zero on \([h,1]\). With the substitution \(t=1-x/h\),
\[
\int_0^1[v_h^{\mathrm{dir}}(x)]^2\,dx
=\int_0^h(1-x/h)^2\,dx
=h\int_0^1t^2\,dt
=\frac h3.
\]
It follows that
\[
\chi^2(\mathsf P_{0,h}^+,\mathsf P_{0,h}^-)
\le\frac{16}{3}\kappa^2h^{2\beta+1}.
\]
At \(h=h_0\), the balance identity therefore yields
\[
n\chi^2(\mathsf P_{0,h_0}^+,\mathsf P_{0,h_0}^-)
\le\frac{16}{3}\kappa^2
=\frac{16}{30000}
\le\frac1{100}.
\]
% lean: CausalSmith.Stat.RdTruesideNoiseFrontier.directAltLaw_observed_ac_zero
% lean: CausalSmith.Stat.RdTruesideNoiseFrontier.directAltLaw_observed_likelihood_square_integrable_zero
% lean: CausalSmith.Stat.RdTruesideNoiseFrontier.directAltLaw_observed_chiSqDiv_zero_le

\item \textbf{Product testing bound and Gaussian contraction.}
Define the sample laws by
\[
\mathsf S_{\sigma,h}^{\pm}
:=(\mathsf P_{\sigma,h}^{\pm})^{\otimes n},
\]
and abbreviate the one-observation divergence at the selected resolution as
\[
\chi_{\mathrm{dir}}
:=\chi^2(\mathsf P_{0,h_0}^+,\mathsf P_{0,h_0}^-).
\]
The product likelihood ratio is
\[
\frac{d\mathsf S_{0,h_0}^+}{d\mathsf S_{0,h_0}^-}
\bigl((o_i)_{i=1}^n\bigr)
=\prod_{i=1}^n\Lambda_{h_0}(o_i),
\qquad
(o_i)_{i=1}^n\in(\mathbb R\times\{0,1\}^2)^n,
\]
almost everywhere under \(\mathsf S_{0,h_0}^-\). Since the one-observation likelihood ratio integrates to one,
\[
\int\Lambda_{h_0}\,d\mathsf P_{0,h_0}^-=1,
\qquad
\chi_{\mathrm{dir}}
=\int(\Lambda_{h_0}-1)^2\,d\mathsf P_{0,h_0}^-
=\int\Lambda_{h_0}^2\,d\mathsf P_{0,h_0}^- -1.
\]
Product integration likewise gives
\[
\int\prod_{i=1}^n\Lambda_{h_0}(o_i)
\,d\mathsf S_{0,h_0}^-\bigl((o_i)_{i=1}^n\bigr)
=\left(\int\Lambda_{h_0}\,d\mathsf P_{0,h_0}^-\right)^n
=1.
\]
Expanding the squared likelihood deviation and factoring its second moment over the product measure therefore yields
\[
\begin{aligned}
1+\chi^2(\mathsf S_{0,h_0}^+,\mathsf S_{0,h_0}^-)
&=\int\left(\prod_{i=1}^n\Lambda_{h_0}(o_i)\right)^2
\,d\mathsf S_{0,h_0}^-\bigl((o_i)_{i=1}^n\bigr)\\
&=\left(\int\Lambda_{h_0}^2\,d\mathsf P_{0,h_0}^-\right)^n\\
&=(1+\chi_{\mathrm{dir}})^n.
\end{aligned}
\]
Since \(\chi_{\mathrm{dir}}\ge0\) and \(n\chi_{\mathrm{dir}}\le1/100\),
\[
(1+\chi_{\mathrm{dir}})^n
\le e^{n\chi_{\mathrm{dir}}}
\le e^{1/100}
\le\sum_{j=0}^\infty(1/100)^j
=\frac{100}{99}.
\]
Here the penultimate inequality follows from the exponential series and \(j!\ge1\). Thus
\[
\chi^2(\mathsf S_{0,h_0}^+,\mathsf S_{0,h_0}^-)
\le\frac1{99}\le\frac4{25}.
\]
For absolutely continuous probability measures, the signed density difference has integral zero, and Cauchy--Schwarz gives
\[
\operatorname{TV}(\mathsf Q_1,\mathsf Q_2)
=\frac12\int
\left|\frac{d\mathsf Q_1}{d\mathsf Q_2}-1\right|
\,d\mathsf Q_2
\le\frac12\sqrt{\chi^2(\mathsf Q_1,\mathsf Q_2)}.
\]
Applying this to the product laws gives
\[
\operatorname{TV}(\mathsf S_{0,h_0}^+,\mathsf S_{0,h_0}^-)
\le\frac12\sqrt{\frac4{25}}=\frac15.
\]

For a sample
\[
\mathbf s=((x_i,d_i,y_i))_{i=1}^n
\in(\mathbb R\times\{0,1\}^2)^n,
\]
define the common Gaussian observation kernel by
\[
\mathsf K_{\sigma,n}(\mathbf s,\mathsf F)
:=
\Pr\!\left\{
((x_i+\sigma\varepsilon_i,d_i,y_i))_{i=1}^n
\in\mathsf F
\right\},
\qquad
(\varepsilon_i)_{i=1}^n\sim N(0,1)^{\otimes n}.
\]
The Gaussian variables in this definition are independent of the input sample. Gaussian error independence under the direct laws implies
\[
\mathsf S_{\sigma,h_0}^{\pm}
=\mathsf K_{\sigma,n}\mathsf S_{0,h_0}^{\pm}.
\]
For every measurable output event \(\mathsf F\), the function
\(\mathsf K_{\sigma,n}(\cdot,\mathsf F)\) takes values in \([0,1]\). The layer-cake formula therefore bounds the difference of its integrals by
\[
\begin{aligned}
&\left|
\int\mathsf K_{\sigma,n}(\mathbf s,\mathsf F)
\,d\mathsf S_{0,h_0}^+(\mathbf s)
-
\int\mathsf K_{\sigma,n}(\mathbf s,\mathsf F)
\,d\mathsf S_{0,h_0}^-(\mathbf s)
\right|\\
&\qquad\le
\int_0^1
\left|
\mathsf S_{0,h_0}^+
\{\mathsf K_{\sigma,n}(\cdot,\mathsf F)>t\}
-
\mathsf S_{0,h_0}^-
\{\mathsf K_{\sigma,n}(\cdot,\mathsf F)>t\}
\right|\,dt\\
&\qquad\le
\operatorname{TV}(\mathsf S_{0,h_0}^+,\mathsf S_{0,h_0}^-).
\end{aligned}
\]
Taking the supremum over output events proves
\[
\operatorname{TV}(\mathsf S_{\sigma,h_0}^+,\mathsf S_{\sigma,h_0}^-)
\le
\operatorname{TV}(\mathsf S_{0,h_0}^+,\mathsf S_{0,h_0}^-)
\le\frac15.
\]
The construction and this inequality also apply at \(\sigma=0\).
% lean: Causalean.Stat.iid_tv_le_one_fifth_of_chiSqDiv
% lean: CausalSmith.Stat.RdTruesideNoiseFrontier.sampleLaw_tv_le_zero_of_model

\item \textbf{Transfer to estimation risk and honest length.}
Let
\[
\mathsf U_{\mathrm{seed}}:=\operatorname{Uniform}[0,1],
\qquad
\mathsf E_{\sigma,h_0}^{\pm}
:=\mathsf S_{\sigma,h_0}^{\pm}\otimes\mathsf U_{\mathrm{seed}}.
\]
The product total-variation inequality gives
\[
\begin{aligned}
\operatorname{TV}(\mathsf E_{\sigma,h_0}^+,\mathsf E_{\sigma,h_0}^-)
&\le
\operatorname{TV}(\mathsf S_{\sigma,h_0}^+,\mathsf S_{\sigma,h_0}^-)
+\operatorname{TV}(\mathsf U_{\mathrm{seed}},\mathsf U_{\mathrm{seed}})\\
&\le\frac15.
\end{aligned}
\]
Define the common submeasure by
\[
d\mathsf H_{\sigma,h_0}
:=
\min\!\left\{
\frac{d\mathsf E_{\sigma,h_0}^+}
{d(\mathsf E_{\sigma,h_0}^++\mathsf E_{\sigma,h_0}^-)},
\frac{d\mathsf E_{\sigma,h_0}^-}
{d(\mathsf E_{\sigma,h_0}^++\mathsf E_{\sigma,h_0}^-)}
\right\}
d(\mathsf E_{\sigma,h_0}^++\mathsf E_{\sigma,h_0}^-).
\]
It satisfies
\[
\mathsf H_{\sigma,h_0}\le\mathsf E_{\sigma,h_0}^{\pm},
\qquad
\int d\mathsf H_{\sigma,h_0}
=
1-\operatorname{TV}(\mathsf E_{\sigma,h_0}^+,\mathsf E_{\sigma,h_0}^-)
\ge\frac45.
\]
The mass identity follows from
\(\min\{u,v\}=(u+v-|u-v|)/2\).

For any admissible estimator \(\widehat\theta\), domination by the common submeasure and the triangle inequality imply
\[
\begin{aligned}
&\mathbb E_{P_{h_0}^{\mathrm{dir},+},n}
|\widehat\theta-\theta(P_{h_0}^{\mathrm{dir},+})|
+
\mathbb E_{P_{h_0}^{\mathrm{dir},-},n}
|\widehat\theta-\theta(P_{h_0}^{\mathrm{dir},-})|\\
&\quad\ge
\int\left(
|\widehat\theta-\theta(P_{h_0}^{\mathrm{dir},+})|
+
|\widehat\theta-\theta(P_{h_0}^{\mathrm{dir},-})|
\right)\,d\mathsf H_{\sigma,h_0}\\
&\quad\ge
2\kappa h_0^\beta\int d\mathsf H_{\sigma,h_0}
\ge\frac85\kappa h_0^\beta.
\end{aligned}
\]
At least one of the two extended risks is therefore at least
\(\frac45\kappa h_0^\beta\), and both laws belong to the model class. Hence
\[
\inf_{\widehat\theta}
\sup_{P\in\mathcal M_\beta(\sigma)}
\mathbb E_{P,n}|\widehat\theta-\theta(P)|
\ge\frac45\kappa h_0^\beta.
\]

For any \(I\in\mathcal I_{\beta,n,\sigma}\), define its two coverage events by
\[
\mathsf F_{h_0,I}^{\pm}
:=
\{(S_n,U):\theta(P_{h_0}^{\mathrm{dir},\pm})\in I(S_n,U)\}.
\]
Honesty as defined in \cref{obj:def:honest-class} gives
\[
\mathsf E_{\sigma,h_0}^{+}(\mathsf F_{h_0,I}^{+})\ge\frac9{10},
\qquad
\mathsf E_{\sigma,h_0}^{-}(\mathsf F_{h_0,I}^{-})\ge\frac9{10}.
\]
Total variation transfers the first probability to the minus law:
\[
\mathsf E_{\sigma,h_0}^{-}(\mathsf F_{h_0,I}^{+})
\ge\frac9{10}-\frac15=\frac7{10}.
\]
Using the union bound under that law,
\[
\mathsf E_{\sigma,h_0}^{-}
(\mathsf F_{h_0,I}^{+}\cap\mathsf F_{h_0,I}^{-})
\ge\frac7{10}+\frac9{10}-1=\frac35.
\]
On the intersection, the connected interval contains both targets and has length at least their separation. Thus
\[
\mathbb E_{P_{h_0}^{\mathrm{dir},-},n}\operatorname{length}(I)
\ge2\kappa h_0^\beta
\mathsf E_{\sigma,h_0}^{-}
(\mathsf F_{h_0,I}^{+}\cap\mathsf F_{h_0,I}^{-})
\ge\frac65\kappa h_0^\beta.
\]
Taking the infimum over honest intervals yields
\[
\inf_{I\in\mathcal I_{\beta,n,\sigma}}
\sup_{P\in\mathcal M_\beta(\sigma)}
\mathbb E_{P,n}\operatorname{length}(I)
\ge\frac65\kappa h_0^\beta.
\]
% lean: CausalSmith.Stat.RdTruesideNoiseFrontier.experiment_tv_le
% lean: CausalSmith.Stat.RdTruesideNoiseFrontier.twoPoint_absRisk_sum
% lean: CausalSmith.Stat.RdTruesideNoiseFrontier.twoPoint_interval_lower
% lean: CausalSmith.Stat.RdTruesideNoiseFrontier.direct_two_point_transfer

\item \textbf{Real lower bounds.}
For every \(P\in\mathcal M_\beta(\sigma)\), the mean bounds give
\[
|\theta(P)|
=|\mu_1(P,0)-\mu_0(P,0)|
\le\frac12.
\]
The constant-zero estimator consequently satisfies
\[
\sup_{P\in\mathcal M_\beta(\sigma)}
\mathbb E_{P,n}|0-\theta(P)|\le\frac12.
\]
The constant interval \([-1,1]\) satisfies
\[
\mathbb P_{P,n}\{\theta(P)\in[-1,1]\}=1,
\qquad
\mathbb E_{P,n}\operatorname{length}([-1,1])=2
\quad(P\in\mathcal M_\beta(\sigma)).
\]
It is therefore an admissible honest interval. These two procedures bound the extended minimax values by finite numbers, so conversion to real numbers preserves the preceding lower inequalities. Since \(h_0^\beta=r_{\mathrm{dir},n}\),
\[
R_\beta(n,\sigma)\ge\frac45\kappa r_{\mathrm{dir},n},
\qquad
\mathcal L_\beta(n,\sigma)\ge\frac65\kappa r_{\mathrm{dir},n}.
\]
Define
\[
c:=\frac45\kappa>0.
\]
As \(\kappa r_{\mathrm{dir},n}\ge0\), both criteria are at least
\(c\,r_{\mathrm{dir},n}\), uniformly over \(n\ge2\) and \(\sigma\in[0,1]\).
% lean: CausalSmith.Stat.RdTruesideNoiseFrontier.risk_value_ne_top
% lean: CausalSmith.Stat.RdTruesideNoiseFrontier.length_value_ne_top
% lean: CausalSmith.Stat.RdTruesideNoiseFrontier.direct_uniform_lower

\item \textbf{Attainment and minimax upper bounds.}
Set
\[
C:=2K_{\mathrm{cert}}>0,
\qquad
P_{\mathrm{unit}}:=P_1^{\mathrm{dir},+}.
\]
The witness verification above, at \(h=1\), gives
\[
P_{\mathrm{unit}}\in\mathcal M_\beta(0).
\]
For every \(n\ge2\), apply \cref{obj:thm:finite-certificate} with
\(\sigma=0\), \(L=1\), and \(P=P_{\mathrm{unit}}\). Its selected-procedure conclusions give
\[
I_{L_\star}\in\mathcal I_{\beta,n,0},
\]
\[
\sup_{P\in\mathcal M_\beta(0)}
\mathbb E_{P,n}
|\widehat\theta_{L_\star}-\theta(P)|
\le E_\beta(n,0),
\qquad
\sup_{P\in\mathcal M_\beta(0)}
\mathbb E_{P,n}\operatorname{length}(I_{L_\star})
\le2E_\beta(n,0).
\]
Since \(K_{\mathrm{cert}}r_{\mathrm{dir},n}\ge0\), the radius bound from the first step implies
\[
\sup_{P\in\mathcal M_\beta(0)}
\mathbb E_{P,n}
|\widehat\theta_{L_\star}-\theta(P)|
\le K_{\mathrm{cert}}r_{\mathrm{dir},n}
\le C r_{\mathrm{dir},n},
\]
\[
\sup_{P\in\mathcal M_\beta(0)}
\mathbb E_{P,n}\operatorname{length}(I_{L_\star})
\le2K_{\mathrm{cert}}r_{\mathrm{dir},n}
=C r_{\mathrm{dir},n}.
\]
Each extended minimax value is bounded above by the worst-case loss of its displayed admissible procedure. The selected interval's honesty supplies its admissibility for the length minimization. The finite upper bounds permit order-preserving conversion to real numbers, giving
\[
R_\beta(n,0)\le C r_{\mathrm{dir},n},
\qquad
\mathcal L_\beta(n,0)\le C r_{\mathrm{dir},n}.
\]
The constants \(c\) and \(C\) depend only on \(\beta\), and the displayed inequalities and honesty statement hold for every integer \(n\ge2\), as required.
% lean: CausalSmith.Stat.RdTruesideNoiseFrontier.risk_le_of_attainerRisk_le
% lean: CausalSmith.Stat.RdTruesideNoiseFrontier.lengthRisk_le_of_attainerLength_le
% lean: CausalSmith.Stat.RdTruesideNoiseFrontier.direct_reduction
\end{enumerate}
\end{proof}

\begin{proof}[Proof of \cref{obj:thm:uniform-frontier-resolution}]
\begin{enumerate}
\item Apply \cref{obj:thm:uniform-frontier} with the given \(\beta\). It supplies positive constants \(c,C\), depending only on \(\beta\), and a unique resolution in \([h_0,1]\). Its defining equation is equivalent to the equation in the present statement:
\[
\log n+\nu\log h_\star=\Psi(h_\star,\sigma)
\quad\Longleftrightarrow\quad
\nu\log(1/h_\star)+\Psi(h_\star,\sigma)=\log n.
\]
The rate and selected procedures are precisely those of
\cref{obj:def:uniform-rate,obj:def:public-selection}. Consequently, the two estimation and expected-length comparison chains, together with interval honesty, follow directly from \cref{obj:thm:uniform-frontier}. Its decision-class conclusions include procedures with an independent uniform seed. The same constants apply simultaneously for all \(n\ge2\) and \(\sigma\in[0,1]\), and hence along every shrinking-noise sequence.
% lean: CausalSmith.Stat.RdTruesideNoiseFrontier.uniform_frontier

\item The direct and intermediate formulas and the strict compact formula are supplied by \cref{obj:thm:uniform-frontier}. At the intermediate--compact equality, that result gives
\[
h_\star=\sigma^4,\qquad z=\sigma^{-2},\qquad \tau=\sigma^{-2},
\]
with both scalar equations satisfied. Thus the compact representation also includes the equality case. Its scalar equation uniquely determines \(\tau\) on the stated domain: for positive \(\sigma\), its left side has derivative
\[
\frac{2\nu}{\tau}+2+\log(\sigma^2\tau)>0
\qquad(\tau\ge\sigma^{-2}).
\]
The direct interface is likewise supplied by the same result:
\[
h_\star=\sigma=h_0.
\]
Finally, its specialization to \(\beta=1\) and \(\sigma=n^{-1/6}\) gives the intermediate branch for every \(n\ge2\) and the asserted comparison
\[
r_1(n,n^{-1/6})
\asymp
\frac{n^{-1/6}}{[1+\tfrac12\log n]^{3/2}}.
\]
These conclusions use the observed side label \(D=\mathbf1\{X\ge0\}\).
% lean: CausalSmith.Stat.RdTruesideNoiseFrontier.uniform_frontier_resolution
\end{enumerate}
\end{proof}

\bibliographystyle{plainnat}

\bibliography{references}

\end{document}
